\part{Chromatic Homotopy Theory}

\chapter{The Lazard Ring and Quillen's Theorem}

\section{Formal group laws and the Lazard ring}

	\begin{defn}\label{Formal_group_law}
		For a commutative ring $R$, we define a formal group law over $R$ to be a formal power series $f\in R[[x,y]]$, satisfying the following identities. 
		\begin{enumerate}
			\item (Unit) $f(x,0)=f(0,x)=x$. 
			\item (Commutativity) $f(x,y)=f(y,x)$. 
			\item (Associativity) $f(x,f(y,z))=f(f(x,y),z)$. 
		\end{enumerate}
		We denote $\FGL(R)$ by the set of all formal group laws over $R$. For a morphism of rings $\varphi:R\to R'$ and a formal group law $f$ over $R$, we denote $\varphi_*f$ as the formal group law over $R'$ givne by the map $\varphi_*:R[[x,y]]\to R'[[x,y]]$. As a consequence, $\FGL$ is a functor from $\Ring$ to $\Set$. 
	\end{defn}

	\begin{prop}
		The functor $\FGL$ is corepresentable. That is, there is a ring $R$ and a formal group law $f_{univ}$ over $R$ such that for any ring $S$ and a formal group law $g$ over $S$, there is a unique map $\varphi:R\to S$ such that $\varphi_*f_{univ}=g$. 
	\end{prop}
	\begin{proof}
		Notice that for $f\in R[[x,y]]=\sum a_{i,j}x^iy^j$ to be a formal group law, the three conditions translates to be $a_{i,0}=a_{0,j}=0$ for $i,j\ge1$, $a_{i,j}=a_{j,i}$, and the associativity to a complicated relation we won't write down here. Therefore, we may simply let $R=\mathbb{Z}[c_{i,j}]/I$, where $I$ is the ideal generated by the above relations. 
	\end{proof}
	\begin{defn}\label{Lazard_ring}
		The universal ring is denoted by $L$, the lazard ring, and the formal group law $f_{univ}$. 
	\end{defn}
	
	Our goal of this section is to characterize the structure of the ring $L$. In fact, we have the following theorem. 
	\begin{thm}\label{Lazard_thm}
		$L\simeq\mathbb{Z}[t_1,t_2,\dots]$, where $t_i$ has degree $2i$. 
	\end{thm}

	The first thing is to explain where the grading comes from. 
	\begin{remark}
		We let the degree of $c_{i,j}$ to be $2(i+j-1)$. Therefore, if we let $x$ and $y$ have degree $-2$, then $\sum c_{i,j}x^iy^j$ has degree $-2$. This grading comes from topology, where $x$ and $y$ are often the first Chern classes, and $f(x,y)$ is often the first Chern class of the tensor product bundle. \\
		Actually, we may consider the multiplicative group scheme $\mathbb{G}_m=\Spec R[t,t^{-1}]$. Notice that we have an action of $R^\times$ on $\FGL(R)$, given by $f(x,y)\mapsto\lambda^{-1}f(\lambda x,\lambda y)$. Later, for a graded ring $R$, we will denote the set of formal group laws on $R$ respecting the gradings by $\FGL^{\gr}(R)$. Clearly, this is corepresented by the set of graded maps $\Hom^{\gr}(L,R)$. 
	\end{remark}
	\begin{const}
		We first construct a map $\varphi:L\to\mathbb{Z}[b_1,b_2,\dots]$. We consider the formal power series $g(x)=x+b_1x^2+b_2x^3+\cdots$. Then we have a formal group law over $\mathbb{Z}[b_1,b_2,\dots]$ given by $f(x,y)=g(g^{-1}(x)+g^{-1}(y))$. Notice if we let $x,y$ have degree -2, then we result in an element of degree 2. We define $\varphi$ as this classifying map, which is a map of graded rings. 
	\end{const}

	The classifying map $\varphi$ turns out to have a very interesting property: it is an isomorphism with rational coefficients, but not with $\mathbb{F}_p$ coefficients. 

	\begin{lemma}\label{Lazard_thm_lem}
		Let $I$ be the ideal in $L$ generated by elements of positive degree, and $J$ the ideal in $\mathbb{Z}[b_1,b_2,\dots]$ generated by elements of positive degree. Notice $J/J^2$ is a free abelian group with one generator (the image of $b_i$) of degree $2i$. Also, notice $I$ inherits a natural grading from $L$. Then, we have $\varphi_{2n}:(I/I^2)_{2n}\to(J/J^2)_{2n}\simeq\mathbb{Z}$ injective, with image $p\mathbb{Z}$ if $n$ is of the form $p^r-1$ and $\mathbb{Z}$ otherwise. 
	\end{lemma}
	We postpone the proof of this lemma and first use it to prove \cref{Lazard_thm}. 

	\begin{proof}
		We now prove \cref{Lazard_thm}. By \cref{Lazard_thm_lem}, the abelian group $(I/I^2)_{2n}$ is isomorphic to $\mathbb{Z}$. Therefore we may choose $t_i\in L_{2n}$ lifting the generator of $(I/I^2)_{2n}$. Therefore we obtain a map of graded rings $\theta:\mathbb{Z}[t_1,\dots]\to L$. \\
		We prove $\theta$ is surjective. We prove by induction that $\theta$ induces a surjection in degree $2n$. Now, the induction hypothesis shows that the image of $\theta$ contains $(I^2)_{2n}$. Also, the image of $\theta$ contains a generator for $(I/I^2)_{2n}$, hence it contains $I_{2n}=L_{2n}$. \\
		We finally prove $\theta$ is injective. In fact, it suffices to prove the composite $\psi=\varphi\circ\theta:\mathbb{Z}[t_1,\dots]\to L\to\mathbb{Z}[b_1,\dots]$ is a rational isomorphism. In fact, since $\psi_{\mathbb{Q}}$ is a map of graded rings, and for each $n$ we have $\mathbb{Q}[t_1,\dots]_n$ and $\mathbb{Q}[b_1,\dots]_n$ a finite $\mathbb{Q}$-vector space of the same finitely many dimensions, it suffices to prove $\psi_{\mathbb{Q}}$ is surjective. This is done by letting $I'$ be the ideal of $\mathbb{Q}[t_1,\dots]$ generated by $t_i$. Again by \cref{Lazard_thm_lem}, we have a surjection $I'/(I')^2\to(J/J^2)\otimes\mathbb{Q}$. By repeating the proof of the surjectivity of $\theta$, we have $\psi_{\mathbb{Q}}$ is surjective. 
	\end{proof}

	Now, to prove \cref{Lazard_thm_lem}, we first need to identify the abelian group $(I/I^2)_{2n}$. We will determine this group by proving that for any Abelian group $M$, we have $\Hom((I/I^2)_{2n},M)\simeq M$. We denote the functor corepresented by $(I/I^2)_{2n}$ by $F(M)$ for a fixed $n$. Obviously, $F$ is related to the functor $\FGL$. 
	\begin{const}
		In this section, for any Abelian group $M$, we let $\mathbb{Z}\oplus M$ denote the graded ring with the $0$-th grade $\mathbb{Z}$ and the $2n$-th grade $M$, with multiplication law given by $(a,m)(a',m')=(aa',am'+a'm)$. Therefore, we have $\FGL^{\gr}(\mathbb{Z}\oplus M)\simeq\Hom^{\gr}(L,\mathbb{Z}\oplus M)\simeq F(M)$. In other words, $F(M)$ can be identified with the following formal group laws
		\[ f(x,y)=x+y+\sum_{i+j=n+1}m_{i,j}x^iy^j \]
		where $m_{i,j}\in M$. For $f(x,y)$ of this form to be a formal group law, the three requirements translates into the following. 
		\begin{enumerate}
			\item $m_{n+1,0}=m_{0,n+1}=0$. 
			\item $m_{i,j}=m_{j,i}$. 
			\item \[ \binom{i+j}{j}m_{i+j,k}=\binom{j+k}{j}m_{i,j+k} \]
		\end{enumerate}
	\end{const}

	Therefore, we have now translated the problem into solving the following equations in an Abelian group $M$. 

	\begin{const}
		Recall that we have a map $(I/I^2)_{2n}\to(J/J^2)_{2n}\simeq\mathbb{Z}$. We consider the map 
		\[ \lambda:M\simeq\Hom((J/J^2)_{2n},M)\to\Hom((I/I^2)_{2n},M)\simeq F(M) \]
		Here, notice that $\Hom((J/J^2)_{2n},M)$ is equal to $\Hom^{\gr}(\mathbb{Z}[b_1,\dots],\mathbb{Z}\oplus M)$, and $m$ corresponds to the map $\psi_m:\mathbb{Z}[b_1,\dots]\to\mathbb{Z}\oplus M$ sending $b_n$ to $m$ and other $b_i$ to $0$. Therefore the change of variables corresponding to $\psi_m$ is $g(x)=x+mx^{n+1}$. Since $m^2=0$ in $\mathbb{Z}\oplus M$, the inverse of $g(x)$ is obviously $g^{-1}(x)=x-mx^{n+1}$. Thus $g$ defines the formal group law 
		\[ f(x,y)=g(g^{-1}(x)+g^{-1}(y))=g(x-mx^{n+1}+y-my^{n+1})=x+y+m((x+y)^{n+1}-x^{n+1}-y^{n+1}) \]
		Therefore, $\lambda$ takes $m$ to the sequence $m_{i,j}$, where $m_{i,j}=\binom{n+1}{i}m$ for $i,j\neq0$. \\
		However, there can be more solutions. For example, if $\binom{n+1}{i}$ has a common divisor $d$, then $m_{i,j}/d$ is also a valid solution. 
	\end{const}
	We then try to prove \cref{Lazard_thm_lem}. Before this, we need some results in elementary number theory. 

	\begin{lemma}
		Suppose $a$ and $b$ are integers with base $p$ expansions $a=\sum a_ip^i$, $b=\sum b_ip^i$. Then we have
		\[ \binom{a}{b}=\prod\binom{a_i}{b_i} \]
	\end{lemma}
	\begin{proof}
		This is elementary. 
	\end{proof}
	\begin{coro}
		$\binom{i+j}{i}$ is not divisible by $p$ iff each digit of the $p$-base expansion of $i+j$ is as large as $i$. In other words, iff $i+j$ is calculated without carrying. 
	\end{coro}
	\begin{coro}
		Let $d$ be the greatest common divisor of $\binom{n+1}{i}\neq 1$. Then $d=p$ if $n+1=p^r$ and $1$ otherwise. 
	\end{coro}

	\begin{prop}
		Let $\lambda':M\to F(M)$ be the map carrying $m$ to the sequence $m_{i,j}=\binom{n+1}{i}m/d$. Then $\lambda'$ is an isomorphism. That is, \cref{Lazard_thm_lem} holds. 
	\end{prop}
	\begin{proof}
		It suffices to prove that $\lambda'$ is an isomorphism after localizing at every $p$. Therefore, we may assume $M$ is a $\mathbb{Z}_{(p)}$ module. 
	\end{proof}


\section{Complex oriented cohomology theories}

	We begin with the definition of a complex oriented cohomology theory. 
	\begin{defn}
		For a multiplicative cohomology theory represented by a ring spectrum $E$, we say that it is complex orientable if the map $\iota:\mathbb{CP}^1\inj\mathbb{CP}^\infty$ induces an surjection $\iota^*:E^2(\mathbb{CP}^\infty)\to E^2(\mathbb{CP}^1)$. \\
		Recall that $E^n(X)\simeq\tilde{E}^n(X)\oplus E^n(*)$. Then notice that $\tilde{E}^2(\mathbb{CP}^1)=\tilde{E}^2(S^2)=\tilde{E}^0(S^0)=E^0(*)=\pi_0(E)$ has a canonical unit. Thus equivalently, $E$ is complex orientable iff there is a $t\in\tilde{E}^2(\mathbb{CP}^\infty)$ such that $t$ restricts to the canonical unit element of $\tilde{E}^2(S^2)$. We define a complex orientation of $E$ as a choice of $t$. In this case, we say $E$ is complex oriented. 
	\end{defn}
	
	Often, for nice spaces $X$, the existence of a complex orientation of $E$ forces the Atiyah Hirzebruch spectral sequence for $E^*(X)$ to degenerate. 
	\begin{thm}\label{Degen_lemma}
		Let $X$ be a space and assume that each of the homology group $H_n(X;\mathbb{Z})$ is a free abelian group on generators $\{h_{\alpha,n}\}_{\alpha\in B_n}$. Let $c_{\alpha,n}$ be the dual basis in $H^n(X;\mathbb{Z})$. Let $E$ be a ring spectrum and let $\tau_{\leq0}E$ denote its truncation. Then the unit map $S\to E$ induces a map $H\mathbb{Z}\simeq\tau_{\leq0}S\to\tau_{\leq0}E$. Then, the homology classes $h_{n,\alpha}$ have images $h'_{n,\alpha}\in(\tau_{\leq0}E)_n(X)$ and the cohomology classes $c_{n,\alpha}$ have images $c'_{n,\alpha}\in(\tau_{\leq0}E)^n(X)$. Assume that one of the following conditions holds: 
		\begin{enumerate}
			\item Each $h'_{n,\alpha}$ can be lifted to a class $h''_{n,\alpha}\in E_n(X)$ under the map $\tau_{\leq0}E\to E$. 
			\item Each $c'_{n,\alpha}$ can be lifted to a class $c''_{n,\alpha}\in E_n(X)$, and $H_n(X,\mathbb{Z})$ is finitely generated. 
		\end{enumerate}
		Then we have the following: 
		\begin{enumerate}
			\item The smash product $E\wedge\Sigma^\infty X_+$ is equivalent as an $E$-module to a coproduct $\bigvee_{n,\alpha\in B^n}\Sigma^nE$.
			\item The function spectrum $\Map(X,E)$ is equivalent to a product $\prod_{n,\alpha\in B^n}\Sigma^{-n}E$. 
			\item We have isomorphisms $E_*(X)\simeq\pi_*(E)\otimes H_*(X)$ and $E^*(X)\simeq\Hom(H_*(X),\pi_*(E))$. 
		\end{enumerate}
	\end{thm}
	\begin{proof}
		It suffices to prove the first claim, since the second follows from the first from the following argument: 
		\begin{align*}
			&\Map(X,E) \\
			\simeq&\Map_E(X\wedge E,E) \\
			\simeq&\Map_E(\bigvee_{n,\alpha\in B_n}\Sigma^nE,E) \\
			\simeq&\prod_{n,\alpha\in B_n}\Map(\Sigma^nE,E) \\
			\simeq&\prod_{n,\alpha\in B_n}\Sigma^{-n}E 
		\end{align*}
		and the third is a direct consequence of the first two. \\
		We identify $X$ and its suspension spectrum $\Sigma^\infty X_+$. Now, $X$ is connective and has free homology. We construct a sequence of spectra $X_0\to X_1\to\cdots$ having colimit $X$ with the following additional properties: 
		\begin{enumerate}
			\item The map $X_n\to X$ induces an isomorphism in homology of degree $\leq n$. 
			\item $X_n$ is built from finitely many spheres of dimension $\leq n$. This implies that the cohomology $H^i(X_n;\mathbb{Z})$ vanish for $i\geq n$. 
		\end{enumerate}
		We inductively construct $X_n$. We let $X_{-1}=*$. Then suppose $X_{n-1}$ has been constructed. We denote $Y_n$ the cofiber of the map $X_{n-1}\to X$. Thus $Y_n$ is $(n-1)$-connected, and we have $H_n(X)\to H_n(Y_n)$ is an isomorphism. Thus by the Hurewicz theorem, the image of the homology classes $h_{n,\alpha}$ in $H_n(Y_n)$ are represented by maps $S^n\to Y_n$. We let $Z_n=\bigvee_{\alpha\in B_n}S^n$, and let $\varphi_n:Z_n\to Y_n$ be the induced map. Then we let the cofiber of $\varphi_n$ be $W_n$. Finally, we let $X_n$ be the fiber of the composition $X\to Y_n\to W_n$ (or equivalently, let $X_n=X\times_{Y_n}Z_n$). Now the above properties hold obviously ($X$ has homology in all dimensions, $Y_n$ has homology $\geq n$, $Z_n$ has homology concentrated at $n$, hence $W_n$ has homology $>n$). \\
		Now we assume the first assertion holds. We prove the map $\theta:\bigvee_{n,\alpha\in B_n}\Sigma^nE\to E\wedge X$ is $k$-connected for all $k$. Notice that both spectra are connective, so we have $\theta$ 0-connected. Notice by construction that $Y_0=X$. Therefore $\varphi_0$ induces an $E$-module map $\varphi_0\wedge E:\bigvee_{\alpha\in B_0}E\simeq E\wedge Z_0\to E\wedge X$. This map classifies homology classes in $E_0(X)$, say $b_{0,\alpha}$ by definition. Then, since $X$ is connective, we have $(\tau_{\geq1}E)_0(X)=0$, hence the map $E_0(X)\to(\tau_{\leq0}E)_0(X)$ is injective. Therefore $h''_{0,\alpha}$ is uniquely determined in $E_0(X)$. Notice then that $b_{0,\alpha}$ also restricts to $h'_{0,\alpha}$, hence $b_{0,\alpha}=h''_{0,\alpha}$. Thus the following diagram commutes, and the upper and lower sequences are fiber sequences. 
		% https://q.uiver.app/#q=WzAsNixbMCwwLCJcXGJpZ29wbHVzX3tcXGFscGhhXFxpbiBCXzB9RSJdLFsxLDAsIlxcYmlnb3BsdXNfe24sXFxhbHBoYVxcaW4gQl9ufVxcU2lnbWFebkUiXSxbMiwwLCJcXGJpZ29wbHVzX3tuXFxnZXEwLFxcYWxwaGFcXGluIEJfbn1cXFNpZ21hXm5FIl0sWzAsMSwiRVxcb3RpbWVzIFpfMCJdLFsxLDEsIkVcXG90aW1lcyBYXFxzaW1lcSBFXFxvdGltZXMgWV8wIl0sWzIsMSwiRVxcb3RpbWVzIFdfMCJdLFswLDFdLFsxLDJdLFsyLDVdLFsxLDRdLFswLDNdLFszLDRdLFs0LDVdXQ==
		\[\begin{tikzcd}
			{\bigvee_{\alpha\in B_0}E} & {\bigvee_{n,\alpha\in B_n}\Sigma^nE} & {\bigvee_{n\geq0,\alpha\in B_n}\Sigma^nE} \\
			{E\wedge Z_0} & {E\wedge X\simeq E\wedge Y_0} & {E\wedge W_0}
			\arrow[from=1-1, to=1-2]
			\arrow[from=1-1, to=2-1]
			\arrow[from=1-2, to=1-3]
			\arrow[from=1-2, to=2-2]
			\arrow[from=1-3, to=2-3]
			\arrow[from=2-1, to=2-2]
			\arrow[from=2-2, to=2-3]
		\end{tikzcd}\]
		By definition, the left map is a homotopy equivalence, hence to prove the middle map $k$-connective, it suffices to prove the right map is $k$-connective. This is by the inductive hypothesis applied to the connective spectrum $\Sigma^{-1}W_0$. \\
		Now we suppose the second condition holds. We prove the map $E\wedge X\to E\wedge Y_n$ admits a splitting $s_n:E\wedge Y_n\to E\wedge X$. Then since $E\wedge Y_n\to E\wedge X\to E\wedge Y_n$ is the identity, we have $E\wedge X\simeq(E\wedge Y_n)\vee(E\wedge X_{n-1})$. Also, consider the pullback $E\wedge X_n\simeq E\wedge X\times_{E\wedge Y_n}E\wedge Z_n$, which implies $E\wedge X_n\simeq(E\wedge Z_n)\vee(E\wedge X_{n-1})$. Then we have 
		\[ E\wedge X\simeq\ilim(E\wedge X_n)\simeq\ilim\bigvee_{m\leq n,\alpha\in B_m}\Sigma^mE\simeq\bigvee_{n,\alpha\in B_n}\Sigma^nE \]
		We construct the splitting. Notice the cohomology classes $c''_{\alpha,n}$ induces a map 
		\[ \psi_\alpha:Y_n\to E\wedge Y_n\to E\wedge X\to \Sigma^nE \]
		Hence $\psi_\alpha\in E^n(Y_n)$. Notice that again by definition, we have the image of $\psi_\alpha$ in $(\tau_{\leq0}E)^n(Y_n)$ is $c'_{n,\alpha}$. Then $\psi_\alpha$ determines a map $\psi:Y_n\to\bigvee_\alpha\Sigma^nE\simeq E\wedge Z_n$. Moreover, the compatibility with $c'_{n,\alpha}$ shows that the composition $\psi\circ\varphi:E\wedge Z_n\to E\wedge Y_n\to E\wedge Z_n$ is the identity. Thus $Z_n$ splits off in $X$, hence, inductively, $Y_n$ splits off. 
	\end{proof}
	\begin{coro}\label{Coh_CP}
		For any complex oriented ring spectrum $E$, we have $E^*(\mathbb{CP}^\infty)=(\pi_*(E))[[t]]$, where $t$ has degree $2$. 
	\end{coro}
	\begin{proof}
		Consider $E^*(\mathbb{CP}^n)$. Suppose $t\in E^*(\mathbb{CP}^n)$ is the restriction of the complex orientation of $E$ under the inclusion $\mathbb{CP}^n\inj\mathbb{CP}^\infty$, then the cohomology class $\{1,t,\dots,t^n\}$ satisfies the second condition in \cref{Degen_lemma}. Thus we have $\{1,t,\dots,t^n\}$ form a basis for $E^*(\mathbb{CP}^n)$. We then claim that $t^{n+1}=0$, so that $E^*(\mathbb{CP}^n)\simeq(\pi_*(E))[t]/(t^{n+1})$. This is direct by taking the connective cover of $E$. Therefore by letting $\mathbb{CP}^\infty=\ilim\mathbb{CP}^n$, we obtain the result directly (notice the Mittag-Leffler condition holds here so that the $\plim^1$ term vanishes). 
	\end{proof}
	\begin{coro}
		For a complex oriented ring spectrum $E$, we have \\
		$E^*(BU(n))=(\pi_*(E))[[c_1,\dots,c_n]]$, where $c_i$ has degree $2i$. 
	\end{coro}
	\begin{proof}
		We first recall the analogous statement for ordinary cohomology in \cref{Chern_gen}. Then suppose $E_*(BU(1))$ has dual basis $\beta_i$ for $t^i$, and by \cref{Degen_lemma}, we have $E_*(BU(n))$ freely generated by the classes $\{\beta_{i_1}\cdots\beta_{i_n}\}$, where these are the image of $\beta_{i_1}\times\dots\times\beta_{i_n}$ under the map classifying the universal direct sum bundle $\theta:BU(1)^n\to BU(n)$. Then by the diagonal map $BU(n)\to BU(n)^2$, each $E_*(BU(n))$ is a coalgebra. For $n=1$, this coalgebra structure is determined by the graded ring structure on $E^*(BU(1))$, namely, $\beta_n\mapsto\sum_{i+j=n}\beta_i\otimes\beta_j$. Then, $\theta$ determines the coalgebra structure on $E_*(BU(n))$ completely, and the coalgebra structure is given by the same formulas as in the case of integral homology. This implies that the multiplication in cohomology is defined as the same multiplication as in ordinary cohomology. This proves the claim. 
	\end{proof}
	\begin{const}
		For a complex oriented ring spectrum $E$, and an $n$-dimensional complex vector bundle $\zeta$ over a space $X$, we define its $i$-th Chern class ($1\leq i\leq n$) to be the pullback along the classifying map of $\zeta$ of $c_i\in E^{2i}(BU(n))$. The Chern classes satisfies the direct sum formula, as in the classical case. For bundles $\zeta$ and $\zeta'$ over $X$, we have 
		\[ c_n(\zeta\oplus\zeta')=\sum_{i+j=n}c_i(\zeta)c_j(\zeta') \]
	\end{const}

	\begin{const}\label{cplx_ori_FGL}
		Let $E$ be a complex-oriented ring spectrum. Then notice by \cref{Coh_CP}, we have $E^*(\mathbb{CP}^\infty\times\mathbb{CP}^\infty)\simeq(\pi_*(E))[[x,y]]$. Thus let $R=\bigoplus_n\pi_{2n}E$ be a graded ring, then let $\eta^1_1\otimes\eta^1_2$ denote the tensor product of the pullback of the two universal line bundles. This determines a map $\mathbb{CP}^\infty\times\mathbb{CP}^\infty\to\mathbb{CP}^\infty$, which in turn induces a map $(\pi_*(E))[[t]]\to(\pi_*(E))[[x,y]]$, and the image of $t$ is, say $f$. Thus $f$ is a formal group law on $R$, compatible with the grading. \\
		Equivalently, this formal group law characterizes the first Chern class of the tensor product of line bundles. 
	\end{const}

	Interestingly, we have another characterization of the complex orientability of ring spectra. 
	\begin{defn}\label{defn_Thom_class}
		Let $\zeta$ be a vector bundle of rank $n$, and $E$ a ring spectrum. Then we define an $E$-orientation of $\zeta$ is a cohomology class $u\in\tilde{E}^n(\Th(\zeta))\simeq E^n(B(\zeta),S(\zeta))$ (see \cref{Thom_isom}) such that for each point $x\in X$, the restriction $x^*(u)\in E^n(B(\zeta|_x),S(\zeta|x))\simeq\tilde{E}^n(\Th(\zeta|_x))\simeq\tilde{E}^n(S^n)\simeq E^0(*)$ is a generator of $E^*(*)=\pi_*(E)$. Then we say $u$ is a Thom class or fundamental class for $\zeta$. 
	\end{defn}
	\begin{prop}\label{gen_Thom_isom}
		For a ring spectrum $E$ and an $E$-oriented vector bundle $\zeta$, we have $E^*(X)\simeq E^{*+n}(B(\zeta),S(\zeta))$, in which the map is given by multiplication by $u$. Similarly, we have $E_{*+n}(B(\zeta),S(\zeta))\simeq E_*(X)$ given by the slant product with $u$. 
	\end{prop}
	\begin{proof}
		Directly follows from the Atiyah-Hirzebruch spectral sequence of the pair $E^*(B(\zeta),S(\zeta))$, in which the second page isomorphic to the second page of the spectral sequence of $E^*(X)$. 
	\end{proof}
	\begin{thm}
		A ring spectrum $E$ is complex orientable iff all complex vector bundles over any space is $E$-orientable (i.e. has a Thom class). 
	\end{thm}
	\begin{proof}
		It suffices to prove that $E$ is complex orientable iff all the universal bundles are $E$-orientable. \\
		$\implies$: Recall from the proof of \cref{Chern_gen} that $S(\eta^n)\simeq BU(n-1)$. Also recall that $E^*(BU(n))\simeq(\pi_*(E))[[c_1,\dots,c_n]]$. It follows that the map $\iota^*:E^*(BU(n))\to E^*(BU(n-1))$ is surjective with kernel $(c_n)$. Thus $E^*(BU(n),BU(n-1))$ could be identified with the kernel. We claim that $c_n\in E^2n(BU(n),BU(n-1))$ is a Thom class. Since $BU(n)$ is connected, it suffices to check at any point the condition of \cref{defn_Thom_class}. We consider the universal direct sum bundle on $BU(1)^n$. Under the classifying map $\theta$, $c_n\in E^2n(BU(n))$ is mapped to $t_1\cdots t_n\in E^2n(BU(1)^n)\simeq(\pi_*(E))[[t_1,\dots,t_n]]$. Therefore, $c_n\in E^{2n}(BU(n),BU(n-1))$ is mapped to 
		\[ t_1\cdots t_n\in E^2n(B((\eta^1)^n),S((\eta^1)^n))=(t_1\cdots t_n)(\pi_*E)[[t_1,\dots,t_n]] \]
		Now, if each $t_i$ restricts to a unit in $E^2(B(\eta^1|_x),S(\eta^1|_x))$, then clearly $t_1\cdots t_n$ restricts to a unit at any point. Thus it suffices to prove for the claim for $n=1$. But it follows by definition, since $c_1\in\tilde{E}^2(\mathbb{CP}^\infty)$ restricts to $1\in\tilde{E}^2(S^2)=\pi_0E$ by definition. \\
		$\impliedby$: Notice the Thom space of $\eta^1$ is homotopy equivalent to $\mathbb{CP}^\infty$. Therefore the Thom class of $\eta^1$ is the complex orientation. 
	\end{proof}
	\begin{coro}\label{prod_Thom_class}
		For a complex oriented ring spectrum $E$ and complex vector bundles $\zeta$ and $\zeta'$ of dimension $a$ and $b$ on $X$, suppose their Thom classes are given by $u_\zeta\in E^{2a}(B(\zeta),S(\zeta))$ and $u_{\zeta'}\in E^{2b}(B(\zeta'),S(\zeta'))$, we have the Thom class of $\zeta\oplus\zeta'$ given by 
		\[ (u_\zeta,u_{\zeta'})\in \tilde{E}^{2a}(\Th(\zeta))\oplus\tilde{E}^{2b}(\Th(\zeta'))\subset\tilde{E}^{2a+2b}(\Th(\zeta)\vee\Th(\zeta'))=\tilde{E}^{2a+2b}(\Th(\zeta\oplus\zeta')) \]
	\end{coro}
	\begin{proof}
		Since the Thom class is given by the top Chern class, this follows directly from the product formula. 
	\end{proof}

\section{The complex coboridsm spectrum $MU$}

	\begin{defn}
		We define $MU(n)$ to be the spectrum $\Sigma^{-2n}\Th(\eta^n)\simeq\Sigma^{-2n}(BU(n)/BU(n-1))$, where $\eta^n$ is the universal complex $n$-bundle on $BU(n)$, and we identify $\Th(\eta^n)$ with its suspension spectrum. Recall that we have for any real bundle $\zeta$ on $X$, we have $\Sigma\Th(\zeta)\simeq\Th(\zeta\oplus\varepsilon^1)$ where $\varepsilon^1$ is the trivial bundle of dimension $1$. Thus we have a map $MU(n)\to MU(n+1)$ defined by 
		\[ MU(n)\simeq\Sigma^{-2n}\Th(\eta^n)\simeq\Sigma^{-2n-2}\Th(\eta^n\oplus\varepsilon^1_{\mathbb{C}})\to \Sigma^{-2n-2}\Th(\eta^{n+1})=MU(n+1) \]
		We define $MU$ to be the colimit of the sequence $MU(1)\to MU(2)\to\cdots$. 
	\end{defn}

	We have the following basic properties of $MU$. 
	\begin{lemma}
		$MU(0)\simeq S$, the sphere spectrum. $MU(1)\simeq\Sigma^{-2}\mathbb{CP}^\infty$. Moreover, the inclusion $MU(0)\to MU(1)$ is given by 
		\[ S=\Sigma^{-2}\Sigma^2S\inj\Sigma^{-2}\mathbb{CP}^\infty=MU(1) \]
	\end{lemma}
	\begin{proof}
		$MU(0)$ is by definition the suspension spectrum of the Thom space of the trivial bundle on one point, hence is the suspension spectrum of $S^0$. It is easy to see that $\Th(\eta^1)\simeq\mathbb{CP}^\infty$, hence we have the second claim. The third claim is by the construction directly. 
	\end{proof}
	\begin{lemma}
		$MU$ is a ring spectrum. 
	\end{lemma}
	\begin{proof}
		Notice we have a map $m_{a,b}:BU(a)\times BU(b)\to BU(a+b)$ classifying the direct sum bundle. Passing to Thom spaces we have maps $MU(a)\wedge MU(b)\to MU(a+b)$ (notice the shifting cancels out directly). Then by passing to colimits, we obtain a map $MU\wedge MU\to MU$. 
	\end{proof}
	\begin{lemma}\label{structure_map_MU}
		The map $MU(n)\to MU(n+1)$ is given by the following composition.
		\[ MU(n)\simeq MU(n)\wedge MU(0)\to MU(n)\wedge MU(1)\to MU(n+1) \]
		Where the last map is the map induced by $BU(n)\wedge BU(1)\to BU(n+1)$. 
	\end{lemma}
	\begin{proof}
		Follows from the construction. 
	\end{proof}
	\begin{lemma}\label{MU_cplx_oriented}
		$MU$ is complex oriented. 
	\end{lemma}
	\begin{proof}
		Notice the map $\mathbb{CP}^\infty\simeq MU(1)\to MU$ represents a cohomology class $t\in\widetilde{MU}^2(\mathbb{CP}^\infty)$. This map restricts to a unit on $\widetilde{MU}^2(S^2)$ by the natural map $S^0\simeq MU(0)\to MU(1)$, since $MU(0)\to MU$ is the unit of the ring spectrum. 
	\end{proof}

	\begin{remark}
		In fact, we may also define $MU$ using the following method. We let $MU(n)$ denote the Thom space of the universal bundle on $BU(n)$. Then we let $MU$ be the spectrum associated to the prespectrum $MU(0),MU(1),\dots$. 
	\end{remark}
	Having the ``$n$-th space'' $MU(n)\simeq BU(n)/BU(n-1)$, we may see that $MU$ is closely related to orientations on complex vector bundles. In fact, we have the following theorem. 
	\begin{thm}
		$MU$ is the universal complex oriented cohomology theory. That is, for every ring spectrum $E$ with a chosen orientation $t\in\tilde{E}^2(\mathbb{CP}^\infty)$, there exists a unique map $\varphi:MU\to E$, such that for a fixed element $t_{univ}\in\widetilde{MU}^2(\mathbb{CP}^\infty)$, $\varphi_*(t_{univ})=t$. 
	\end{thm}
	\begin{proof}
		We let $E$ be a complex oriented cohomology theory with Thom classes $u_n\in\tilde{E}^{2n}(\Th(\eta^n))$ \\
		$\simeq\tilde{E}^{2n}(BU(n)/BU(n-1))$. Then we have maps $\varphi_n:MU(n)\to E$ classifying $u_n$. We claim the $\varphi_n$ satisfies the following properties. 
		\begin{enumerate}
			\item The map $\varphi_1:MU(1)\to E$ is given by the complex orientation of $E$. 
			\item The following diagram commutes up to homotopy. 
			% https://q.uiver.app/#q=WzAsNCxbMCwwLCJNVShhKVxcd2VkZ2UgTVUoYikiXSxbMSwwLCJNVShhK2IpIl0sWzAsMSwiRVxcd2VkZ2UgRSJdLFsxLDEsIkUiXSxbMCwxXSxbMSwzLCJcXHZhcnBoaV97YStifSJdLFswLDIsIlxcdmFycGhpX2FcXHdlZGdlXFx2YXJwaGlfYiIsMl0sWzIsM11d
			\[\begin{tikzcd}
				{MU(a)\wedge MU(b)} & {MU(a+b)} \\
				{E\wedge E} & E
				\arrow[from=1-1, to=1-2]
				\arrow["{\varphi_a\wedge\varphi_b}"', from=1-1, to=2-1]
				\arrow["{\varphi_{a+b}}", from=1-2, to=2-2]
				\arrow[from=2-1, to=2-2]
			\end{tikzcd}\]
			\item The map $MU(n)\to MU(n+1)\to E$ coincides with $\varphi_n$. 
		\end{enumerate}
		The first claim is by definition. The second claim is by \cref{prod_Thom_class}. For the third claim, we consider \cref{structure_map_MU}, and obtain that the composition $MU(n)\to MU(n+1)\to E$ is given by 
		\[ MU(n)\simeq MU(n)\wedge MU(0)\to MU(n)\wedge MU(1)\to E \]
		where the last map is given by the smash product of $\varphi_n$ and $\varphi_1$ followed by the product map of $E$. Thus this is equal to $\varphi_n$ by the second claim and that $\varphi_1$ restricts to $\varphi_0$. \\
		By definition, $\Map(MU,E)\simeq\Map(\ilim MU(n),E)\simeq\plim\Map(MU(n),E)$. Therefore we have a Milnor long exact sequence
		\[ \plim{^1}E^{-1}(MU(n))\to E^0(MU)\to\plim E^0(MU(n))\to\plim{^1}E^0(MU(n)) \]
		The $\plim^1$ terms vanish by the Mittag-Leffler condition (notice $E^*(MU(n+1))\to E^*(MU(n))$ is surjective since it is the process of killing $c_{n+1}$ under a shift). Therefore $\varphi_n$ uniquely determine $\varphi:MU\to E$. \\
		We then show that $\varphi$ is a map of ring spectra, that is, the following diagram commutes. 
		% https://q.uiver.app/#q=WzAsNCxbMCwwLCJNVVxcd2VkZ2UgTVUiXSxbMSwwLCJNVSJdLFswLDEsIkVcXHdlZGdlIEUiXSxbMSwxLCJFIl0sWzAsMV0sWzEsM10sWzAsMl0sWzIsM11d
		\[\begin{tikzcd}
			{MU\wedge MU} & MU \\
			{E\wedge E} & E
			\arrow[from=1-1, to=1-2]
			\arrow[from=1-1, to=2-1]
			\arrow[from=1-2, to=2-2]
			\arrow[from=2-1, to=2-2]
		\end{tikzcd}\]
		But since we have $E^0(MU)\simeq\plim E^0(MU(n))$, this diagram commutes iff all diagrams in claim (2) commutes. Thus $\varphi$ is a map of ring spectra. \\
		Finally we are left to show that $\varphi$ carries the complex orientation of $MU$ to the complex orientation of $E$. However this is by definition of $\varphi_1$. 
	\end{proof}

	Now, as we have shown in \cref{cplx_ori_FGL}, a complex orientation on $E$ induces a formal group law on $\bigoplus\pi_{2n}(E)$. As the following theorem presents, the universal complex oriented spectrum classifies the universal formal group law. 
	\begin{thm}\label{Quillen_thm}
		The map $\varphi:L\to\pi_*(MU)$ classifying is an isomorphism. Here $L$ is the Lazard ring. 
	\end{thm}
	This will be the main topic of the two following sections. But now, we first calculate the homology of $MU$, which is much easier. 
	\begin{thm}
		For any complex oriented cohomology theory $E$, we have $E_*(MU)\simeq$ \\
		$(\pi_*(E))[b_1,b_2,\dots]$, where $b_i$ has degree $2i$. 
	\end{thm}
	\begin{proof}
		Since $MU=\ilim MU(n)$, we have $E_*(MU)=\ilim E_*(MU(n))$. By \cref{gen_Thom_isom}, we have $E_*(MU(n))\simeq E_*(BU(n))$. Then recall $E_*(BU(n))\simeq\mathrm{Sym}^n(\pi_*(E){\beta_0,\beta_1,\dots})$, where $\beta_i$ is the dual basis of $t^i$ for $E^*(BU(1))$. Then we have $E_*(MU(n))\simeq\mathrm{Sym}^n(\pi_*(E)){b_0,b_1,\dots}$, where $b_i$ is the dual basis of $t^{i+1}$. Notice that the map $E_*(MU(n))\to E_*(MU(n+1))$ is given by multiplying with $b_0$. Then finally notice that we may identify $\mathrm{Sym}^n(\pi_*(E)){b_0,b_1,\dots}$ with $\bigoplus{i\leq n}\mathrm{Sym}^i(\pi_*(E)){b_1,\dots}$ by mapping $b_0$ to 1. Then by passing to colimits we obtain $E_*(MU)\simeq(\pi_*(E))[b_1,b_2,\dots]$. 
	\end{proof}
	\begin{remark}
		This is analogous to the calculation of $H_*(MO)$. 
	\end{remark}

	In particular, we have $H_*(MO)\simeq L$, the Lazard ring. 
	\begin{const}
		Notice that for any complex oriented cohomology theory $E$, the smash product $MU\wedge E$ has two complex orientations, associated to the pushforward of the orientation of $MU$ and $E$ respectively. Notice $\pi_*(MU\wedge E)\simeq E_*(MU)\simeq(\pi_*(E))[b_1,b_2,\dots]$. Explicitly, we have two classes $t_{MU},t_E\in\widetilde{MU\wedge E}^2(\mathbb{CP}^\infty)$ such that
		\[ (\pi_*(E))[b_1,b_2,\dots][[t_E]]\simeq(MU\wedge E)^*(\mathbb{CP}^\infty)\simeq(\pi_*(E))[b_1,b_2,\dots][[t_{MU}]] \]
		Therefore we may express $t_{MU}$ as a power series $t_{MU}=\sum_{i\geq 1}a_it_E^{i}$. 
	\end{const}
	\begin{prop}
		We have $t_{MU}=t_E+b_1t_E^2+\cdots$. 
	\end{prop}
	\begin{proof}
		We first recall that a class in $\widetilde{MU\wedge E}^2(\mathbb{CP}^\infty)$ can be represented by a map $MU(1)\to MU\wedge E$. Equivalently, we may consider a $E$-module map $MU(1)\wedge E\to MU\wedge E$. We suppose $t_MU$ and $t_E$ are represented by maps $\varphi_{MU}$ and $\varphi_E$. Then, notice $b_i\in E_{2i}(MU(1))$ determines a map $\Sigma^{2i}E\to MU(1)\wedge E$. By \cref{Degen_lemma}, we have an equivalence of $E$-module spectra $\bigvee_{i\geq0}\Sigma^{2i}E\simeq MU(1)\wedge E$. Thus, to specify a map from $MU(1)\wedge E$ to $MU\wedge E$, it suffices to specify its restriction to $\Sigma^{2i}E$ for all $i$. \\
		Let the map $\lambda:E\wedge MU(1)\to E$ be the complex orientation on $E$, then we have $\varphi_E$ the composition $E\wedge MU(1)\to E\to E\wedge MU$. By construction, $\lambda$ vanishes on $\Sigma^{2i}E$ for $i>0$ and restricts to the identity when $i=0$. Then, the map $\varphi_{MU}$ is given by smashing with $E$ the map $MU(1)\to MU$. Thus $\varphi_{MU}$ is the coproduct of the maps $\varphi_{MU}^i:\Sigma^{2i}E\to MU\wedge E$ classified by $b_i$. \\
		Finally, notice that the complex orientation $t_E$ induces a map $\Sigma^{-2}(MU(1)\wedge E)\to MU(1)\wedge E$ by the slant product. This action translates $MU(1)\otimes E$ by one degree downward. Therefore $\varphi_{MU}^i$ can be identified with the map translated downward $i$ times composed with $b_i\varphi_E$. Finally, we return to $t_{MU}$ and $t_E$ and obtain $t_{MU}=\sum b_it_E^{i+1}$. 
	\end{proof}

	Then, recall from \cref{cplx_ori_FGL} that every complex orientation induces a formal group law on $R=\pi_*(MU\wedge E)=(\pi_*(E))[b_1,b_2,\dots]$. By the above proposition we clearly have the following. 
	\begin{coro}
		Let $R=(\pi_*(E))[b_1,b_2,\dots]$, and $f_E$, $f_{MU}$ be the formal group laws associated to $t_E$ and $t_{MU}$. Then let $g(x)\in R[[x]]$ be $x+b_1x^2+\cdots$, we have $t_{MU}=t_E$. Hence we have 
		\[ f_{MU}(x,y)=g\circ f_E(g^{-1}(x),g^{-1}(y)) \]
	\end{coro}
	Then, by restricting to $H\mathbb{Z}$, we get the following. 
	\begin{coro}
		Let $E=MU\wedge H\mathbb{Z}$, equipped with the complex orientation coming from $MU$. Then $\pi_*(E)=\mathbb{Z}[b_1,\dots]$, and the formal group law is given by $f(x,y)=g(g^{-1}(x)+g^{-1}(y))$, where $g=x+b_1x^2+\cdots$. 
	\end{coro}
	\begin{proof}
		Recall the formal group law on $H\mathbb{Z}$ is given by addition. 
	\end{proof}
	Notice the map $L\to H_*(MU;\mathbb{Z})$ is exactly the same as in \cref{Lazard_thm}, thus by \cref{Lazard_thm}, we have the composition $L\to\pi_*(MU)\to H_*(MU;\mathbb{Z})$ is a rational isomorphism, where the second map is the Hurewicz map. Recall that the Hurewicz map is always an rational isomorphism, hence we have the following. 
	\begin{coro}
		The map $L\otimes\mathbb{Q}\to\pi_*(MU)\otimes\mathbb{Q}$ is an isomorphism.  
	\end{coro}
	Thus to prove the theorem $L\simeq\pi_*(MU)$, it suffices to prove that they are isomorphic after $p$-adic completion for all $p$ (this is an purely algebraic result). We will use the Adams spectral sequence for $MU$ to prove this result. 

\section{The Adams spectral sequence for $MU$}
	
	We first recall some basic facts about the Adams spectral sequence. 
	\begin{const}
		Let $R=H\mathbb{F}_p$ for some prime $p$. Then we have a cosimplicial object $R^\bullet$ in $\Sp$, given by 
		\[ [n]\mapsto R\wedge\cdots\wedge R=R^{\wedge n+1} \]
		and the boundary and degeneracy maps the obvious ones. Moreover, we may extend this cosimplicial object to a augmented simplicial object $\overline{R}^\bullet$ by letting $\overline{R}^{-1}=S$, the sphere spectrum. For any other spectrum $X$, we let $\overline{X}^\bullet=X\wedge\overline{R}^\bullet$, and $X^\bullet=X\wedge R^\bullet$. Then, we have a map $X=\overline{X}^{-1}\to\plim X^\bullet$. Finally, we let $E_1^{a,b}=\pi_a(X^b)=\pi_a(X\wedge R^{\wedge b+1})$, with the differential given on the first page to be 
		\[ \pi_a(X^0)\xrightarrow{d^0-d^1}\pi_a(X^1)\xrightarrow{d^0-d^1+d^2}\cdots \]
		For simplicity, we denote $\plim X^\bullet$ by $\Tot(X^\bullet)$ and the partial limit $\plim_{1\leq i\leq n}X^\bullet$ by $\Tot_n(X^\bullet)$. Notice that in general context, the totalization of a simplicial object is defined via ends, and the author is not clear of the difference of the two notions. \\
		We have the following theorem. 
	\end{const}
	\begin{thm}\label{Adams_spec_seq}
		Suppose $X$ is a connective spectrum with $\pi_n(X)$ finitely generated for all $n$. Then the map $X\to \Tot X^\bullet$ has the following properties. 
		\begin{enumerate}
			\item For all $n$, we have the following filtration
			\[ \cdots\subseteq F^2\pi_n(X)\subseteq F^1\pi_n(X)\subseteq F^0\pi_n(X)\simeq\pi_n(X) \]
			Where $F^i\pi_n(X)$ is the kernel of $\pi_n(X)\to\pi_n(\Tot_i X^\bullet)$. 
			\item For each $i\geq0$, there is some $j>i$ such that $(p^i)\pi_n(X)\subseteq F^j\pi_n(X)\subset (p^j)\pi_n(X)$. In particular, we have $\plim(\pi_n(X)/F^j\pi_n(X))\simeq(\pi_n(X))_p$, the $p$-adic completion. 
			\item As in the case for general spectral sequences, for fixed $a$ and $b$, we have $E^{a,b}_r$ stablize to $E^{a,b}_\infty$ for $r\gg0$. Moreover, we have $E^{a,b}_\infty\simeq F^b\pi_{a+b}(X)/F^{b+1}\pi_{a+b}(X)$. 
			If $X$ is a ring spectrum, we have the following additional properties. 
			\item The filtration agrees with the multiplicative structure on $\pi_*(X)$. Explicitly, the multiplication map $\pi_m(X)\times\pi_n(X)\to\pi_{m+n}(X)$ restricts to a map $F^i\pi_m(X)\times F^j\pi_n(X)\to F^{i+j}\pi_{m+n}(X)$. 
			\item The spectral sequence is multiplicative. That is, for each $r$ we have maps $E^{a,b}_r\times E^{a',b'}_r\to E^{a+a',b+b'}_r$. These maps are associative and satisfies the Leibniz rule when applied the differential. Finally, the map $E^{a,b}_\infty\times E^{a',b'}_\infty\to E^{a+a',b+b'}_\infty$ agrees with the quotient map of $F^b\pi_{a+b}(X)\times F^{b'}\pi_{a'+b'}(X)\to F^{b+b'}\pi_{a+b+a'+b'}(X)$. 
		\end{enumerate}
	\end{thm}
	\begin{const}
		Again let $R=H\mathbb{F}_p$. We recall that the Steenrod algebra is the Hopf algbera $\mathcal{A}\simeq H^*(R;\mathbb{F}_p)$. Then we have the dual Steenrod algebra $\mathcal{A}_*\simeq\pi_*(R\wedge R)$. \\
		Also recall that we define for a ring spectrum $E$ and module spectrums $X$ and $Y$, we have the definition 
		% https://q.uiver.app/#q=WzAsMyxbMCwwLCJYXFx3ZWRnZSBZIl0sWzEsMCwiWFxcd2VkZ2UgRVxcd2VkZ2UgWSJdLFsyLDAsIlxcY2RvdHMiXSxbMCwxXSxbMSwwLCIiLDAseyJvZmZzZXQiOi0yfV0sWzEsMCwiIiwxLHsib2Zmc2V0IjoyfV0sWzIsMV0sWzIsMSwiIiwxLHsib2Zmc2V0IjotNH1dLFsyLDEsIiIsMSx7Im9mZnNldCI6NH1dLFsxLDIsIiIsMSx7Im9mZnNldCI6Mn1dLFsxLDIsIiIsMSx7Im9mZnNldCI6LTJ9XV0=
		\[ X\wedge_EY=\ilim\left( 
		\begin{tikzcd}
			{X\wedge Y} & {X\wedge E\wedge Y} & \cdots
			\arrow[from=1-1, to=1-2]
			\arrow[shift left=2, from=1-2, to=1-1]
			\arrow[shift right=2, from=1-2, to=1-1]
			\arrow[shift right=2, from=1-2, to=1-3]
			\arrow[shift left=2, from=1-2, to=1-3]
			\arrow[from=1-3, to=1-2]
			\arrow[shift left=4, from=1-3, to=1-2]
			\arrow[shift right=4, from=1-3, to=1-2]
		\end{tikzcd} \right) \]
		Then, notice that we have 
		\[ X^b=X\wedge R^{\wedge b+1}=(X\wedge R)\wedge_R(R\wedge R)\wedge_R\cdots \]
		Therefore, we have $E_1^{a,b}\simeq H_a(X;\mathbb{F}_p)\otimes_{\mathbb{F}_p}(\mathcal{A}_*)^{\otimes b}$, and the differentials is given by the coalgebra structure of $\mathcal{A}_*$, and the comodule structure of $H_*(X;\mathbb{F}_p)$. 
	\end{const}

	\begin{const}
		From the above construction, we fix $p=2$ for convenience in applying algebro-geometric languages. let $\mathbb{G}$ denote the spectrum of the commutative ring $\mathcal{A}_*$, and the Hopf algebra structure on $\mathcal{A}_*$ induces a group scheme structure on $\mathbb{G}$. Therefore, we have $\mathbb{G}$ a group scheme over $\Spec\mathbb{F}_2$. Then let $V$ denote the $\mathbb{F}_2$-vector space $H_*(X;\mathbb{F}_2)$. The coaction $V\to V\otimes\mathcal{A}_*$ induces a structure of a representation of the algebraic group $\mathbb{G}$ on $V$. In this term, the $E_1^{*,*}$ term is simply given by the cohomology of the complex (with each term graded) 
		\[ V\to\Gamma(\mathbb{G};V\otimes\mathcal{O}_\mathbb{G})\to\Gamma(\mathbb{G}^2;V\otimes\mathcal{O}_{\mathbb{G}^2})\to\cdots \]
		And therefore the $E_2^{*,b}$ term can be identified with $H^b(\mathbb{G};H_*(X;\mathbb{F}_2))$. \\
		If $X$ is a ring spectrum, then $H_*(X;\mathbb{F}_2)$ is itself a commutative $\mathbb{F}_2$-algebra, hence we may form the spectrum $Y=\Spec(H_*(X;\mathbb{F}_2))$. Therefore the action of $\mathbb{G}$ on $Y$ gives a quotient stack $Y/\mathbb{G}$, whose cohomology, by definition, can be identified with the $E_2$ page. Therefore we simply have $E_2^{*,*}=H^*(Y/\mathbb{G};\mathcal{O}_{Y/\mathbb{G}})$ (with natural grading on $Y$). 
	\end{const}

	\begin{const}
		We let $X(n)=\Map(\mathbb{RP}^n,S)$. Then $X(n)$ is an $\EE_\infty$ ring spectrum, and we have $H_*(X(n);\mathbb{F}_2)\simeq H^*(\mathbb{RP}^n;\mathbb{F}_2)\simeq\mathbb{F}_2[x]/(x^{n+1})$. Then we pass through the limit and get $\ilim\Spec H_*(X(n);\mathbb{F}_2)\simeq\Spf H^*(\mathbb{RP}^\infty;\mathbb{F}_2)=\mathbb{F}_2[[x]]$. Then we naturally have an action of $\mathbb{G}$ on $\Spf H^*(\mathbb{RP}^\infty;\mathbb{F}_2)$. Notice that the map $\mathbb{RP}^\infty\times\mathbb{RP}^\infty\to\mathbb{RP}^\infty$ is the colimit of the maps $\mathbb{RP}^m\times\mathbb{RP}^n\to\mathbb{RP}^{m+n}$, and hence this induces a map 
		\[ \Spf\mathbb{F}_2[[x]]\times_{\Spec\mathbb{F}_2}\Spf\mathbb{F}_2[[x]]\to\Spf\mathbb{F}_2[[x]] \]
		given by a power series $f(x,y)=x+y$. By definition, $f$ is a formal group law over $\mathbb{F}_2$, and $\mathbb{G}$ respects the group structure given by $f$. Therefore the $\mathbb{G}$ action on $\Spf\mathbb{F}_2\times_{\Spec\mathbb{F}_2}\Spec A=\Spf A[[x]]$ can be viewed as an automorphism of the formal group $x+y$. Therefore, the automorphism must be of the form $x\mapsto x+a_1x^2+a_2x^4+a_3x^8+\cdots$ since they are the only series preserving addition in characteristic 2. Recall also the computation of the dual Steenrod algebra \cref{Struct_A_*}. This implies that $\mathbb{G}$, viewed as a functor $\Ring\to\Set$, is given by $\Hom(\Spec A,\mathbb{G})$ to be all power series
		\[ x\mapsto x+a_1x^2+a_2x^4+a_3x^8+\cdots \]
		Now, we apply the above constructions to the case where $X=MU$. Recall that we have $H_*(MU;\mathbb{F}_2)=\mathbb{F}_2[b_i]$ for $b_i$ of degree $2i$. We let $Z$ denote the spectrum $\Spec H_*(MU;\mathbb{F}_2)$, then for any $\mathbb{F}_2$ algebra $R$, the map $\Hom(\Spec R,Z)$ can be identified with the set consisting of all power series $y+b_1y^2+b_2y^3+\cdots$. Therefore we see that $Z$ classifies the coordinates of the additive formal group on $R$ preserving the first order terms. By the description of $\mathbb{G}$ as the group of automorphism of the additive formal group, we have a natural action of $\mathbb{G}$ on $Z$. 
	\end{const}

	However, we notice that this action is different to the desired action in the Adams spectral sequence of $MU$, due to the problems of cohomological gradings. We now construct the correct action. 

	\begin{const}
		For a complex oriented ring spectrum $E$ where $\pi_*(E)$ is a $\mathbb{F}_2$ algebra, notice that $\Hom(\mathcal{A}_*,\pi_*(E))$ (respect gradings) is the group of automorphisms on the additive formal group on $\pi_*(E)$. But the ring $\Hom(H_*(MU;\mathbb{F}_2),\pi_*(E))$ is the coordinates of the additive formal group on $\bigoplus_{i}\pi_{2i}(E)$. This implies that the correct action of $\mathbb{G}$ on $Z$ should differ by our construction with a Frobenius map $F:\mathbb{G}\to\mathbb{G}$. Explicitly, consider $\mathbb{G}'\subseteq\mathbb{G}$ be the image of the Frobenius map. Then for any ring $R$, we define the $R$-point $f=x+a_1x^2+\cdots$ of $\mathbb{G}'$ acts on the $R$-points of $Z$ by $x\mapsto x+a_1^2x^2+\cdots$. Finally we define the action of $\mathbb{G}$ on $Z$ to be the action factoring through the Frobenius map. 
	\end{const} 
	\begin{lemma}
		The action described above is indeed the action induced by the $\mathcal{A}_*$ coalgebra structure of $H_*(MU;\mathbb{F}_2)$. 
	\end{lemma}
	\begin{proof}
		This can be done by direct computation of the coalgebra structure of $H_*(\mathbb{CP}^\infty;\mathbb{F}_2)$ (if someone sees this, remind me to finish the computation). 
	\end{proof}

	We first calculate the cohomology of $Z/\mathbb{G}'$. 
	\begin{lemma}
		Let $Z_0=\Spec\mathbb{F}_2[b_2,b_4,\dots]$ (with generators of degree not of the form $2^r-1$), then the action of $\mathbb{G}'$ on $Z$ determines an isomorphism $a:\mathbb{G}'\times Z_0=Z$, that is, we have an isomorphism of schemes $Z/\mathbb{G}'=Z_0$. 
	\end{lemma}
	\begin{proof}
		It suffices to prove that for all $R$, we have the following: for all $R$-points of $Z$, say $h(x)=x+c_1x^2+\cdots$, $h$ can be uniquely factored into $f\circ g$, where $f$ is of the form $f(x)=x+a_1x^2+a_2x^4+\cdots$ and $g(x)=x+b_2x^3+b_4x^5+\cdots$. This is obvious.  
	\end{proof}
	\todo{Correct Lurie!!!}

	Finally we see the following results. 
	\begin{thm}
		We have the second page of the mod-2 Adams spectral sequence of $MU$ given by 
		\[ E_2^{*,*}\simeq\mathbb{F}_2[b_2,b_4,\dots,\varepsilon_1,\varepsilon_2,\dots] \]
		where $b_i$ ($i\neq 2^r-1$) has bidegree $(2i,0)$, and $\varepsilon_j$ has bidegree $(2^j-1,1)$. 
	\end{thm}
	Therefore, the Adams spectral sequence degenerates since there is no room for non-trivial maps. 
	The mod-$p$ Adams spectral sequence behaves exactly the same as the mod-2 case, namely, we have
	\begin{thm}
		We have the second page of the mod-$p$ (for all prime $p$) Adams spectral sequence of $MU$ given by 
		\[ E_2^{*,*}\simeq\mathbb{F}_p[c_1,c_2,\dots] \]
		where $c_i=\varepsilon_j$ if $i+1=p^j$, and $c_i=b_i$ otherwise (again $b_i$ has bidegree $(2i,0)$ and $\varepsilon_j$ has bidegree $(p^j-1,1)$). 
	\end{thm}
	\begin{proof}
		By essentially the same proof, but with super algebraic geometry instead since for odd $p$, $\mathcal{A}_*$ is supercommutative (maybe I will add an appendix). 
	\end{proof}
	
	Therefore, by \cref{Adams_spec_seq}, we deduce the following result. 
	\begin{coro}
		The associated graded ring $\gr(\pi_*(MU))$ of the mod-$p$ Adams filtration on $\pi_*(MU)$ is isomorphic to a polynomial algebra $\mathbb{F}_p[c_0,\dots]$, where $c_i\in\gr_0(\pi_{2i}(MU))$ for $i+1\neq p^j$ and $c_i\in\gr_1(\pi_{2i}(MU))$ for $i+1=p^j$. 
	\end{coro}
	\begin{coro}\label{p_adic_cmplt_piMU}
		$(\pi_*MU)^\vee=\mathbb{Z}_p[u_1,u_2,\dots]$. Here $\mathbb{Z}_p$ denotes the $p$-adic completion of $\mathbb{Z}$. 
	\end{coro}

	We are finally ready to prove \cref{Quillen_thm}. 
	\begin{proof}
		Recall that we have by \cref{Lazard_thm_lem}, the map $L\to\pi_*(MU)\to H_*(MU;\mathbb{Z})$ is injective. Thus it suffices to show that the map $L\to\pi_*(MU)$ is surjective. In fact, we show that it is surjective after $p$-adic completion for all $p$. By \cref{Lazard_thm}, we have $L^\vee\simeq\mathbb{Z}_p[t_1,t_2,\dots]$, and we have by \cref{p_adic_cmplt_piMU} that $(\pi_*(MU))^\vee\simeq\mathbb{Z}_p[u_1,u_2,\dots]$. Let $I$, $J$ be defined as in \cref{Lazard_thm_lem}, and $K$ the likewise ideal for $(\pi_*(MU))^\vee$. As in the proof of \cref{Lazard_thm}, it suffices to prove the map $I/I^2\to K/K^2$ is surjective in each degree. \\
		Recall that $(I/I^2)_{2n}=\mathbb{Z}_pt_n$ and $(K/K^2)_{2n}=\mathbb{Z}u_n$, hence $\theta(t_n)=\lambda u_n+c$, where $c$ is decomposable in $(\pi_*(MU))^\vee$. It suffices to prove that $\lambda$ is a $p$-adic unit. But then, the Hurewicz map maps $u_n$ to $\lambda'b_n+d$, where $d$ is again decomposable. Thus by \cref{Lazard_thm}, $\lambda\lambda'$ is $p$ if $n+1=p^r$ and $1$ otherwise. If $n+1$ is not a power of $p$, then $\lambda$ is obviously a unit. If $n+1=p^r$, it suffices to prove $\lambda'$ is divisable by $p$, which is equivalent to the image of $u_n$ vanishes in the ring $H_*(MU;\mathbb{F}_p)$. This is equivalent to the requirement that $u_n$ has Adams filtration degree $\geq1$, which is true. 
	\end{proof}

	Finally, we return to the Adams spectral sequence. 
	\begin{const}
		Recall in \cref{Adams_spec_seq}, $R=H\mathbb{F}_p$ for some $p$. We now let $R=MU$. By the same construction we obtain a cosimplicial spectrum $X^\bullet$. Now, for connective $X$, the map $X\to\plim X^\bullet\simeq\plim_n(\plim_{i\leq n}X^i)$ is an equivalence. We call the spectral sequence associated to this cosimplicial spectrum the Adams-Novikov spectral sequence. \\
		Since $\pi_0(MU)=\pi_0(S)=Z$, the convergence is faster than the mod-$p$ Adams spectral sequence. In fact, if we denote $F^n\pi_*X$ to be the kernel of $\pi_*X\to\pi_*\plim_{i\leq n+1}X^\bullet$, for each $n$ we have a finite filtration
		\[ 0=F^{n+1}\pi_n(X)\subseteq F^n\pi_n(X)\subseteq\cdots\subseteq F^0\pi_n(X)=\pi_n(X) \]
	\end{const}

	\begin{const}
		Again, like the mod-$p$ Adams spectral sequence, the chain complex 
		\[ MU_*(X)\to (MU\wedge MU)_*(X)\to (MU\wedge MU\wedge MU)_*X\to\cdots \]
		has an algebro-geometric interpretation. Recall $\pi_*(E\wedge MU)=E_*(MU)=\pi_*(E)[b_1,b_2,\dots]$, for all complex oriented cohomology theory $E$. Then, letting $\pi_*(E)=R$, $\Spec\pi_*(E\wedge MU)$ is an infinite dimensional affine space over $\Spec R$. Explicitly, it is the affine space parametrizing all coordinates $g(t)=t+b_1t^2+\cdots$ on $R[[t]]$ agreeing with the standard coordinate to the first order. \\
		Then, we let $G=\Spec\mathbb{Z}[b_1,b_2,\dots]$. This is a group scheme, whose group structure on $R$-points (namely, power series of the form $g(t)=t+b_1t^2+\cdots$) is given by composition of power series. Therefore, we have $\Spec\pi_*(E\wedge MU)=G\times\Spec\pi_*(E)$. In particular, we have $\Spec\pi_*(MU\wedge MU)=G\times\Spec L$. The ring structure of $MU$ encodes a $G$ action on $\Spec L$, namely, a power series $g$ acts on $\Spec L$ by sending the formal group law $f(x,y)$ to $g(f(g^{-1}(x),g^{-1}(y)))$. 
	\end{const}
	
	\begin{defn}\label{Mod_stack_FGLs}
		We define $\mathcal{M}_{FG}^s$ to be the quotient stack $\Spec L/G$. Equivalently, this is the moduli stack sending a ring $R$ to the groupoid of all formal group laws with strict isomorphisms. Note that this is indeed a stack w.r.t the flat topology. 
	\end{defn}

	\begin{const}
		Let $X$ be any spectrum. Then, $\pi_*(X^n)=\pi_*(X\wedge MU^{\wedge n+1})$ is a module over the commutative ring $\pi_*(MU^{\wedge n+1})$, hence can be identified with a quasi-coherent sheaf on $\Spec L\times G^n$. Notice that the quasi-coherent sheaves $\pi_*(X^n)$ are compatible with one another, hence they form a quasi-coherent sheaf over $\mathcal{M}_{FG}^s$. We call this quasi-coherent sheaf $\mathcal{F}_X$. Moreover, the cochain complex 
		\[ MU_*(X)\to(MU\wedge MU)_*(X)\to\cdots \]
		is the standard complex for computing the cohomology of $\mathcal{M}_{FG}^s$ with $\mathcal{F}_X$ coefficients. Thus we have the following. 
	\end{const}
	\begin{thm}\label{Sec_page_Adams_Nov}
		For any spectrum $X$, the second page of the Adams-Novikov spectral sequence of $X$ is given by 
		\[ E_2^{*,b}\simeq H^b(\mathcal{M}_{FG}^s;\mathcal{F}_X)=H^b(G;MU_*(X)) \]
	\end{thm}
	We will focus on this point of view in the next chapter. 

\chapter{The moduli stack of formal groups}

\section{The moduli stack $\mathcal{M}_{FG}$}

	\begin{defn}
		Besides the $G$ action on $\Spec L$, we have an action of the multiplicative group scheme $\mathbb{G}_m$ on $\Spec L$, given by $f(x,y)\mapsto \lambda f(\lambda^{-1}x,\lambda^{-1}y)$. We define the group scheme $G^+$ to be the semidirect product of $G$ and $\mathbb{G}_m$, whose $R$-points are given by series $g(x)=b_0x+b_1x^2+\cdots$, where $\lambda$ is invertible. $g$ has an action on the formal group laws, given by $g(f(g^{-1}(x),g^{-1}(y)))$. These kinds of isomorphisms between formal group laws are called isomorphisms (in contrast to strict isomorphisms), and we define $\mathcal{M}_{FG}=\Spec L/G^+$ to be the moduli stack of formal group laws with isomorphisms. 
	\end{defn}

	Recall from the appendix that a formal group is simply a group object in the category of formal schemes. Therefore, we have the following definition. 
	\begin{defn}
		For a formal group law $f\in R[[x,y]]$, we define $\mathcal{G}_f$ to be the formal group over $\Spec R$ whose $A$-points (here $A$ is an $R$-algebra) is given by 
		\[ \mathcal{G}_f(A)=\{ a\in A|\exists n,a^n=0 \} \]
		and the group structure is given by $a\cdot b=f(a,b)$. We call $\mathcal{G}_f$ the formal group associated to $f$. \\
		We define a formal group over $R$ to be coordinatizable, iff it is of the form $\mathcal{G}_f$ for some $f$. 
	\end{defn}

	\begin{prop}
		Isomorphisms between coordinatizable formal groups are given by isomorphisms between the corresponding formal group laws. 
	\end{prop}
	\begin{proof}
		Suppose we have $\alpha:\mathcal{G}_f\simeq\mathcal{G}_{f'}$ an isomorphism. Then, consider $\alpha_n=\alpha(R[t]/(t^{n+1}))$. Then $\alpha_n$ induces a map from nilpotent elements in $R[t]/(t^{n+1})$ to itself. We consider the element $\alpha_n(t)=b_0t+b_1t^2+\cdots+b_{n-1}t^n$. Then by functorality of $\alpha$ and $\mathcal{G}_f$, we have $\alpha_{n+1}(t)=b_0t+\cdots+b_nt^{n+1}$. These elements put together to a formal power series $g(t)=b_0t+\cdots$. This is the isomorphism between $f$ and $f'$. 
	\end{proof}

	It is now important to make precise what the moduli stacks $\mathcal{M}_{FG}^s$ and $\mathcal{M}_{FG}$ represents. The main obstruction is that the functor assigning a ring $R$ to its coordinatizable formal group laws with isomorphisms does not satisfy descent. Therefore we need the following definition. 

	\begin{defn}\label{Defn_FG}
		For a ring $R$, we define a formal group over $R$ to be a functor $\mathcal{G}:\Alg_R\to\Ab$ satisfying the two following properties. 
		\begin{enumerate}
			\item $\mathcal{G}$ is a Zariski sheaf. 
			\item $\mathcal{G}$ is Zariski locally isomorphic to coordinatizable formal groups. 
		\end{enumerate}
	\end{defn}
	\begin{thm}
		$\mathcal{M}_{FG}$ sends each ring $R$ to the groupoid of formal groups in the sense of \cref{Defn_FG} with isomorphisms as morphisms. 
	\end{thm}
	\begin{proof}
		This can be viewed as a stackification. 
		\todo{Make this clear!}
	\end{proof}


	\begin{defn}
		For a ring $R$ and a formal group $\mathcal{G}$ over $R$, we define its Lie algebra $\mathfrak{g}$ to be the abelian group $\ker(\mathcal{G}(R[t]/(t^2))\to\mathcal{G}(R))$. In fact, it is an $R$-module, with $\lambda\in R$ acting on $\mathfrak{g}$ by the endomorphism on $R[t]/(t^2)$ given by $t\mapsto \lambda t$. 
	\end{defn}
	
	\begin{lemma}\label{Coord_FGL_endo}
		For a coordinatizable formal group $\mathcal{G}_f$, we have its Lie algebra $\mathfrak{g}\simeq R$. Moreover, for $g(x)=b_0x+b_1x^2+\cdots$, If $g$ is an endomorphism of $f$, then $g$ induces a map on the Lie algebra given by multiplication with $b_0$. 
	\end{lemma}
	\begin{proof}
		Direct. 
	\end{proof}

	\begin{thm}
		$\mathcal{M}_{FG}^s$ represents the functor sending $R$ to all coordinatizable formal groups, with morphisms isomorphisms between the coordinatizable formal groups preserving the trivialization of the Lie algebra. 
	\end{thm}
	\begin{proof}
		By definition. 
	\end{proof}
	
	\begin{prop}\label{FGL_coord_condition}
		A formal group $\mathcal{G}$ is coordinatizable iff $\mathfrak{g}\simeq R$. 
	\end{prop}
	\begin{proof}
		$\implies$: Obvious. \\
		$\impliedby$: We only present an idea here. We fix an isomorphism $\mathfrak{g}\simeq R$. Since $\mathcal{G}$ is a formal group, we localize it by the second property, and we have $\mathcal{G}$ is coordinatizable on each open set $D(t_i)$, say isomorphic to $\mathcal{G}_{f_i}$. Notice that this identification makes $\mathfrak{g}$ a locally free module of rank 1 on $R$. Then, we attempt to glue all $\mathcal{G}_{f_i}$ together. Suppose that the transition between the restriction of $\mathcal{G}_{f_i}$ and $\mathcal{G}_{f_j}$ on $D(t_it_j)$ is given by $g_{i,j}(t)=b_0t+b_1t^2+\cdots$. Then the transition between the restriction of $\mathfrak{g}_i$ and $\mathfrak{g}_j$ on $D(t_it_j)$ is exactly given by $b_0$, by \cref{Coord_FGL_endo}. Then, $\mathfrak{g}_i$ glueing to $\mathfrak{g}$ is isomorphic to $R$ iff all $b_0$ term agrees, which is also equivalent to the condition that all $\mathcal{G}_{f_i}$ glues together to some $\mathcal{G}_f$. 
		\todo{Make this precise!}
	\end{proof}

	Now we have seen that ordinary formal group laws equipped with strict isomorphisms corresponds to coordinatizable formal groups, hence to $\mathcal{M}_{FG}^s$. We now present a concept so as to make non-coordinatizable formal groups more concrete. 
	\begin{defn}\label{Defn_twisted_FGL}
		Let $R$ be a ring and $\mathcal{L}$ be an invertible $R$-module. An $\mathcal{L}$-twisted formal group law is a formal series $f(x,y)=\sum a_{i,j}x^iy^j$, where $a_{i,j}\in\mathcal{L}^{\otimes i+j-1}$, and satisfies the usual identities of formal group laws. 
	\end{defn}
	\begin{remark}
		Here the degree $a_{i,j}\in\mathcal{L}^{\otimes i+j-1}$ is since this term behaves like an $(i+j-1)$-fold tensor. Consider the isomorphism $g(t)=\lambda t$, then $a_{i,j}$ is multiplied by $\lambda^{i+j-1}$ times. 
	\end{remark}

	\begin{thm}
		$\mathcal{L}$-twisted formal group laws corresponds to formal groups with their Lie algebra equal to $\mathcal{L}^{-1}$. 
	\end{thm}
	\begin{proof}
		Similar. 
	\end{proof}

	\begin{remark}
		Notice the notion of $\mathcal{L}$-twisted formal group laws is merely the formal group laws respecting gradings on the ring $\bigoplus_{n\in\mathbb{Z}}\mathcal{L}^{\otimes n}$, where $\mathcal{L}^{\otimes n}$ has grade $2n$. Therefore $\bigoplus\mathcal{L}^{\otimes n}$ is equipped with a natural action by the lazard ring $L$ given a $\mathcal{L}$-twisted formal group law $f$. 
	\end{remark}

	\begin{const}
		We may consider the assignment $(R,\mathcal{G})\mapsto\mathfrak{g}^{-1}$ as a line bundle $\omega$ on the moduli stack $\mathcal{M}_{FG}$. With the above \cref{FGL_coord_condition}, we see that $\mathcal{M}_{FG}^s$ is equivalently, the total space of $\omega$ with the zero section removed (in other words, the space of sections/trivializations of $\omega$). This gives a more precise version of \cref{Sec_page_Adams_Nov}. 
	\end{const}
	\begin{thm}\label{Sec_page_Adams_Nov_2}
		Let $X$ be a spectrum. Then notice $MU_{2*}(X)$ and $MU_{2*+1}(X)$ are both modules over $L=\pi_*(MU)$ carrying compatible actions of $G^+$, and therefore determine sheaves $\mathcal{F}_X^{\mathrm{even}}$ and $\mathcal{F}_X^{\mathrm{odd}}$. Then the second page of the Adams-Novikov spectral sequence is given by 
		\[ E_2^{2a,b}=H^b(\mathcal{M}_{FG};\mathcal{F}_X^{\mathrm{even}}\otimes\omega^a) \]
		and 
		\[ E_2^{2a+1,b}=H^b(\mathcal{M}_{FG};\mathcal{F}_X^{\mathrm{odd}}\otimes\omega^a) \]
	\end{thm}

\section{The stratification on $\mathcal{M}_{FG}$}

	Recall from \cref{Lazard_thm} that we have a rational isomorphism $L\to\mathbb{Z}[b_1,\dots]$. This implies that over $\mathbb{Q}$, all formal group laws are uniquely isomorphic to the additive formal group law. This implies the following. 
	\begin{coro}
		$\mathcal{M}_{FG}^s\times\Spec\mathbb{Q}\simeq\Spec\mathbb{Q}$, and $\mathcal{M}_{FG}\times\Spec\mathbb{Q}\simeq B\mathbb{G}_m\times\mathbb{Q}$. Here $B\mathbb{G}_m$ is the classifying stack of a group scheme, as usual. 
	\end{coro}
	\begin{exam}
		Let $f(x,y)=x+y+xy$, the multiplicative formal group law. Then the isomorphism between $f$ and the additive formal group law is given by $t\mapsto e^t-1$, which has rational coefficients. This implies that the two formal group laws are not generally isomorphic. We would want an invariant to tell the two formal group laws apart. To do this, we make the following definition. 
	\end{exam}

	\begin{defn}
		For $f(x,y)$ a formal group law over a ring $R$, we define its $n$-series $[n](t)\in R[[t]]$ as follows. 
		\begin{enumerate}
			\item $[0](t)=0$. 
			\item $[n](t)=f([n-1](t),t)$. 
		\end{enumerate}
	\end{defn}
	\begin{defn}
		We define a homomorphism $h$ between formal group laws $f$ and $f'$ to be a power series in $tR[[t]]$ such that $f(h(x),h(y))=h(f'(x,y))$. 
	\end{defn}
	\begin{prop}
		$[n]$ is a homomorphism from $f$ to itself, that is, $[n]f(x,y)=f([n]x,[n]y)$. 
	\end{prop}
	\begin{proof}
		We prove by induction. The case $n=0$ is obvious. Suppose the proposition is correct for $n-1$, then we have the following equations. 
		\begin{align*}
			&[n]f(x,y) \\
			=&f(f(x,y),[n-1]f(x,y)) \\
			=&f(f(x,y),f([n-1]x,[n-1]y)) \\
			=&f(x,f(y,f([n-1]y,[n-1]x)))~~~~~\text{by~associativity} \\
			=&f(x,f(f(y,[n-1]y),[n-1]x)) \\
			=&f(x,f([n]y,[n-1]x)) \\
			=&f(f(x,[n-1]x),[n]y) \\
			=&f([n]x,[n]y)
		\end{align*}
	\end{proof}

	We wish to develop other properties of $[n]$. Therefore we introduce the following definition. 
	\begin{defn}
		Let $R[[t]]\der t$ denote a free module of rank $1$ over $R[[t]]$. For a formal group law $f(x,y)$, we denote 
		\[ f^*(g(t)\der t)=g(f(x,y))\left(\pd{f}{x}\der x+\pd{f}{y}\der y\right)\in R[[x,y]]\{\der x,\der y\} \]
		We define $g(t)\der t$ to be a translation invariant differential if $f^*(g(t)\der t)=g(x)\der x+g(y)\der y$. 
	\end{defn}

	\begin{exam}
		$\der t$ is a translation invariant differential for $f(x,y)=x+y$. \\
		$\der t/(1+t)$ is a translation invariant differential for $f(x,y)=x+y+xy$. 
	\end{exam}
	\begin{lemma}\label{Trans_inv_diff_gen}
		For any formal group law $f$, there is a unique translation invariant differential $\omega_f$ of the form $(1+c_1t+\cdots)\der t$. Moreover, $R[[t]]\der t\simeq R[[t]]\omega_f$. 
	\end{lemma}
	\begin{proof}
		Direct. 
	\end{proof}
	\begin{lemma}\label{Pullback_trans_inv_diff}
		Let $f$ and $f'$ be formal group laws over $R$. Let $h$ be a homomorphism from $f$ to $f'$. Then $h^*\omega_{f'}=\lambda\omega_f$, where $\lambda\in R$ is a constant. Here 
		\[ h^*(g(t)\der t)=g(h(t))(\pd{h}{t}\der t) \]
	\end{lemma}
	\begin{proof}
		Direct. 
	\end{proof}
	\begin{coro}\label{Homomorphism_FGL_power}
		For $h$ a homomorphism between formal group laws $f$ and $f'$ over a ring with characteristic $p$, one of the following holds. 
		\begin{enumerate}
			\item There is some $\lambda\neq0$ such that $h(t)=\lambda t+\cdots$. 
			\item There is some $n$ and $\lambda\neq0$ such that $h(t)=h'(t^{p^n})$, and $h'(t)=\lambda t+\cdots$. 
		\end{enumerate} 
	\end{coro}
	\begin{proof}
		If the $\lambda$ in \cref{Pullback_trans_inv_diff} is nonzero, then this fits in the first case. Otherwise, $h^*$ sends $\omega_{f'}$ to $0$. However, since $\omega_{f'}$ is a generator of $R[[t]]\der t$, we have $h^*=0$. Then unwinding the definitions, we have $\der h=0$. Since $R$ has characteristic $p$, this implies all non-zero terms in $h$ can only have the degree divisable by $p$. Therefore $h(t)=h_1(t^p)$. Then we notice that $h_1$ is also a homomorphism between formal group laws, namely, $f_1(x,y)=f^p(x,y)$ and $f'_1(x,y)=f'(x^p,y^p)$. Explicitly, we have 
		\[ h_1f^p(x,y)\!=\!hf(x,y)\!=\!f'(h(x),h(y))\!=\!f'(h_1(x^p),h_1(y^p))\!=\!f'(h_1^p(x),h_1^p(y))\!=\!f'_1(h_1(x),h_1(y)) \]
		Therefore, we have either $h_1$ having non-zero $\lambda$ or $h_1(t)=h_2(t^p)$. This process ends somewhere, say at $h_n$. Then $h_n$ is the $h'$ required. 
	\end{proof}

	\begin{coro}
		Either $[p]=0$ or $[p]=\lambda t^{p^n}+\cdots$. 
	\end{coro}
	\begin{proof}
		$[n]$ is a homomorphism. 
	\end{proof}
	\begin{defn}
		For a formal group law $f$ and a fixed prime $p$, we let $v_n$ denote the coefficient of $t^{p^n}$ in the $p$-series $[p]$. We say $f$ is of height $\geq n$ if $v_i=0$ for $i<n$, and height exactly $n$ if $f$ has height $\geq n$ and $v_n$ is invertible. 
	\end{defn}
	\begin{exam}
		$v_0=p$. Thus $f$ has height $\geq 0$ iff $p=0$ in $R$, and height exactly $0$ iff $p$ is invertible in $R$. \\
		The additive formal group law has height infinity if $p=0$ in $R$, and the multiplicative formal group law has height $1$. Therefore, the two formal group laws are not isomorphic in rings with some $p=0$. 
	\end{exam}

	From now on, to focus on the prime $p$, we always assume the ring $R$ we are considering is a $\mathbb{Z}_{(p)}$-algebra, so that any other prime is invertible. 
	We then consider the universal $v_n$'s associated to the universal formal group law on the Lazard ring. Recall that in the proof of \cref{Lazard_thm}, the $t_i$'s are not canonically determined, but chosen to be generators of the canonically deterimined $(I/I^2)_{2n}$. We consider the image of $v_i$ in it. 
	\begin{prop}
		The image of $v_n\in (I/I^2)_{2(p^n-1)}\simeq\mathbb{Z}$ is $p^{p^n-1}-1$. 
	\end{prop}
	\begin{proof}
		Let $k=p^n-1$. The homomorphism $L\to\mathbb{Z}\oplus\mathbb{Z}t_k$ we constructed in the proof of \cref{Lazard_thm} classifies the formal group law 
		\[ f(x,y)=x+y+\sum_{0<i<p^n}\frac{1}{p}\binom{p^n}{i}t_kx^iy^{p^n-i}=x+y+\frac{t_k}{p}((x+y)^{p^n}-x^{p^n}-y^{p^n}) \]
		Then notice by induction that 
		\[ [a]=at+\frac{t_k}{p}((at)^{p^n}-at^{p^n}) \]
		Thus the coefficient of $t^{p^n}$ in $[p]t$ is 
		\[ \frac{t_k}{p}(p^{p^n}-p)=(p^{p^n-1}-1)t_k \]
	\end{proof}
	\begin{coro}
		In the identification $L\otimes\mathbb{Z}_{(p)}\simeq\mathbb{Z}_{(p)}[t_1,t_2,\dots]$, we may let each $t_{p^n-1}$ be given by $v_n$. Equivalently, $v_n$ are the generators of the first degree Adams filtration of $\pi_*(MU)$. 
	\end{coro}
	\begin{coro}
		For a field $k$ of characteristic $p$, for all $1\leq n\leq\infty$, there is a formal group law $f$ having height exactly $n$. 
	\end{coro}
	\begin{proof}
		The case $n=\infty$ is given by the additive formal group law. Otherwise, we let $f$ be the formal group law classified by $\mathbb{Z}_{(p)}[t_1,\dots]$ with $t_i\mapsto0$ for $i<p^n-1$, and $v_n\mapsto 1$. 
	\end{proof}

	\begin{lemma}
		The condition that a formal group law has height $\geq n$ or exactly $n$ is Zariski local. 
	\end{lemma}
	\begin{proof}
		Consider an element $x\in R$. Notice that $x$ is zero or invertible iff its image in $R[a_i^{-1}]$ is zero or invertible for all $i$ for some $\sum a_i=1$. 
	\end{proof}
	This gives rise to the following definition. 
	\begin{defn}
		For a not necessarily coordinatizable formal group $\mathcal{F}:\Alg_R\to\Ab$ over $R$, we define it to have height $\geq n$ or exactly $n$ iff $\mathcal{F}_{\Alg_{R'}}$ has height $\geq n$ or exactly $n$ in the above sense for all $R'$ such that $\mathcal{F}_{\Alg_{R'}}$ is coordinatizable. 
	\end{defn}
	\begin{defn}\label{Defn_strat_MFG}
		We define closed substacks of $\mathcal{M}_{FG}\otimes\mathbb{Z}_{(p)}$ with $\mathcal{M}_{FG}^{\geq n}$ given by \\
		$\Spec(L_{(p)}/(v_0,v_1,\dots,v_{n-1}))/G^+$, and $\mathcal{M}_{FG}^{n}$ given by $\Spec(L_{p}/(v_0,v_1,\dots,v_{n-1})[v_n^{-1}])/G^+$. By their constructions, we see that $\mathcal{M}_{FG}^{\geq n}$ classifies formal groups over $R$ of height $\geq n$ and $\mathcal{M}_{FG}^{n}$ classifies formal groups of height exactly $n$. 
	\end{defn}
	We actually have the following results (the proofs are elementary). 
	\begin{thm}
		For any $R$ in which $p=0$, and $f$ is a formal group law over $R$ having infinite height, then $f$ is isomorphic to the additive formal group law. 
	\end{thm}
	\begin{thm}\label{Iso_FGL_height}
		Let $f$ and $f'$ be formal group laws of height exactly $n$ over $R$, and let $R'$ be defined as the ring parametrizing the isomorphisms between $f$ and $f'$ (explicitly, this is the ring $R[b_0^{\pm1},b_1,\dots]/I$, where $I$ is the ideal generated by the coefficients of $x^iy^j$ of the expression $gf(x,y)-f'(g(x),g(y))$ in which $g=b_0t+b_1t^2+\cdots$). Then $R'$ is isomorphic to the direct limit of a system of injective finite etale maps 
		\[ R=R(0)\inj R(1)\inj R(2)\inj\cdots \]
		In particular, if $R$ is strictly Henselian (for example when $R$ is an algebraically closed field of characteristic $p$), we have a map $R'\to R$, that is, $f$ and $f'$ are always isomorphic formal group laws. 
	\end{thm}
	\begin{proof}
		Omitted. 
	\end{proof}

\section{The Landweber exact functor theorem}

	Recall in previous sections, each complex oriented cohomology theory determines an algebra $\pi_*(E)$ over $L\simeq\pi_*(MU)$. We will answer in this section the converse problem, that is, construct complex oriented cohomology theory using algebras over $L$. However, due to many reasons (we will see this later), we will consider quasi-coherent sheaves over $\mathcal{M}_{FG}$ instead of $\Spec L$. 
	Then, recall the definition of flat quasi-coherent sheaves over an algebraic stack from \cref{QCoh_sheaf_stack}. 
	\begin{lemma}
		A coherent sheaf $M$ over $\mathcal{M}_{FG}$ is flat iff for every $\eta\in\mathcal{M}_{FG}(R)$ representing a coordinatizable formal group over $\Spec R$, we have $M(\eta)$ is flat. 
	\end{lemma}
	\begin{proof}
		Recall the notion of flatness (or faithfully flatness) of a coherent sheaf on a scheme is Zariski local. Since $\eta$ represents a formal group law (not necessarily coordinatizable), we consider its localization at $a_i$, where $\sum a_i=1$, where it is coordinatizable on $D(a_i)$. Hence $M(\eta)$ is flat over $\Spec R$ iff $M(\eta|_{D(a_i)})$ is flat over $\Spec R[a_i^{-1}]$, which proves the lemma. 
	\end{proof}
	\begin{coro}
		Consider the $L$-point $\eta_0$ representing the universal formal group law. Then we have $M$ is flat iff $M(\eta_0)$ is flat over $L$. 
	\end{coro}
	\begin{proof}
		For $\eta\in\mathcal{M}_{FG}(R)$ representing a coordinatizable formal group law, $M(\eta)$ is clearly an image of $M(\eta_0)$. Then by the definition of a quasi-coherent sheaf over an algebraic stack, we have $M(\eta)\simeq M(\eta_0)\otimes_LR$. Therefore $M(\eta)$ is flat over $R$ iff $M(\eta_0)$ is flat over $L$. 
	\end{proof}

	\begin{const}
		For an $R$-point $\eta$ of $\mathcal{M}_{FG}$ corresponding to a morphism $q:\Spec R\to\mathcal{M}_{FG}$, the forgetful functor $q^*:\QCoh(\mathcal{M}_{FG})\to\QCoh(\Spec R)$ given by $M\mapsto M(\eta)$ has a left adjoint $q_*:\QCoh(\Spec R)\to\QCoh(\mathcal{M}_{FG})$, which can be written as $q_*(N)(\eta')=N\otimes_RB$, where $\eta'\in\mathcal{M}_{FG}(R')$ classifies a morphism $q':\Spec R'\to\mathcal{M}_{FG}$ and $B$ is the $R\otimes R'$ algebra determined by the pullback 
		% https://q.uiver.app/#q=WzAsNCxbMCwwLCJcXFNwZWMgQiJdLFsxLDAsIlxcU3BlYyBSIl0sWzAsMSwiXFxTcGVjIFInIl0sWzEsMSwiXFxtYXRoY2Fse019X3tGR30iXSxbMCwxXSxbMSwzXSxbMCwyXSxbMiwzXSxbMCwzLCIiLDEseyJzdHlsZSI6eyJuYW1lIjoiY29ybmVyLWludmVyc2UifX1dXQ==
		\[\begin{tikzcd}
			{\Spec B} & {\Spec R} \\
			{\Spec R'} & {\mathcal{M}_{FG}}
			\arrow[from=1-1, to=1-2]
			\arrow[from=1-1, to=2-1]
			\arrow["\ulcorner"{anchor=center, pos=0.125}, draw=none, from=1-1, to=2-2]
			\arrow[from=1-2, to=2-2]
			\arrow[from=2-1, to=2-2]
		\end{tikzcd}\]
		Notice that $B$ is actually the algebra classifying the universal isomorphism of the two formal groups over $R$ and $R'$. We will use this fact later. 
	\end{const}

	\begin{defn}
		We say an $R$-module $N$ is flat (faithfully flat) over $\mathcal{M}_{FG}$ iff $q_*(N)$ is flat (faithfully flat) over $\mathcal{M}_{FG}$. We say $q$ is flat (faithfully flat) if $R$ (as a module) is flat (faithfully flat) over $\mathcal{M}_{FG}$. 
	\end{defn}

	\begin{prop}\label{Flat_mod_homology_lemma}
		Let $q:\Spec R\to\mathcal{M}_{FG}$ be a map classifying the formal group $\eta\in\mathcal{M}_{FG}(R)$, and let $N$ be an $R$-module flat over $\mathcal{M}_{FG}$. Then the functor $M\mapsto M(\eta)\otimes_RN$ is an exact functor from $\QCoh(\mathcal{M}_{FG})\to\QCoh(\Spec R)\simeq\Mod_R$.  
	\end{prop}
	\begin{proof}
		The problem is local on $\Spec R$ since exactness of quasi-coherent sheaves is a local property. Therefore we may assume $\eta$ classifies a coordinatizable formal group $\mathcal{G}_f$. Let $p:\Spec L\to\mathcal{M}_{FG}$ classifies the universal formal group law. Notice that the universal isomorphism between $f_{univ}$ and $f$ is given by the ring $R[b_0^{\pm1},b_1,\dots]$, with structure morphisms from $R$ the obvious map and from $L$ classifying the formal group law given by the action of $g(t)=b_0t+b_1t^2+\cdots$ on the image of $f$. That is, we have a pullback diagram 
		% https://q.uiver.app/#q=WzAsNCxbMCwwLCJcXFNwZWMgUltiXzBeXFxwbSxiXzEsXFxkb3RzXSJdLFsxLDAsIlxcU3BlYyBMIl0sWzAsMSwiXFxTcGVjIFIiXSxbMSwxLCJcXG1hdGhjYWx7TX1fe0ZHfSJdLFswLDEsInEnIl0sWzEsMywicCJdLFswLDIsInAnIiwyXSxbMiwzLCJxIiwyXSxbMCwzLCIiLDEseyJzdHlsZSI6eyJuYW1lIjoiY29ybmVyLWludmVyc2UifX1dXQ==
		\[\begin{tikzcd}
			{\Spec R[b_0^{\pm1},b_1,\dots]} & {\Spec L} \\
			{\Spec R} & {\mathcal{M}_{FG}}
			\arrow["{q'}", from=1-1, to=1-2]
			\arrow["{p'}"', from=1-1, to=2-1]
			\arrow["\ulcorner"{anchor=center, pos=0.125}, draw=none, from=1-1, to=2-2]
			\arrow["p", from=1-2, to=2-2]
			\arrow["q"', from=2-1, to=2-2]
		\end{tikzcd}\]
		Since the polynomial algebra is a free module over $R$, we have $p'$ a faithfully flat morphism. Hence, to show $M\mapsto q^*(M)\otimes_RN$ is exact, it suffices to show that $M\mapsto (p')^*(q^*(M)\otimes_RN)$ is exact. Explicitly, we have $(p')^*(q^*(M)\otimes_RN)\simeq (qp')^*(M)\otimes_{R[b_0^{\pm1},b_1,\dots]}(p')^*(N)\simeq p^*(M)\otimes_L(p')^*(N)$. However, flatness of $N$ over $\mathcal{M}_{FG}$ guarantees the flatness of $(p')^*N$ over $L$, which proves the theorem. 
	\end{proof}
	\begin{coro}\label{Flat_mod_homology}
		Let $M$ be a graded module over $L$. If $M$ is flat over $\mathcal{M}_{FG}$, then the functor $X\mapsto MU_*(X)\otimes_LM$ is a homology theory representable by some spectrum $E$. Later, we will always assume the grading on $M$ is even (though the even grading is not required for this corollary). 
	\end{coro}
	\begin{proof}
		Directly verify the axioms of homology theories. 
	\end{proof}
	\begin{prop}\label{BP_spectra_const}
		$R=\mathbb{Z}_{(p)}[v_1,v_2,\dots]$ as an $L$-module is flat over $\mathcal{M}_{FG}$. 
	\end{prop}
	\begin{proof}
		We prove the morphism $\Spec R\to\Spec L\to\mathcal{M}_{FG}$ is exact. Again, we consider the diagram of the above form. 
		% https://q.uiver.app/#q=WzAsNCxbMCwwLCJcXFNwZWMgUltiXzBeXFxwbSxiXzEsXFxkb3RzXSJdLFsxLDAsIlxcU3BlYyBMIl0sWzAsMSwiXFxTcGVjIFIiXSxbMSwxLCJcXG1hdGhjYWx7TX1fe0ZHfSJdLFswLDEsInEnIl0sWzEsMywicCJdLFswLDIsInAnIiwyXSxbMiwzLCJxIiwyXSxbMCwzLCIiLDEseyJzdHlsZSI6eyJuYW1lIjoiY29ybmVyLWludmVyc2UifX1dXQ==
		\[\begin{tikzcd}
			{\Spec R[b_0^{\pm1},b_1,\dots]:=\Spec B} & {\Spec L} \\
			{\Spec R} & {\mathcal{M}_{FG}}
			\arrow["{q'}", from=1-1, to=1-2]
			\arrow["{p'}"', from=1-1, to=2-1]
			\arrow["\ulcorner"{anchor=center, pos=0.125}, draw=none, from=1-1, to=2-2]
			\arrow["p", from=1-2, to=2-2]
			\arrow["q"', from=2-1, to=2-2]
		\end{tikzcd}\]
		Now the left map is faithfully flat, hence the bottom map is flat iff the composition $q\circ p'$ is flat. Since $\Spec L\to\mathcal{M}_{FG}$ is flat, it suffices to prove $q'$ is flat. However, the map $q'$ exhibits $R[b_0^{\pm1},b_1,\dots]$ as the quotient of $L_{(p)}[b_0^{\pm1},b_1,\dots]$ with the ideal generated by the image of $t_i$ ($i\neq p^r-1$). Finally recall that for $i\neq p^r-1$, the image of $t_i$ is given by $\lambda b_i+\cdots$, where $\lambda$ is invertible in $\mathbb{Z}_{(p)}$ and the rest terms are decomposable. Thus we may replace $b_i$ by the image of $t_i$. Therefore $B$ is a polynomial over $L_{(p)}[b_0^{\pm1}]$ and hence flat over $L$. 
	\end{proof}
	\begin{coro}\label{ff_BP}
		$\Spec\mathbb{Z}_{(p)}[v_1,\dots]$ is faithfully flat over the localization $\mathcal{M}_{FG}\times\mathbb{Z}_{(p)}$. 
	\end{coro}
	\begin{coro}
		The construction $X\mapsto MU_*(X)\otimes_LR\simeq MU_*(X)_{(p)}/(t_i,i\neq p^r-1)$ is represented by a homology theory $BP_*(X)$. This is called the Brown-Peterson spectrum. 
	\end{coro}

	Then we present an extremely useful theorem called the Landweber exact functor theorem. This gives a criterion on flatness. 
	\begin{defn}\label{Reg_seq}
		Recall $x_i\in R$ is called a regular sequence for a module $M$ if for all $n$, $x_n$ is not a zero divisor on $M/(x_0,x_1,\dots,x_{n-1})M$. In particular, every sequence is a regular sequence for a trivial module. 
	\end{defn}
	\begin{thm}\label{Landweber_exact_func_thm}
		For a module $M$ over $L$, $M$ is flat over $\mathcal{M}_{FG}$ iff for every prime $p$, $v_0,v_1,\dots$ is a regular sequence for $M$. We denote such an $M$ to be Landweber-exact. 
	\end{thm}
	\begin{defn}\label{Landweber_exact_spectra}
		By \cref{Flat_mod_homology}, we obviously have a Landweber-exact evenly graded module $M$ corresponds to a spectrum $E$. We also call the spectra arising from this construction Landweber-exact. 
	\end{defn}
	\begin{coro}
		Let $R=\mathbb{Z}[\beta,\beta^{-1}]$, where $\beta$ has degree $2$. Then  the map $L\to R$ classifying the formal group law $x+y+\beta xy$ exhibits $R$ as a Landweber-exact ring. 
	\end{coro}
	\begin{proof}
		We fix a prime $p$. Then we obviously have $v_0=p$ not a zero divisor. Then notice the $[p]$ series is given by 
		\[ [p]t=\sum_{n=1}^p\binom{p}{n}\beta^{n-1}t^n \]
		hence $v_1$ is invertible on $R/v_0R=R/pR$. Therefore $R$ is Landweber-exact. 
	\end{proof}
	\begin{remark}
		We will later see that this represents the complex $K$-theory. 
	\end{remark}
	\begin{const}
		Let $R$ be a commutative ring and $E$ be an elliptic curve defined over $R$. We can associate to $E$ a formal group $\widehat{E}$ given by $\widehat{E}(A)\subseteq\Hom(\Spec A,E)$ to be all points such that the diagram commutes 
		% https://q.uiver.app/#q=WzAsNCxbMCwwLCJcXFNwZWMgQS9OKEEpIl0sWzEsMCwiXFxTcGVjIEEiXSxbMCwxLCJcXFNwZWMgUiJdLFsxLDEsIkUiXSxbMiwzLCIwIiwyXSxbMSwzXSxbMCwxXSxbMCwyXV0=
		\[\begin{tikzcd}
			{\Spec A/\mathfrak{n}(A)} & {\Spec A} \\
			{\Spec R} & E
			\arrow[from=1-1, to=1-2]
			\arrow[from=1-1, to=2-1]
			\arrow[from=1-2, to=2-2]
			\arrow["0"', from=2-1, to=2-2]
		\end{tikzcd}\]
		where the left map is given by the composition $\Spec A/\mathfrak{n}(A)\to\Spec A\to E\to\Spec R$. Moreover, the multiplication is given by the operation on the elliptic curve. This construction gives a map $\mathcal{M}_{Ell}\to\mathcal{M}_{FG}$. 
		\todo{Finish the construction of elliptic cohomology}
	\end{const}

	We now present a proof of \cref{Landweber_exact_func_thm}. Again, flatness is a local condition, hence it suffices to prove $M_{(p)}$ is flat over $\mathcal{M}_{FG}\otimes\mathbb{Z}_{(p)}$ for all $p$. We fix a prime $p$. Before the proof, we first establish a few conclusions. 

	\begin{lemma}\label{Landweber_exact_func_thm_lem_1}
		Recall the construction $q:\mathbb{Z}_{(p)}[v_1,\dots]\to\mathcal{M}_{FG}$ representing the $BP$ spectrum in \cref{BP_spectra_const}. Then a quasi-coherent sheaf $M$ over $\mathcal{M}_{FG}$ is flat iff $q^*(M)$ is a flat module over $\mathbb{Z}_{(p)}[v_1,\dots]$. 
	\end{lemma}
	\begin{proof}
		$\implies$: Obvious since $\mathbb{Z}_{(p)}[v_1,\dots]$ is flat over $\mathcal{M}_{FG}$. \\
		$\impliedby$: We fix a map $f:\Spec R\to\mathcal{M}_{FG}$. We wish to show $f^*(M)$ is flat. By \cref{ff_BP}, we form the following pullback 
		% https://q.uiver.app/#q=WzAsNCxbMCwwLCJcXFNwZWMgQiJdLFsxLDAsIlxcU3BlYyBMIl0sWzAsMSwiXFxTcGVjIFxcbWF0aGJie1p9X3socCl9W3ZfMSxcXGRvdHNdIl0sWzEsMSwiXFxtYXRoY2Fse019X3tGR31cXHRpbWVzXFxtYXRoYmJ7Wn1feyhwKX0iXSxbMCwxLCJxJyJdLFsxLDMsImYiXSxbMCwyLCJmJyIsMl0sWzIsMywicSIsMl0sWzAsMywiIiwxLHsic3R5bGUiOnsibmFtZSI6ImNvcm5lci1pbnZlcnNlIn19XV0=
		\[\begin{tikzcd}
			{\Spec B} & {\Spec L} \\
			{\Spec \mathbb{Z}_{(p)}[v_1,\dots]} & {\mathcal{M}_{FG}\times\mathbb{Z}_{(p)}}
			\arrow["{q'}", from=1-1, to=1-2]
			\arrow["{f'}"', from=1-1, to=2-1]
			\arrow["\ulcorner"{anchor=center, pos=0.125}, draw=none, from=1-1, to=2-2]
			\arrow["f", from=1-2, to=2-2]
			\arrow["q"', from=2-1, to=2-2]
		\end{tikzcd}\]
		and the top morphism is faithfully flat. Therefore it suffices to prove the $(f\circ q')^*(M)$ is flat over $B$. But this is equivalent to the module $f^*(M)\otimes_RB\simeq q^*(M)\otimes_{\mathbb{Z}_{(p)}[v_1,\dots]}B$ by considering the other composition, which is flat over $B$ iff $q^*(M)$ is flat over $\mathbb{Z}_{(p)}$. 
	\end{proof}
	\begin{lemma}\label{Landweber_exact_func_thm_lem_2}
		Let $R$ be a commutative ring containing a non-zero-devisor $x$, and let $M$ be an $R$-module. Then $M$ is flat over $M$ iff the three following conditions are satisfied. 
		\begin{enumerate}
			\item $x$ is a non-zero-divisor on $M$. 
			\item The quotient $M/xM$ is a flat $R/(x)$-module. 
			\item The module $M[x^{-1}]$ is a flat $R[x^{-1}]$-module. 
		\end{enumerate}
	\end{lemma}
	\begin{proof}
		$\implies$: Obvious. \\
		$\impliedby$: We prove $\Tor_i^R(M,N)\simeq 0$ for all $N$. \\
		We first prove the case where $N$ consists of $x$-power-torsions. We may write $N$ as a filtered colimit $N_i$, where $N_i$ consists of the elements annihilated by $x^i$. Then, consider the exact sequence
		\[ 0\lto K\lto N_i\lto xN_i\lto 0 \]
		Now, $N_i$ is $x^{i-1}$-torsion and $K$ is $x$ torsion, and by the long exact sequence of $\Tor$, it suffices to prove for $N$ with $xN=0$. These $N$'s are $R/(x)$ modules, hence by the universal coefficient theorem we have $\Tor_i^R(M,N)\simeq\Tor_i^{R/(x)}(M/xM,N)=0$, by the second assumption. \\
		Then, we prove for general $N$. We consider $K=\ker(N\to N[x^{-1}])$. Then $K$ is obviously $x$-power-torsion, hence to prove $\Tor_i^R(M,N)=0$ it suffices to prove $\Tor_i^R(M,N/K)=0$. That is, we may assume $N\to N[x^{-1}]$ is injective. Then we may consider the exact sequence 
		\[ 0\lto N\lto N[x^{-1}]\lto C\lto 0\]
		Since $C$ is $x$-power-torsion, it suffices to show $\Tor_i^R(M,N[x^{-1}])=0$. Finally, in this case, $N\simeq N[x^{-1}]$ is a module over $R[x^{-1}]$. Then again by the third assumption and the universal coefficient theorem $\Tor_i^R(M,N)\simeq\Tor_i^{R[x^{-1}]}(M[x^{-1}],N)=0$. 
	\end{proof}
	\begin{lemma}\label{Landweber_exact_func_thm_lem_3}
		Every quasi-coherent sheaf on the stack $\mathcal{M}_{FG}^m$ is flat (recall from \cref{Defn_strat_MFG} the definition of $\mathcal{M}_{FG}^m$). 
	\end{lemma}
	\begin{proof}
		For any map $q:\Spec A\to\mathcal{M}_{FG}^m$ classifying a formal group of height exactly $m$ on $A$, and $X$ a quasi-coherent sheaf on $\mathcal{M}_{FG}$, we prove $q^*$ is flat over $A$. Since again flatness is a local notion, we may suppose the formal group classified by $q$ is coordinatizable. Then, choose a formal group law of height $m$ over $\mathbb{F}_p$, classified by a map $f:\Spec\mathbb{F}_p\to\mathcal{M}_{FG}^m$. We form the pullback 
		% https://q.uiver.app/#q=WzAsNCxbMCwxLCJcXFNwZWMgQSJdLFsxLDEsIlxcbWF0aGNhbHtNfV97Rkd9Il0sWzEsMCwiXFxTcGVjXFxtYXRoYmJ7Rn1fcCJdLFswLDAsIlxcU3BlYyBBJyJdLFszLDJdLFsyLDEsImYiXSxbMywwXSxbMCwxLCJxIiwyXV0=
		\[\begin{tikzcd}
			{\Spec A'} & {\Spec\mathbb{F}_p} \\
			{\Spec A} & {\mathcal{M}_{FG}}
			\arrow[from=1-1, to=1-2]
			\arrow[from=1-1, to=2-1]
			\arrow["f", from=1-2, to=2-2]
			\arrow["q"', from=2-1, to=2-2]
		\end{tikzcd}\]
		In fact, by \cref{Iso_FGL_height}, $A'$ is a direct limit of injective finite etale ring extensions. Therefore $A'$ is faithfully flat over $A$. Therefore it suffices to prove $q^*(X)\otimes_AA'$ is flat over $A'$, and $q^*(X)\otimes_AA'\simeq f^*(X)\otimes_{\mathbb{F}_p}A'$, which is flat over $\mathbb{F}_p$, hence flat over $A'$. 
	\end{proof}

	\begin{proof}
		We prove \cref{Landweber_exact_func_thm}. We first prove $v_i$ being a regular sequence implies flantess. Now, let $M$ be a module over $L_{(p)}$. Since for $p'\neq p$, $v_0=p'$ is already invertible hence $v_i$ is a regular sequence, it suffices to consider the case $p'=p$. We wish to prove for $M$ such that $v_0=p,v_1,\dots$ is a regular sequence, the pushforward of $M$ along $\Spec L_{(p)}\to\mathcal{M}_{FG}\times\Spec\mathbb{Z}_{(p)}$ is flat. We form the pullback square 
		% https://q.uiver.app/#q=WzAsNCxbMCwwLCJcXFNwZWMgQiJdLFsxLDAsIlxcU3BlYyBMX3socCl9Il0sWzAsMSwiXFxTcGVjIFxcbWF0aGJie1p9X3socCl9W3ZfMSxcXGRvdHNdIl0sWzEsMSwiXFxtYXRoY2Fse019X3tGR31cXHRpbWVzXFxTcGVjXFxtYXRoYmJ7Wn1feyhwKX0iXSxbMCwxXSxbMSwzXSxbMCwyXSxbMiwzXSxbMCwzLCIiLDEseyJzdHlsZSI6eyJuYW1lIjoiY29ybmVyLWludmVyc2UifX1dXQ==
		\[\begin{tikzcd}
			{\Spec B} & {\Spec L_{(p)}} \\
			{\Spec \mathbb{Z}_{(p)}[v_1,\dots]} & {\mathcal{M}_{FG}\times\Spec\mathbb{Z}_{(p)}}
			\arrow[from=1-1, to=1-2]
			\arrow[from=1-1, to=2-1]
			\arrow["\ulcorner"{anchor=center, pos=0.125}, draw=none, from=1-1, to=2-2]
			\arrow[from=1-2, to=2-2]
			\arrow[from=2-1, to=2-2]
		\end{tikzcd}\]
		By \cref{Landweber_exact_func_thm_lem_1}, it suffices to prove $M_B:=M\otimes_{L_{(p)}}B$ is a flat module over the ring $\mathbb{Z}_{(p)}[v_1,\dots]$. We attempt to prove $\Tor_i^{\mathbb{Z}_{(p)}[v_1,\dots]}(M_B,N)$ vanish. \\
		Recall the functor $\Tor$ commutes with filtered colimits, hence it suffices to show the case where $N$ is finitely presented. The finite presentation gives $N\simeq N_0[v_{n+1},v_{n+2},\dots]$, where $N_0$ is a module over $\mathbb{Z}_{(p)}[v_1,\dots,v_n]$. By the universal coefficient theorem we have $\Tor_i^{\mathbb{Z}_{(p)}[v_1,\dots]}(M_B,N)\simeq\Tor_i^{\mathbb{Z}_{(p)}[v_1,\dots,v_n]}(M_B,N_0)$. Therefore it suffices to prove $M_B$ is flat over all $\mathbb{Z}_{(p)}$. \\
		Notice that the the homomorphisms $\mathbb{Z}_{(p)}[v_1,\dots]\to B$ and $L_{(p)}\to B$ induces different images of $v_i$ in $B$, say $v'_i$ and $v''_i$. However we still have the isomorphism $(v'_1,\dots,v'_m)\simeq(v''_1,\dots,v''_m)$. We denote this ideal by $I_{m+1}$, then prove inductively $M_B/I_mM_B$ is flat over \\
		$R:=\mathbb{Z}_{(p)}[v_1,\dots,v_n]/(p,v_1,\dots,v_{m-1})$ for $m\leq n+1$. Notice when $m=n+1$, the ring is isomorphic to $\mathbb{F}_p$, hence is obvious. Then notice $M_B/I_mM_B\simeq B\otimes_{L_{(p)}}(M/(v_0,\dots,v_{m-1}))$. Since $v_m$ is not a zero divisor on $M/(v_0,\dots,v_{m-1})$ by assumption, and since $B$ is flat over $L_{(p)}$, we have an exact sequence 
		\[ 0\lto M_B/I_mM_B\xrightarrow{v''_m} M_B/I_mM_B\lto M_B/I_{m+1}M_B \]
		Since $v''_m\equiv \lambda v'_m(\mod I_m)$ where $\lambda$ is invertible, $v_m\in R$ is a non-zero-divisor on $M_B/I_mM_B$. Moreover, the quotient $M_B/(I_m,v_m)M_B\simeq M_B/I_{m+1}M_B$ is flat over $R/(v_m)$ by the induction hypothesis. By \cref{Landweber_exact_func_thm_lem_2}, it suffices to prove $(M_B/I_mM_B)[v_m^{-1}]$ is flat over $R[v_m]^{-1}$, which follows from \cref{Landweber_exact_func_thm_lem_3}. The converse follows exactly from reversing the proof.  
	\end{proof}

\section{Morava stablizer groups}
	
	\begin{const}
		Recall from \cref{Iso_FGL_height} we have a unique formal group law over $\overline{\mathbb{F}}_p$ of height $n$, say $f$. We consider the group $G=\Aut(\Spf\overline{\mathbb{F}}_p[[t]])$, where the automorphism is required to respect the formal group structure and not required to be compatible with the structure map $\Spf\overline{\mathbb{F}}_p[[t]]\to\Spec\overline{\mathbb{F}}_p$. Obviously we have an exact sequence 
		\[ 0\lto\Aut_{\overline{\mathbb{F}}_p}(\Spf\overline{\mathbb{F}}_p[[t]])\lto\Aut(\Spf\overline{\mathbb{F}}_p[[t]])\lto\Gal(\overline{\mathbb{F}}_p/\mathbb{F}_p)\lto 0 \]
		where $\Aut_{\overline{\mathbb{F}}_p}(\Spf\overline{\mathbb{F}}_p[[t]])$ is the subgroup of $\Aut(\Spf\overline{\mathbb{F}}_p[[t]])$ respecting the structure map. 
	\end{const}
	\begin{thm}
		We have $\mathcal{M}_{FG}^n\simeq\Spec\overline{\mathbb{F}}_p/G$, where $G$ acts on $\overline{\mathbb{F}}_p$ by the projection $G\to\Gal(\overline{\mathbb{F}}_p/\mathbb{F}_p)$. 
	\end{thm}

	Before considering $G$, we first consider its subgroup $\Aut_{\overline{\mathbb{F}}_p}(\Spf\overline{\mathbb{F}}_p[[t]])$. We recall that $\Aut_{\overline{\mathbb{F}}_p}(\Spf\overline{\mathbb{F}}_p[[t]])$ is the group of units of the ring $\End_{\overline{\mathbb{F}}_p}(\Spf\overline{\mathbb{F}}_p[[t]])$, whose elements are power series $g(t)\in\overline{\mathbb{F}}_p[[t]]$ satisfying $g(f(x,y))=f(g(x),g(y))$. 
	
	\begin{const}
		Let $f^p$ denote the formal group law over $\overline{\mathbb{F}}_p$ given by applying the Frobenius map to the coefficients of $f$. Then $f^p$ is still a formal group law of height $n$, hence non-canonically isomorphic to $f$ by \cref{Iso_FGL_height}, say $\nu$. Since $f(x,y)^p\simeq f^p(x^p,y^p)$, we have 
		\[ \nu(f(x,y)^p)\simeq\nu(f^p(x^p,y^p))\simeq f(\nu(x^p),\nu(y^p)) \]
		Thus $\pi(t)=\nu(t^p)$ is an endomorphism of $f$. \\
		Then, let $g\in\End_{\overline{\mathbb{F}}_p}(\Spf\overline{\mathbb{F}}_p[[t]])$ be arbitrary, and let $g(t)=b_0t+b_1t^2+\cdots$. If $b_0\neq0$, then $g\in\Aut_{\overline{\mathbb{F}}_p}(\Spf\overline{\mathbb{F}}_p[[t]])$. Otherwise $g(t)=g_0(t^p)$ by \cref{Homomorphism_FGL_power}, and $g_0$ is an endomorphism of $f^p$. Thus $g_0\circ\nu^{-1}$ is an endomorphism of $f$, and we have $g=g_0\circ\nu^{-1}\circ\nu\circ\Frob_p=(g_0\circ\nu^{-1})\circ\pi$. Thus, every non-invertible element $g$ of the ring $\End_{\overline{\mathbb{F}}_p}(\Spf\overline{\mathbb{F}}_p[[t]])$ can be written uniquely of the form $g'\pi$. \\
		Therefore, the ring $\End_{\overline{\mathbb{F}}_p}(\Spf\overline{\mathbb{F}}_p[[t]])$ is a (non-commutative) local ring with the non-invertible elements a two-sided ideal. This ideal is the left ideal generated by $\pi$. Moreover, by \cref{Homomorphism_FGL_power}, we see that every endomorphism $g$ is of the form $u\pi^k$ for some $k$. We will refer to $k$ as the valuation of $g$, and denoted by $v(g)$. By convention we set $v(0)=\infty$. Notice we then have $v([p])=n$, where $n$ is the height of $f$. Also, although the ring is non-commutative, we still have $v(gg')=v(g)+v(g')$. 
	\end{const}

	\begin{const}
		We let $\lambda:\End_{\overline{\mathbb{F}}_p}(\Spf\overline{\mathbb{F}}_p[[t]])\to\overline{\mathbb{F}}_p$ sending $g(t)=b_0t+b_1t^2+\cdots$ to $b_0$. Then, consider $[p](t)=\mu t^{p^n}+\cdots$. Since $g$ is a homomorphism of $f$, we have $[p](g(t))=g([p](t))$, which implies $b_0\mu t^{p^n}+\cdots=\mu b_0^{p^n}t^{p^n}$, hence $b_0\in\mathbb{F}_{p^n}$. Thus $\lambda:\End_{\overline{\mathbb{F}}_p}(\Spf\overline{\mathbb{F}}_p[[t]])\to\mathbb{F}_{p^n}$ is surjective.  
	\end{const}

	\begin{const}
		Obviously we have $\End_{\overline{\mathbb{F}}_p}(\Spf\overline{\mathbb{F}}_p[[t]])\simeq\plim(\End_{\overline{\mathbb{F}}_p}(\Spf\overline{\mathbb{F}}_p[[t]])/(\pi^k))$. Then, notice that $\End_{\overline{\mathbb{F}}_p}(\Spf\overline{\mathbb{F}}_p[[t]])/(\pi^k)$ is finite of cardinality $p^{nk}$ (notice that this is true for $k=1$, and the rest follows from the factorization $g=u\pi^k$). Then, notice $[p^k]\in\pi^k$. This induces a well-defined map $\mathbb{Z}/(p^k)\to\End_{\overline{\mathbb{F}}_p}(\Spf\overline{\mathbb{F}}_p[[t]])/(\pi^k)$, which is compatible with the quotient maps. This induces a map $\mathbb{Z}_p\to\End_{\overline{\mathbb{F}}_p}(\Spf\overline{\mathbb{F}}_p[[t]])$, which exhibits a $\mathbb{Z}_p$ algebra structure. Then, we consider the non-commutative ring $D=\End_{\overline{\mathbb{F}}_p}(\Spf\overline{\mathbb{F}}_p[[t]])[p^{-1}]$. Notice that $[p]$ is a non-zero-divisor in $\End_{\overline{\mathbb{F}}_p}(\Spf\overline{\mathbb{F}}_p[[t]])$, hence it can be viewed as a subring of $D$. Moreover, since $u\pi^n=p$, $\pi$ is invertible in $D$. The valuation $v$ extends to $D$ by $v(\lambda/p^k)=v(\lambda)-nk$. 
	\end{const}

	\begin{lemma}
		$\End_{\overline{\mathbb{F}}_p}(\Spf\overline{\mathbb{F}}_p[[t]])\simeq\{x\in D|v(x)\geq0\}$. 
	\end{lemma}
	\begin{proof}
		This is obvious. 
	\end{proof}

	\begin{lemma}
		$\End_{\overline{\mathbb{F}}_p}(\Spf\overline{\mathbb{F}}_p[[t]])$ is $n^2$ dimensional as a vector space over $\mathbb{Q}_p$. 
	\end{lemma}
	\begin{proof}
		Let $\{x_i\}$ denote the set of $\mathbb{F}_p$ basis of $\mathbb{F}_{p^n}$. Then let $g_i$ denote an element in \\
		$\End_{\overline{\mathbb{F}}_p}(\Spf\overline{\mathbb{F}}_p[[t]])$ such that $\lambda(g_i)=x_i$. Then clearly $\{\pi^jg_i\}$ is a basis. 
	\end{proof}

	\begin{lemma}
		Let $g\in D$. Then the conjugate action of $g$ on $\End_{\overline{\mathbb{F}}_p}(\Spf\overline{\mathbb{F}}_p[[t]])/(\pi)\simeq\mathbb{F}_{p^n}$ is given by $b\mapsto b^{v(g)}$. 
	\end{lemma}
	\begin{proof}
		WLOG $g\in\End_{\overline{\mathbb{F}}_p}(\Spf\overline{\mathbb{F}}_p[[t]])$. Then suppose $g(t)=\lambda t^{p^{v(g)}}+\cdots$. Fix $b\in\mathbb{F}_{p^n}$, and let $h(t)=bt+\cdots$. Then let $h'(t)=g\circ h\circ g^{-1}(t)$, from $g\circ h=h'\circ g^{-1}$ we have $\lambda b^{p^{v(g)}}=b'\lambda$. This gives the result. 
	\end{proof}

	\begin{lemma}
		$D$ has center $\mathbb{Q}_p$. 
	\end{lemma}
	\begin{proof}
		Let $g$ be in the center of $D$. We may assume $g\in\End_{\overline{\mathbb{F}}_p}(\Spf\overline{\mathbb{F}}_p[[t]])$, and prove $g$ is in $\mathbb{Z}_p$. Since $\mathbb{Z}_p$ is a closed subset, it suffices to show there is an integer $m$ such that $g\equiv m(\mod p^k)$ for all $k$. We prove by inductin. Since $\pi g\pi^{-1}=g$, the reduction of $g$ mod $\pi$ belongs to $\mathbb{F}_p\subseteq\mathbb{F}_{p^n}$. Therefore we may subtract this integer, and assume $v(g)>0$. Then by the preceding lemma, the conjugate action of $g$ on $\pi$ is trivial, hence $v(g)$ must be divisible by $n$. Thus $g=g'p$ for some $g'$ in the center of $\End_{\overline{\mathbb{F}}_p}(\Spf\overline{\mathbb{F}}_p[[t]])$, and thus $g'$ is congruent to an integer mod $p^{k-1}$, which is by the induction hypothesis. 
	\end{proof}

	\begin{const}
		By the preceding lemma, $D$ is a central simple algebra over $\mathbb{Q}_p$, hence can be indentified with an element of $\mathrm{Br}(\mathbb{Q}_p)$. We construct an isomorphism $\mathrm{Br}(\mathbb{Q}_p)\to\mathbb{Q}/\mathbb{Z}$ as follows. For a Brauer class, we represent it with a division algebra $D'$. Then it contains a ring of integers $\mathcal{O}$ and its maximal ideal $\mathfrak{m}$. We have a valuation $v:D'-\{0\}\to\mathbb{Z}$ with $\mathcal{O}=v^{-1}(\mathbb{Z}_{\geq0})$ and $\mathfrak{m}=v^{-1}(\mathbb{Z}_{>0})$. Finally, the conjugation action gives a surjective homomorphism $D'-\{0\}\to\Gal((\mathcal{O}/\mathfrak{m})/\mathbb{F}_p)$. In particular, the frobenius automorphism on $\mathcal{O}/\mathfrak{m}$ is given by conjugation by some $x\in D'$. We define $\mu(D')=v(x)/v(p)$. Hence, $D$ is the unique central division algebra over $\mathbb{Q}_p$ with $\mu(D)=1/n$. 
	\end{const}

	\begin{const}
		Recall that we have $\End_{\overline{\mathbb{F}}_p}(\Spf\overline{\mathbb{F}}_p[[t]])^{\times}=\Aut_{\overline{\mathbb{F}}_p}(\Spf\overline{\mathbb{F}}_p[[t]])$. We attempt to extend this isomorphism to a map $\chi:D^\times\to\Aut(\Spf\overline{\mathbb{F}}_p[[t]])$. In fact, $\chi$ is given on $\End_{\overline{\mathbb{F}}_p}(\Spf\overline{\mathbb{F}}_p[[t]])-\{0\}$ by sending $g$ to the map $g_0\circ\Frob_p^{v(g)}$ where $g_0$ is the isomorphism of $f$ with $f^{p^{v(g)}}$ determined by $g$ and $\nu$. That is, we have a commutative diagram of short exact sequences. 
		% https://q.uiver.app/#q=WzAsMTAsWzAsMCwiMCJdLFswLDEsIjAiXSxbNCwwLCIwIl0sWzQsMSwiMCJdLFsxLDAsIlxcRW5kX3tcXG92ZXJsaW5le1xcbWF0aGJie0Z9fV9wfShcXFNwZlxcb3ZlcmxpbmV7XFxtYXRoYmJ7Rn19X3BbW3RdXSlee1xcdGltZXN9Il0sWzIsMCwiRF5cXHRpbWVzIl0sWzMsMCwiXFxtYXRoYmJ7Wn0iXSxbMSwxLCJcXEF1dF97XFxvdmVybGluZXtcXG1hdGhiYntGfX1fcH0oXFxTcGZcXG92ZXJsaW5le1xcbWF0aGJie0Z9fV9wW1t0XV0pXntcXHRpbWVzfSJdLFsyLDEsIlxcQXV0KFxcU3BmXFxvdmVybGluZXtcXG1hdGhiYntGfX1fcFtbdF1dKV57XFx0aW1lc30iXSxbMywxLCJcXEdhbChcXG92ZXJsaW5le1xcbWF0aGJie0Z9fV9wLFxcbWF0aGJie0Z9X3ApIl0sWzAsNF0sWzQsNV0sWzUsNiwidiJdLFs2LDJdLFsxLDddLFs3LDhdLFs4LDldLFs5LDNdLFs0LDddLFs1LDhdLFs2LDldXQ==
		\[\begin{tikzcd}
			0 & {\End_{\overline{\mathbb{F}}_p}(\Spf\overline{\mathbb{F}}_p[[t]])^{\times}} & {D^\times} & {\mathbb{Z}} & 0 \\
			0 & {\Aut_{\overline{\mathbb{F}}_p}(\Spf\overline{\mathbb{F}}_p[[t]])^{\times}} & {\Aut(\Spf\overline{\mathbb{F}}_p[[t]])^{\times}} & {\Gal(\overline{\mathbb{F}}_p,\mathbb{F}_p)} & 0
			\arrow[from=1-1, to=1-2]
			\arrow[from=1-2, to=1-3]
			\arrow[from=1-2, to=2-2]
			\arrow["v", from=1-3, to=1-4]
			\arrow[from=1-3, to=2-3]
			\arrow[from=1-4, to=1-5]
			\arrow[from=1-4, to=2-4]
			\arrow[from=2-1, to=2-2]
			\arrow[from=2-2, to=2-3]
			\arrow[from=2-3, to=2-4]
			\arrow[from=2-4, to=2-5]
		\end{tikzcd}\]
		The left map is an isomorphism by definition, and the right morphism is the profinite completion $\mathbb{Z}\to\widehat{\mathbb{Z}}$. Therefore, the middle map is somewhat of a completion. 
	\end{const}

	\begin{prop}
		The formal group law of height $n$ over $\overline{\mathbb{F}}_p$ is in the image of the pushforward map along the inclusion $\mathbb{F}_{p^n}\inj\overline{\mathbb{F}}_p$. 
	\end{prop}
	\begin{const}
		Recall $\Gal(\overline{\mathbb{F}}_p/\mathbb{F}_{p^k})$ can be identified with $k\widehat{Z}$ with respect to the inclusion $\Gal(\overline{\mathbb{F}}_p/\mathbb{F}_{p^k})\inj\Gal(\overline{\mathbb{F}}_p/\mathbb{F}_p)\simeq\widehat{\mathbb{Z}}$. \\
		By descent theory, descending the formal group is equivalent to giving a $\Gal(\overline{\mathbb{F}}_p/\mathbb{F}_{p^k})\simeq k\widehat{\mathbb{Z}}$ action on $\Spf\overline{\mathbb{F}}_p[[t]]$, which is equivariant w.r.t the projection $\Spf\overline{\mathbb{F}}_p[[t]]\to\Spec\overline{\mathbb{F}}_p$. Also, $k\widehat{\mathbb{Z}}$ is topologically cyclic, hence this condition is equivalent to choosing a single element of $\Aut(\Spf\overline{\mathbb{F}}_p[[t]])$ lying over the integer $k$, that is, giving an element of $D^\times$ with $v(x)=k$. Such element exists for all $k\geq 1$, and in particular, when $k=n$, we have a canonical choice $[p]$. Unwinding the definitions, we have proven the proposition. 
	\end{const}

	\begin{coro}
		$\mathcal{M}_{FG}^n$ is equivalently $\mathbb{F}_{p^n}/G'$, where $G'$ fits in an exact sequence 
		\[ 0\lto\Aut_{\mathbb{F}_{p^n}}(\Spf\mathbb{F}_{p^n}[[t]])\lto G'\lto\Gal(\mathbb{F}_{p^n}/\mathbb{F}_p)\lto 0 \]
		The group $G'$ is sometimes called also the morava stabilizer group. 
	\end{coro}

\section{Lubin-Tate theory}

	After we have finished analysing the structure of $\mathcal{M}_{FG}^n$, we attempt to put these strata together. Therefore, we will have to figure out what $\mathcal{M}_{FG}\times\mathbb{Z}_{(p)}$ looks like near one strata, which is where deformation theory gets involved. In this section, we fix a perfect field of characteristic $p$, say $k$, and fix a formal group law of height $n$ over $k$, say $f$. 
	We first recall that the infinitesimal objects in $\Sch$ (or $\Aff$) are the Artinian local rings. 
	\begin{defn}
		We define an infinitesimal thickening of a field $k$ to be an Artinian local ring $A$ whose residue field is $k$. Intuitively, $\Spec A$ is $\Spec k$ with a tiny, pointless neighborhood. 
	\end{defn}
	\begin{defn}
		Let $A$ be an infinitesimal thickening of $k$. For a formal group law $f$ over $k$, we define a deformation of $f$ over $A$ to be a formal group law $f_A$ over $A$, and for $\varphi:A\to k$ the quotient map, $\varphi_*(f_A)=f$. We define two deformations of $f$, say $f_A$ and $f'_A$, to be isomorphic if there is an invertible power series $g(t)\in A[[t]]$ such that $gf_A(x,y)=f'_A(g(x),g(y))$, and $g\equiv t\mod\mathfrak{m}_A$. We denote the groupoid of deformations of $f$ with isomorphisms by $\Def_A(f)$. 
	\end{defn}

	\begin{lemma}\label{Disc_deformation}
		The groupoid $\Def_A(f)$ is discrete. 
	\end{lemma}
	\begin{proof}
		To prove this, it suffices to prove that the loopspace at each point is trivial, that is, let $f_A$ be a deformation of $f$, then $g$ an automorphism of $f_A$ such that $g\equiv t\mod\mathfrak{m}_A$ is the identity. We prove by induction that $g\equiv t\mod\mathfrak{m}_A^a$. Since $A$ is Artinian, for sufficiently large $a$ $\mathfrak{m}_A^a=0$, and would imply the lemma. The case $a=1$ is by assumption. \\
		Then, let $A'$ be the quotient of $A[b_0^{\pm1},b_1,\dots]$ classifying automorphisms of $f_A$. Then the map $g$ is classified by a map $\psi:A'\to A$, and the identity automorphism is classified by $\psi_0:A'\to A$. Assume that $g\equiv t\mod\mathfrak{m}_A^a$. Then the composite maps $\psi,\psi_0:A'\to A\to A/\mathfrak{m}_A^a$ agree. Thus $d=\psi-\psi_0:A'\to\mathfrak{m}_A^a\to V=\mathfrak{m}_A^a/\mathfrak{m}_A^{a+1}$ is a well-defined $A$-derivation. This map factors as a composition $A'\to A'\otimes_Ak\xrightarrow{d'} V$, where the second map, $d'$, is a $k$-derivation. However notice that we have $A'\otimes_Ak$ the ring classifying the automorphisms of $f$. Finally, by \cref{Iso_FGL_height}, we have $A'\otimes_Ak$ etale over $k$, hence $d'=0$. This implies $g\equiv t\mod\mathfrak{m}_A^{a+1}$. 
	\end{proof}

	\begin{defn}
		For a ring $A$, we define the Witt vectors over $A$ to be infinite sequences $a=(a_0,a_1,\dots)\in A^{\mathbb{N}}$. This set admits a ring structure, with addition given by 
		\[ (a_0,a_1,\dots)+(b_0,b_1,\dots)=(S_0(a_0,b_0),S_1(a_0,b_0,a_1,b_1),\dots) \]
		and multiplication 
		\[ (a_0,a_1,\dots)(b_0,b_1,\dots)=(M_0(a_0,b_0),M_1(a_0,b_0,a_1,b_1),\dots) \]
		where $S_n$ and $M_n$ are polynomials of $x_0,\dots,x_n,y_0,\dots,y_n$ with integer coefficients determined by 
		\[ \Phi_n(S_0,\dots,S_n)=\Phi_n(x_0,\dots,x_n)+\Phi_n(y_0,\dots,y_n) \]
		and 
		\[ \Phi_n(M_0,\dots,M_n)=\Phi_n(x_0,\dots,x_n)\Phi_n(y_0,\dots,y_n) \]
		and 
		\[ \Phi_n(z_0,\dots,z_n)=z_0^{p^n}+pz_1^{p^{n-1}}+\cdots+p^nz_n \]
		This ring is denoted by $W(A)$. Notice the zero element is $(0,0,\dots)$ and the identity is $(1,0,\dots)$. 
	\end{defn}

	\begin{const}
		Let $R=W(k)[[v_1,\dots,v_{n-1}]]$. Then we have a canonical map $R\to k$ with kernel the maximal ideal $\mathfrak{m}_R=(p,v_1,\dots,v_{n-1})$. The formal group law $f$ over $k$ is classified by a map $\varphi_0:L_{(p)}\to k$. Since $f$ has height $n$, we obviously have $t_{p^i-1}$ is send to $0$ by $\varphi_0$ for $i<n$. We let $\varphi:L_{(p)}\to R$ be any map that lifts $\varphi_0$ and sends $t_{p^i-1}$ to $v_i$. Then this map determines a formal group law over $R$, denoted by $\overline{f}$. whose image is $f$ under the canonical map $R\to k$. 
	\end{const}

	\begin{thm}\label{Lubin_Tate}
		The formal group law $\overline{f}$ over $R$ is a universal deformation of $f$. That is, for any infinitesimal thickening of $k$, $\overline{f}$ induces a bijection 
		\[ \Hom_{k}(R,A)\xrightarrow{\simeq}\Def_A(f) \]
	\end{thm}

	\begin{lemma}\label{Lubin_Tate_lem_1}
		If $A\to A'$ is a surjective map between infinitesimal thickenings of $k$, then the induced map $\Def_A(f)\to\Def_{A'}(f)$ is surjective. 
	\end{lemma}
	\begin{proof}
		This is since the Lazard ring is a polynomial ring, and we may simply lift image of the generators. 
	\end{proof}

	\begin{lemma}\label{Lubin_Tate_lem_2}
		Given a pair of surjective maps $A\to B\from C$ between infinitesimal thickenings of $k$, the canonical map $\Def_{A\times_BC}(f)\to\Def_A(f)\times_{\Def_B(f)}\Def_F(f)$ is a bijection. 
	\end{lemma}
	\begin{proof}
		Notice that since Artinian local rings has only trivial Zariski open covers, a formal group over $A$ is equivalent to a formal group law over $A$. Then this conclusion is equivalent to the fact that formal groups satisfies descent. 
	\end{proof}

	\begin{lemma}\label{Lubin_Tate_lem_3}
		\cref{Lubin_Tate} holds for $A=k[x]/(x^2)$. 
	\end{lemma}
	\begin{proof}
		We first recall that $\Hom_k(W(k),k[x]/(x^2))\simeq 0$ by commutative algebra, hence \\
		$\Hom_k(R,k[x]/(x^2))\simeq k^{n-1}$. Notice an element in $\Def_A(f)$ is equivalently a map $L\to A$ with compatibility conditions. In this case, the compatibility conditions is simply $v_i\mapsto c_ix$. We construct a map $\theta':\Def(k[x]/(x^2))\to k^{n-1}$ given by $f:L\to k[x]/x^2$ mapping to $(c_1,\dots,c_{i-1})$. Thus the composition 
		\[ \Hom_k(R,k[x]/(x^2))\to\Def_A(f)\to k^{n-1} \]
		is an isomorphism by construction. Then it suffices to prove the second map is an isomorphism. We have seen it is surjective, and we prove injectivity. But this is by \cref{Disc_deformation}.  
	\end{proof}

	\begin{proof}
		We now prove \cref{Lubin_Tate}. We proceed by induction on the length of the Artinian ring $A$. If $A$ has length $1$, then $A=k$ and the conclusion obviously holds. If $A$ has length $>1$, we pick $x$ annihilated by $\mathfrak{m}_A$. By \cref{Lubin_Tate_lem_2}, we have the pullback diagram 
		% https://q.uiver.app/#q=WzAsNCxbMSwxLCJcXERlZl97QS8oeCl9KGYpIl0sWzEsMCwiXFxEZWZfQShmKSJdLFswLDEsIlxcRGVmX0EoZikiXSxbMCwwLCJcXERlZl97QVxcdGltZXNfe0EvKHgpfUF9KGYpIl0sWzMsMV0sWzEsMF0sWzMsMl0sWzIsMF0sWzMsMCwiIiwxLHsic3R5bGUiOnsibmFtZSI6ImNvcm5lci1pbnZlcnNlIn19XV0=
		\[\begin{tikzcd}
			{\Def_{A\times_{A/(x)}A}(f)} & {\Def_A(f)} \\
			{\Def_A(f)} & {\Def_{A/(x)}(f)}
			\arrow[from=1-1, to=1-2]
			\arrow[from=1-1, to=2-1]
			\arrow["\ulcorner"{anchor=center, pos=0.125}, draw=none, from=1-1, to=2-2]
			\arrow[from=1-2, to=2-2]
			\arrow[from=2-1, to=2-2]
		\end{tikzcd}\]
		Notice $A\times_{A/(x)}A\simeq k[x]/(x^2)\times_kA$. The addition map $k[x]/(x^2)\times_kk[x]/(x^2)\to k[x]/(x^2)$ determines a group structure on $\Def_{k[x]/(x^2)}(f)$ by \cref{Lubin_Tate_lem_2}. And the multiplication $k[x]/(x^2)\times_kA\to A$ determines an action of $\Def_{k[x]/(x^2)}(f)$ on $\Def(A)$, and the pullback diagram shows that we have an embedding $\Def_A(f)/\Def_{k[x]/(x^2)}(f)\inj\Def_{A/(x)}(f)$. By \cref{Lubin_Tate_lem_1} this map is also surjective. Similarly, we may show that $\Hom_k(R,A)/\Hom_k(R,k[x]/(x^2))\simeq\Hom_k(R,A/(x))$. Now by \cref{Lubin_Tate_lem_3} and the induction hypothesis for $\Hom_k(R,A/(x))$ we obtain the result. 
	\end{proof}

	\begin{coro}
		Let $A$ be a complete Noetherian local ring with residue field $k$ and maximal ideal $\mathfrak{m}_A$. Then \cref{Lubin_Tate} also holds for $A$. 
	\end{coro}
	\begin{proof}
		Directly write $A\simeq\plim A/\mathfrak{m}_A^a$, where each term is a Artinian local ring. 
	\end{proof}

	\begin{const}\label{Univ_def_FGL}
		By the preceding corollary, $W(k)[[v_1,\dots,v_{n-1}]]$ is Yonedianly determined by the residue field $k$ together with a choice of a formal group law of height $n$ over $k$. Taking $k=\overline{\mathbb{F}}_p$, we have a action of $G$ on $W(k)[[v_1,\dots,v_{n-1}]]$, where $G$ is the Morava stabilizer group. 
	\end{const}

\chapter{Morava $E$ theory and Morava $K$ theory}

\section{Landweber-exact spectra}

	\begin{defn}
		Let $f:E\to E'$ be a map of spectra. We say $f$ is a phantom if the underlying map $f_*:E_*\to E'_*$ between homology theories is the zero map. 
	\end{defn}
	
	\begin{prop}\label{Equiv_cond_phantom}
		Let $f:E\to E'$ be a map of spectra. The following conditions are equivalent. 
		\begin{enumerate}
			\item The map $f$ is a phantom. 
			\item For every spectrum $X$, the map $E_*(X)\to E'_*(X)$ is zero. 
			\item For every finite spectrum $X$, the map $E_*(X)\to E'_*(X)$ is zero. 
			\item For every finite spectrum $X$, the map $E^*(X)\to {E'}^*(X)$ is zero. 
			\item For every finite spectrum $X$ and every map $g:X\to E$, the composition $fg:X\to E'$ is zero. 
		\end{enumerate}
	\end{prop}
	\begin{proof}
		$(1)\iff(2)\iff(3)$ is by the fact that every spectra can be written as a filtered colimit of $\Sigma^nS$. Then $(3)\iff(4)$ by Spanier-Whitehead duality, and $(4)\iff(5)$ is obvious. 
	\end{proof}

	\begin{defn}\label{Defn_even_spectra}
		We say a finite spectrum $X$ is even if the homology groups $H\mathbb{Z}_*(X)$ are free abelian groups concentrated on even degrees. Equivalently, a finite spectrum $X$ is even if it admits a finite cell decomposition using only even dimensional cells. \\
		We say a spectrum $E$ is evenly generated if for every map $X\to E$ where $X$ is finite, we have a factorization of the map $X\to X'\to E$, where $X'$ is even. 
	\end{defn}

	\begin{prop}
		Any Landweber-exact spectrum $E$ is evenly generated. 
	\end{prop}
	\begin{proof}
		Let $E$ be the Landweber-exact spectrum associated to a graded $L$-module $M$. Then let $f:X\to E$ be a map where $X$ is finite. We may associate $f$ to an element in $E^0(X)\simeq E_0(\dual X)=(MU\otimes_LM)_0(X)=(MU\otimes_LM)^0(X)\simeq\sum c_i\otimes m_i$, where $c_i\in MU^{d_i}(X)$ and $m_i\in M_{d_i}$. Thus $f$ factors as a composition
		\[ X\xrightarrow{(c_i)}\bigoplus\Sigma^{d_i}MU\xrightarrow{(m_i)}E \]
		Therefore it suffices to prove $MU$ is evenly generated. Then, since $MU\simeq\ilim MU(n)$, it suffices to prove each $MU(n)$ is evenly generated (recall that the finite spectra are the compact objects). Recall $MU(n)=\Th(\eta^n-\varepsilon^n)$, where $\eta^n$ is the universal bundle, and $\varepsilon^n$ is the trivial bundle of rank $n$ on $BU(n)$. Finally, since $BU(n)\simeq\ilim Gr_n(\mathbb{C}^{n+k})$, by the Thom isomorphism, it suffices to prove the thom spectra of each $Gr_n(\mathbb{C}^{n+k})$ is an even spectrum, which is obvious by the Bruhat decomposition of $Gr_n$ into even dimensional cells. 
	\end{proof}

	\begin{prop}
		Let $E$ be an evenly generated spectrum and $E'$ be any spectrum whose homotopy groups are concentrated in even degrees. Then every phantom map $E\to E'$ is null. 
	\end{prop}
	\begin{proof}
		We let $A$ denote the set of all finite even spectra equipped with a map to $E$, say $u_\alpha:X_\alpha\to E$. Then we consider the fiber 
		\[ K\xrightarrow{\iota}\bigoplus X_\alpha\xrightarrow{u=\bigoplus u_\alpha} E \]
		Then $u$ is classified by a map $u':E\to\Sigma K$. Since $E$ is evenly generated, every map from a finite spectrum $X\to E$ factors through $u$, hence $X\to E\to\Sigma K$ is null. Thus $u'$ is a phantom. Conversly, if $f:E\to E'$ is any phantom map, then $f\circ u$ is null, hence $f$ factors as a composition $E\to\Sigma K\to E'$. Thus it now suffices to prove every map $\Sigma K\to E'$ is null, that is, ${E'}^{-1}(K)=0$. \\
		Since the homotopy groups of $E'$ are concentrated in even degrees, the Atiyah-Hirzebruch spectral sequence shows that $E'^{-1}(X)\simeq0$ for $X$ a finite even spectrum. Thus it suffices to show that $K$ is the retract of a direct sum of finite even spectra. Let $B$ denote the collection of triples $(\alpha,\alpha',f)$, where $\alpha,\alpha'\in A$ and $u$ ranges over all homotopy classes of maps $X_\alpha\to X_{\alpha'}$ commuting with $u_\alpha$ and $u_{\alpha'}$. For $\beta=(\alpha,\alpha',f)$, we let $Y_\beta=X_\alpha$. Then clearly we have a map $\varphi:\bigoplus Y_\beta\to\bigoplus X_\alpha$, given on each factor $Y_\beta$ by the difference map of $Y_\beta\simeq X_\alpha\inj\bigoplus X_\alpha$ and $Y\beta\to X_{\alpha'}\inj\bigoplus X_\alpha$. Let $F$ denote the cofiber of $\varphi$. Since $\bigoplus Y_\beta$ is obviously evenly generated, we have a map of fiber sequences 
		% https://q.uiver.app/#q=WzAsNixbMCwwLCJcXGJpZ29wbHVzIFlfXFxiZXRhIl0sWzEsMCwiXFxiaWdvcGx1cyBYX1xcYWxwaGEiXSxbMiwwLCJGIl0sWzIsMSwiRSJdLFsxLDEsIlxcYmlnb3BsdXMgWF9cXGFscGhhIl0sWzAsMSwiSyJdLFswLDFdLFsxLDJdLFsyLDNdLFsxLDRdLFswLDVdLFs1LDRdLFs0LDNdXQ==
		\[\begin{tikzcd}
			{\bigoplus Y_\beta} & {\bigoplus X_\alpha} & F \\
			K & {\bigoplus X_\alpha} & E
			\arrow[from=1-1, to=1-2]
			\arrow[from=1-1, to=2-1]
			\arrow[from=1-2, to=1-3]
			\arrow[from=1-2, to=2-2]
			\arrow[from=1-3, to=2-3]
			\arrow[from=2-1, to=2-2]
			\arrow[from=2-2, to=2-3]
		\end{tikzcd}\]
		We construct a map $E\to F$ right inverse to the map $F\to E$, so as to exhibit $K$ as the retract of $\bigoplus Y_\beta$. It suffices to construct a map between the homology theories $E_*\to F_*$. By passing to filtered colimits and Spanier-Whitehead duality, it suffices to construct for every finite spectrum $X$ and a map $f:X\to E$ a map $q(f):X\to F$. \\
		Since $E$ is evenly generated, $f$ factors as $X\to X_{\alpha'}\to E$ for some $\alpha'$. Then we define $q(f)$ to be the map $X\to X_{\alpha'}\to\bigoplus X_\alpha\to F$. We prove $f$ is well-defined. For another factorization $X\to X_{\alpha''}\to E$, let $Y$ denote the pushout $X_{\alpha'}\amalg_XX_{\alpha''}$. Then $Y$ is a finite spectrum, hence the map $Y\to E$ induced by our assumptions factors as $Y\to X_\alpha\to E$. Let $h'$ denote the composition $X_{\alpha'}\to Y\to X_\alpha$ and $h''$ similarly. Then $(\alpha',\alpha,h')$ and $(\alpha'',\alpha,h'')$ are elements of $B$. Thus the composite maps $X\to X_{\alpha'}\to\bigoplus X_\alpha\to F$ and $X\to X_{\alpha''}\to\bigoplus X_\alpha\to F$ can be identified with $X\to Y\to X_\alpha\to\bigoplus X_\alpha\to F$, hence $q$ is well-defined. 
	\end{proof}

	\begin{coro}\label{Trivial_phantom}
		Let $E$ be a Landweber-exact spectrum, and let $E'$ be a spectrum such that $\pi_{2k+1}(E')=0$. Then every phantom map $f:E\to E'$ is null-homotopic. 
	\end{coro}

	\begin{lemma}
		Recall the definition of $\mathcal{L}$-twisted formal group laws from \cref{Defn_twisted_FGL}. Then the following conditions are equivalent for a $\mathcal{L}$-twisted formal group law $f$. 
		\begin{enumerate}
			\item The associated formal group $\mathcal{G}_f$ is classified by a flat map $q:\Spec R\to\mathcal{M}_{FG}$. 
			\item The graded $L$-algebra $\bigoplus\mathcal{L}^{\otimes n}$ is Landweber-exact. 
		\end{enumerate}
	\end{lemma}
	\begin{proof}
		$(2)\implies(1)$: The condition is equivalent to $\bigoplus\mathcal{L}^{\otimes n}$ flat over $\mathcal{M}_{FG}$ by \cref{Landweber_exact_func_thm}. Thus we obviously have $\mathcal{L}^{\otimes0}\simeq R$ flat over $\mathcal{M}_{FG}$. \\
		$(1)\implies(2)$: $\bigoplus\mathcal{L}^{\otimes n}$ is flat over $R$. 
	\end{proof}

	\begin{coro}\label{Const_flat_ring_spectra}
		In the situation of the above lemma, we may define a spectrum $E_R$ given by $(E_R)_*(X)=MU_*(X)\otimes_L(\bigoplus\mathcal{L}^{\otimes n})$. 
	\end{coro}
	\begin{proof}
		Directly by \cref{Flat_mod_homology}. 
	\end{proof}

	\begin{const}
		Let $R=L$ and $\mathcal{L}=R$ be trivial. Then $\bigoplus\mathcal{L}^{\otimes n}\simeq L[\beta^{\pm1}]$. Then the above construction gives a spectrum $E_L$ whose homology theory is given by $(E_L)_*(X)=MU_*(X)[\beta^{\pm1}]$. We denote this spectrum by $MP$, the periodic complex cobordism spectra. In particular, $MP_0(X)\simeq\bigoplus MU_{2n}(X)$. \\
		Moreover, the above construction generalizes as follows. We let $\mathcal{L}\simeq R$, then we define $(E_R)_*(X)=MU_*(X)\otimes_LR[\beta^{\pm1}]=MP_*(X)\otimes_LR$. 
	\end{const}

	\begin{const}
		The construction in the corollary can be reformulated as follows. Let $q:\Spec R\to\mathcal{M}_{FG}$ be a flat map. Then $E_R$ is uniquely determined by the definition $(E_R)_0(X)=q^*\mathcal{F}_X$, where $\mathcal{F}_X$ is the associated quasi-coherent sheaf of $X$ over $\mathcal{M}_{FG}$. 
	\end{const}

	\begin{const}\label{Well_def_map_spectra_from_rings}
		For a commutative diagram 
		% https://q.uiver.app/#q=WzAsMyxbMCwwLCJcXFNwZWMgUiJdLFsyLDAsIlxcU3BlYyBSJyJdLFsxLDEsIlxcbWF0aGNhbHtNfV97Rkd9Il0sWzAsMSwiZiJdLFsxLDIsInEnIl0sWzAsMiwicSIsMl1d
		\[\begin{tikzcd}
			{\Spec R'} && {\Spec R} \\
			& {\mathcal{M}_{FG}}
			\arrow["f", from=1-1, to=1-3]
			\arrow["{q'}"', from=1-1, to=2-2]
			\arrow["q", from=1-3, to=2-2]
		\end{tikzcd}\]
		where $q$ and $q'$ are flat, we have an obvious induced map in homology $(E_R)_*(X)\to(E_{R'})_*(X)$. Then by \cref{Trivial_phantom}, we have a uniquely determined map up to homotopy $f:E_R\to E_{R'}$. 
	\end{const}

	\begin{prop}
		Suppose we have flat maps $q:\Spec R\to\mathcal{M}_{FG}$ and $q':\Spec R'\to\mathcal{M}_{FG}$, the smash product $E_R\wedge E_{R'}$ is homotopy equivalent to $E_B$, where $B$ fits in the following Cartesian diagram 
		% https://q.uiver.app/#q=WzAsNCxbMSwwLCJcXFNwZWMgUiciXSxbMCwxLCJcXFNwZWMgUiJdLFsxLDEsIlxcbWF0aGNhbHtNfV97Rkd9Il0sWzAsMCwiXFxTcGVjIEIiXSxbMSwyLCJxIiwyXSxbMCwyLCJxJyJdLFszLDFdLFszLDBdLFszLDIsIiIsMSx7InN0eWxlIjp7Im5hbWUiOiJjb3JuZXItaW52ZXJzZSJ9fV1d
		\[\begin{tikzcd}
			{\Spec B} & {\Spec R'} \\
			{\Spec R} & {\mathcal{M}_{FG}}
			\arrow[from=1-1, to=1-2]
			\arrow[from=1-1, to=2-1]
			\arrow["\ulcorner"{anchor=center, pos=0.125}, draw=none, from=1-1, to=2-2]
			\arrow["{q'}", from=1-2, to=2-2]
			\arrow["q"', from=2-1, to=2-2]
		\end{tikzcd}\]
	\end{prop}
	\begin{proof}
		Clearly $\Spec B$ is flat over $\mathcal{M}_{FG}$. For simplicity, we assume the formal groups over $R$ and $R'$ classified by $q$ and $q'$ are coordinatizable, given by maps $L\to R$ and $L\to R'$. Then, recall from the construction in the proof of \cref{Flat_mod_homology_lemma}, we have the pullback 
		% https://q.uiver.app/#q=WzAsNCxbMSwwLCJcXFNwZWMgTCJdLFswLDEsIlxcU3BlYyBMIl0sWzEsMSwiXFxtYXRoY2Fse019X3tGR30iXSxbMCwwLCJcXFNwZWMgTVBfMChNUCk9XFxTcGVjIExbYl8wXntcXHBtMX0sYl8xLFxcZG90c10iXSxbMSwyXSxbMCwyXSxbMywxXSxbMywwXSxbMywyLCIiLDEseyJzdHlsZSI6eyJuYW1lIjoiY29ybmVyLWludmVyc2UifX1dXQ==
		\[\begin{tikzcd}
			{\Spec MP_0(MP)=\Spec L[b_0^{\pm1},b_1,\dots]} & {\Spec L} \\
			{\Spec L} & {\mathcal{M}_{FG}}
			\arrow[from=1-1, to=1-2]
			\arrow[from=1-1, to=2-1]
			\arrow["\ulcorner"{anchor=center, pos=0.125, rotate=45}, draw=none, from=1-1, to=2-2]
			\arrow[from=1-2, to=2-2]
			\arrow[from=2-1, to=2-2]
		\end{tikzcd}\]
		Thus for general $R$ and $R'$, we have $B\simeq R\otimes_LMP_0(MP)\otimes_LR'$ by pulling back. Finally, let $X$ be any spectrum, we have 
		\begin{align*}
			(E_R\wedge E_{R'})_0(X)\simeq&(E_R)_0(E_{R'}\wedge X) \\
			\simeq&R\otimes_LMP_0(E_{R'}\wedge X) \\
			\simeq&R\otimes_L(E_{R'})_0(MP\wedge X) \\
			\simeq&R\otimes_LMP_0(MP\wedge X)\otimes_LR' \\
			\simeq&R\otimes_L(MP\wedge MP)_0(X)\otimes_LR'
		\end{align*}
		Thus we have $E_R\wedge E_{R'}\simeq E_B$. 
		\todo{Finish the general case after calculating the general pullback.}
	\end{proof}

	\begin{coro}
		For any flat map $\Spec R\to\mathcal{M}_{FG}$, there is a canonical multiplication $E_R\wedge E_R\to E_R$ making $E_R$ into a commutative and associative algebra in the homotopy category of spectra. 
	\end{coro}
	\begin{proof}
		We consider the pullback diagram 
		% https://q.uiver.app/#q=WzAsNCxbMSwxLCJcXG1hdGhjYWx7TX1fe0ZHfSJdLFsxLDAsIlxcU3BlYyBSIl0sWzAsMSwiXFxTcGVjIFIiXSxbMCwwLCJcXFNwZWMgQiJdLFszLDFdLFsxLDBdLFszLDJdLFsyLDBdLFszLDAsIiIsMSx7InN0eWxlIjp7Im5hbWUiOiJjb3JuZXItaW52ZXJzZSJ9fV1d
		\[\begin{tikzcd}
			{\Spec B} & {\Spec R} \\
			{\Spec R} & {\mathcal{M}_{FG}}
			\arrow[from=1-1, to=1-2]
			\arrow[from=1-1, to=2-1]
			\arrow["\ulcorner"{anchor=center, pos=0.125}, draw=none, from=1-1, to=2-2]
			\arrow[from=1-2, to=2-2]
			\arrow[from=2-1, to=2-2]
		\end{tikzcd}\]
		Then we have obviously a map $\Spec R\to\Spec B$. Then this induces a map $E_R\wedge E_R\simeq E_B\to E_R$. The commutativity and associativity is obvious. 
	\end{proof}

	\begin{defn}
		Recall from \cref{Const_flat_ring_spectra}, we have $\pi_n(E_R)=\mathfrak{g}^k$ for $n=2k$ and $0$ otherwise, where $\mathfrak{g}$ is the Lie algebra of the formal group. Then it makes sense to have the following definition. \\
		Let $E$ be a ring spectrum, we say $E$ is even periodic if the following conditions are satisfied. 
		\begin{enumerate}
			\item $\pi_{2k+1}(E)=0$. 
			\item The pairing $\pi_2(E)\otimes_{\pi_0(E)}\pi_{-2}(E)\to\pi_0(E)$ is an isomorphism (in particular, $\pi_2(E)$ is an invertible $\pi_0(E)$-module $\mathcal{L}$ and $\pi_{2n}(E)=\mathcal{L}^{\otimes n}$). 
		\end{enumerate}
	\end{defn}

	\begin{const}
		For a even periodic ring spectrum $E$, it is automatically complex-orientable by the degeneration of the Atiyah Hirzebruch spectral sequence of $E^*(\mathbb{CP}^\infty)$ on the second page. Then we have the calculation 
		\begin{align*}
			E^d(\mathbb{CP}^\infty)\simeq&(\pi_*(E)[[t]])^d \\
			\simeq&\prod_{n=0}^\infty\pi_{2n-d}(E)t^n \\
			\simeq&\prod_{n=0}^\infty\pi_{-d}(E)\otimes_{\pi_0(E)}\pi_{2n}(E)t^n \\
			\simeq&\pi_{-d}(E)\otimes_{\pi_0(E)}E^0(\mathbb{CP}^\infty) \\
		\end{align*}
		Then recall that from the construction \cref{cplx_ori_FGL} we have a coordinatizable formal group over $\pi_*(E)$ given by $\Spf E^*(\mathbb{CP}^\infty)\simeq\Spf \pi_*(E)[[t]]$. By the above calculation, we have $E^*(\mathbb{CP}^\infty)\simeq E^0(\mathbb{CP}^\infty)\otimes_{\pi_0(E)}\pi_*(E)$. Then notice the comultiplication on $E^*(\mathbb{CP}^\infty)$ preserves dimensions, hence it restricts to a comultiplication on $E^0(\mathbb{CP}^\infty)$. Then, this gives a formal group structure on $\Spf E^0(\mathbb{CP}^\infty)$ over $\pi_0(E)$, which base changes to $\Spf E^*(\mathbb{CP}^\infty)$ over $\pi_*(E)$. 
	\end{const}

	\begin{coro}
		Let $\mathcal{C}$ be the category of pairs $(R,\eta)$, where $R$ is a ring and $\eta:\Spec R\to\mathcal{M}_{FG}$ is a flat map. Then the construction $R\mapsto E_R$ determines a fully faithful embedding of $\mathcal{C}$ into the category of commutative algebras in the category of spectra. Moreover, a ring spectrum $E$ belongs to the image iff $E$ is even periodic, and the map $\pi_0(E)\to\mathcal{M}_{FG}$ is flat. 
	\end{coro}
	\begin{proof}
		The construction $E\mapsto(\pi_0(E),\Spf E^0(\mathbb{CP}^\infty))$ is a left inverse for this operation. Conversely, if $E$ is an even periodic ring spectrum, let $R=\pi_0(E)$, then the homology theory $E_R$ is given by $(E_R)_*(X)=MU_*(X)\otimes_L(\pi_*(E))$. Then we have a map $MU\to E$ given by the complex orientation on $E$, and since $E_*(X)$ is always a $\pi_*(E)$ module, we have a map $(E_R)_*(X)\to E_*(X)$. This map is an isomorphism when $X=*$, and also, recall from \cref{Well_def_map_spectra_from_rings} that we have a map of spectra $E_R\to E$ well-defined up to homotopy equivalence. The two information induces an isomorphism $\pi_*(E_R)\to\pi_*(E)$, which gives an equivalence of spectra. Then finally it is obvious that this equivalence is compatible with the ring spectra structures. 
	\end{proof}

\section{$E(n)$ and $K(n)$}

	\begin{const}\label{E_n_spectra}
		Recall the formal group law over $W(k)[[v_1,\dots,v_{n-1}]]$ in \cref{Univ_def_FGL} exhibits it as a Landweber-exact, since $v_0=p,v_1,\dots,v_{n-1}$ is regular by construction and $v_n$ is invertible in $W(k)[[v_1,\dots,v_{n-1}]]/(v_0,v_1,\dots,v_{n-1})\simeq k$ by the assumption that $f$ has height exactly $n$. Thus, by \cref{Landweber_exact_func_thm} and \cref{Const_flat_ring_spectra}, we may construct a spectra $E(n)$ with $\pi_*(E(n))=W(k)[[v_1,\dots,v_{n-1}]][\beta^{\pm1}]$. Notice the construction depends not only on $n$ but also a choice of $k$ and a formal group law of height $n$ on $k$. 
	\end{const}

	Actually, we have a more intuitive way to view $E(n)$. We consider the property of $L_{E(n)}$. 
	\begin{thm}\label{Bousfield_class_E_n}
		A spectrum $X$ is $E(n)$-acyclic (recall the definition from \cref{E_acyc_loc_equiv}) iff the associated sheaf $\mathcal{F}_X$ and $\mathcal{F}_{\Sigma X}$ over $\mathcal{M}_{FG}\times\Spec\mathbb{Z}_p$ are supported on the closed substack $\mathcal{M}_{FG}^{\geq n+1}$. 
	\end{thm}
	\begin{proof}
		\todo{No idea.} 
	\end{proof}

	This implies that the localization with respect to $E(n)$ has the intuition of restricting $\mathcal{F}_X$ to the substack $\mathcal{M}_{FG}^{\leq n}$. In the rest of this section, we attempt to construct the spectra whose associated localization that restricts $\mathcal{F}_X$ to $\mathcal{M}_{FG}^n$. 
	\begin{const}
		Consider the $p$-local complex cobordism spectrum $MU_{(p)}$. This spectrum is an $E^\infty$-ring, so $MU_{(p)}$-modules and relative smash products $M\wedge_{MU_{(p)}}N$ makes sense in a good way. By definition, we have $\pi_*(MU_{(p)})\simeq L_{(p)}\simeq\mathbb{Z}_{(p)}[[t_1,\dots]]$. Again, we let $t_{p^k-1}=v_k$ and $t_0=p$. We define $M(k)$ to be the cofiber of the map $\Sigma^{2k}(MU)_{(p)}\to MU_{(p)}$ given by multiplication by $t_k$. 
	\end{const}

	\begin{lemma}\label{A_infty_M_k}
		Each $M(k)$ is an $A_\infty$ $MU_{(p)}$-algebra. 
	\end{lemma}
	\begin{proof}
		We let the unit of $M(k)$ to be the evident map $u:MU_{(p)}\to M(k)$. Then the $n$-fold smash product $M(k)\wedge_{MU_{(p)}}\cdots\wedge_{MU_{(p)}}M(k)$ can be realized as the total homotopy cofiber (that is, the cofiber of the map from the colimit of the diagram without the conepoint to the conepoint) of the commutative diagram % https://q.uiver.app/#q=WzAsOCxbMiwzLCJNVV97KHApfSJdLFsyLDEsIlxcU2lnbWFeezJrfU1VX3socCl9Il0sWzAsMywiXFxTaWdtYV57Mmt9TVVfeyhwKX0iXSxbMCwxLCJcXFNpZ21hXns0a31NVV97KHApfSJdLFsxLDAsIlxcY2RvdHMiXSxbMywwLCJcXGNkb3RzIl0sWzMsMiwiXFxjZG90cyJdLFsxLDIsIlxcY2RvdHMiXSxbMywxXSxbMSwwXSxbMywyXSxbMiwwXSxbNywyXSxbNiwwXSxbNCwzXSxbNSwxXV0=
		\[\begin{tikzcd}
			& \cdots && \cdots \\
			{\Sigma^{4k}MU_{(p)}} && {\Sigma^{2k}MU_{(p)}} \\
			& \cdots && \cdots \\
			{\Sigma^{2k}MU_{(p)}} && {MU_{(p)}}
			\arrow[from=1-2, to=2-1]
			\arrow[from=1-4, to=2-3]
			\arrow[from=2-1, to=2-3]
			\arrow[from=2-1, to=4-1]
			\arrow[from=2-3, to=4-3]
			\arrow[from=3-2, to=4-1]
			\arrow[from=3-4, to=4-3]
			\arrow[from=4-1, to=4-3]
		\end{tikzcd}\]
		which is shaped like an $n$-dimensional cube with each vertex $\Sigma^{2nk}MU_{(p)}$, and the map between vertices given by multiplying with $t_k$. We prove that these information together indeed organizes into a simplicial object in the category of $MU_{(p)}$-modules. \\
		We first give the multiplication map $M(k)\wedge_{MU_{(p)}}M(k)\to M(k)$. We let $K$ denote the cofiber of the map $(\Sigma^{2k}M(k))^2\to M(k)$ in which the morphism is given by $(t_k,t_k)$. Or equivalently, we define $K$ to be the total cofiber of the following diagram. 
		% https://q.uiver.app/#q=WzAsNCxbMCwwLCIqIl0sWzEsMCwiXFxTaWdtYV57Mmt9TVVfeyhwKX0iXSxbMCwxLCJcXFNpZ21hXnsya31NVV97KHApfSJdLFsxLDEsIk1VX3socCl9Il0sWzAsMV0sWzEsM10sWzAsMl0sWzIsM11d
		\[\begin{tikzcd}
			{*} & {\Sigma^{2k}MU_{(p)}} \\
			{\Sigma^{2k}MU_{(p)}} & {MU_{(p)}}
			\arrow[from=1-1, to=1-2]
			\arrow[from=1-1, to=2-1]
			\arrow[from=1-2, to=2-2]
			\arrow[from=2-1, to=2-2]
		\end{tikzcd}\]
		Then we have a map $\alpha:K\to M(k)$ given by the morphisms of fiber sequences 
		% https://q.uiver.app/#q=WzAsNixbMCwwLCIoXFxTaWdtYV57Mmt9TVVfeyhwKX0pXjIiXSxbMSwwLCJNVV97KHApfSJdLFswLDEsIlxcU2lnbWFeezJrfU1VX3socCl9Il0sWzEsMSwiTVVfeyhwKX0iXSxbMiwwLCJLIl0sWzIsMSwiTShrKSJdLFswLDEsIih0X2ssdF9rKSJdLFswLDIsIihcXGlkLFxcaWQpIiwyXSxbMiwzLCJ0X2siLDJdLFsxLDMsIlxcc2ltZXEiLDJdLFsxLDRdLFszLDVdLFs0LDUsIlxcYWxwaGEiLDAseyJzdHlsZSI6eyJib2R5Ijp7Im5hbWUiOiJkYXNoZWQifX19XV0=
		\[\begin{tikzcd}
			{(\Sigma^{2k}MU_{(p)})^2} & {MU_{(p)}} & K \\
			{\Sigma^{2k}MU_{(p)}} & {MU_{(p)}} & {M(k)}
			\arrow["{(t_k,t_k)}", from=1-1, to=1-2]
			\arrow["{(\id,\id)}"', from=1-1, to=2-1]
			\arrow[from=1-2, to=1-3]
			\arrow["\simeq"', from=1-2, to=2-2]
			\arrow["\alpha", dashed, from=1-3, to=2-3]
			\arrow["{t_k}"', from=2-1, to=2-2]
			\arrow[from=2-2, to=2-3]
		\end{tikzcd}\]
		and $\beta:K\to M(k)\wedge_{MU_{(p)}}M(k)$ given by the morphisms of squares. Then we attempt to factor $\alpha$ through $\beta$, which would give a multiplication map. Notice that $\alpha$ factors through $\beta$ iff the composition $\fib(\beta)\to K\to M(k)$ is null-homotopic. To see this, we notice that $\fib(\beta)$ is equivalently the total cofiber of the diagram (the fiber of the morphism between the diagram associated to $K$ and to $M(k)\wedge_{MU_(p)}M(k)$) 
		% https://q.uiver.app/#q=WzAsNCxbMCwwLCJcXFNpZ21hXns0a31NVV97KHApfSJdLFsxLDAsIioiXSxbMSwxLCIqIl0sWzAsMSwiKiJdLFswLDFdLFsxLDJdLFswLDNdLFszLDJdXQ==
		\[\begin{tikzcd}
			{\Sigma^{4k}MU_{(p)}} & {*} \\
			{*} & {*}
			\arrow[from=1-1, to=1-2]
			\arrow[from=1-1, to=2-1]
			\arrow[from=1-2, to=2-2]
			\arrow[from=2-1, to=2-2]
		\end{tikzcd}\]
		which implies $\fib(\beta)\simeq\Sigma^{4k+1}MU_{(p)}$. Then, notice that since $M(k)$ is a $MU_{(p)}$-module, a map $\Sigma^{4k+1}MU_{(p)}=0$ iff its corresponding element in $\pi_{4k+1}(M(k))$ is $0$, which is obvious since $M(k)$ has homotopy groups concentrated in even degrees. 
		\todo{Finish the proof}. 
	\end{proof}
	\begin{remark}
		When we give the ring structure on $M(k)$, we have used the fact that $\pi_{4k+1}(M(k))=0$ to imply the corresponding element of the map $\fib(\beta)\simeq\Sigma^{4k+1}MU_{(p)}$ is null-homotopic. However, this null-homotopy corresponds to an element of $\pi_{4k+2}MU_{(p)}$, and thus the collection of multiplications of $M(k)$ form a $\pi_{4k+2}(M(k))$ torsor. Then, we consider the action of $\Sigma_2$ on the torsor (given by the action on $M(k)\wedge_{MU_{(p)}}M(k)$). Then we have $H^1(\Sigma_2,\pi_{4k+2}(M(k)))=0$, hence when $p\neq2$, the multiplication on $M(k)$ can be chosen to be homotopy commutative (i.e, $E_2$). 
	\end{remark}

	\begin{defn}\label{K_n_spectra}
		We define $K(n)$ to be the smash product 
		\[ \bigwedge_{k\neq p^n-1}{_{MU_{(p)}}}M(k)\wedge_{MU_{(p)}} MU_{(p)}[v_n^{-1}] \]
		By \cref{A_infty_M_k}, we see that $K(n)$ has an $A_\infty$ algebra structure, and when $p\neq2$, an $E_2$ structure. 
	\end{defn}

	\begin{coro}\label{K_n_homology}
		By construction, we have 
		\[ \pi_*(K(n))\simeq(\pi_*(MU_{(p)}))[v_n^{-1}]/(t_0,t_1,\dots,t_{p^n-2},t_{p^n},\dots)\simeq\mathbb{F}_p[v_n^{\pm1}] \]
		Moreover, for a finite $p$-local spectra $X$, we can use this fact to compute the $K(n)$-homology of $X$ for $n\gg0$ easily. Since $X$ is finite and $p$-local, $H_*(X;\mathbb{F}_p)$ vanishes in almost all degrees. In particular, we may suppose $H_k(X;\mathbb{F}_p)\neq0$ implies $a\leq k\leq b$. Then for $n\geq b-a$, the Atiyah-Hirzebruch spectral sequence of $K(n)_*(X)$ degenerates, giving the fact that $K(n)_*(X)\simeq H_*(X;\mathbb{F}_p)[v_n^{\pm1}]$ (here $v_n$ has dimension $2(p^n-1)$, as usual). 
	\end{coro}

	The localization with respect to $K(n)$ has the intuition of restricting to the strata $\mathcal{M}_{FG}^n$. We will show this by calculating the Bousfield class of $K(n)$. 
	
	\begin{thm}\label{E_n_equiv_E_and_K}
		$E(n)$ is Bousfield equivalent to $E(n-1)\wedge K(n)$. By convention, $E(0)=H\mathbb{Q}[\beta^\pm]$, which is Bousfield equivalent to $H\mathbb{Q}$. That is, a spectrum is $E(n)$-acyclic iff it is both $E(n-1)$-acyclic and $K(n)$-acyclic. 
	\end{thm}
	\begin{proof}
		We first notice that every $p$-local complex oriented cohomology theory $E$ having \\
		$\Spec\pi_*(E)\to\mathcal{M}_{FG}\times\Spec\mathbb{Z}_{(p)}$ a flat covering of $\mathcal{M}_{FG}^{\leq n}$ is Bousfield equivalent to $E(n)$ by \cref{Bousfield_class_E_n}. This allows us to construct another representative of the Bousfield class of $E(n)$. For $m\leq n$, we define
		\[ Z(m):=\bigwedge_{k\geq0,k\neq p^{m'}-1,m\leq m'\leq n}{_{MU_{(p)}}}M(k)\wedge_{MU_{(p)}}MU_{(p)}[v_n^{-1}] \]
		By definition, $Z(n)=K(n)$, and $Z(0)$ is Bousfield equivalent to $E(n)$. Then, let $X$ be an $E(n)$-acyclic spectrum. Since $\mathcal{M}_{FG}^{\geq n+1}\subseteq\mathcal{M}_{FG}^{\geq n}$, it is $E(n-1)$-acyclic, by \cref{Bousfield_class_E_n}. Then, since $Z(0)$ is Bousfield equivalent to $E(n)$, $X\wedge Z(0)\simeq*$. Then, since $X\wedge K(n)\simeq X\wedge Z(n)=(X\wedge Z(0))\wedge_{MU_{(p)}}\bigwedge_{0\leq k<n}{_{MU_{(p)}}}M(p^k-1)$, we have $X\wedge K(n)\simeq *$. Thus $X$ is $K(n)$ acyclic. \\
		Conversely, suppose $X$ is both $K(n)$-acyclic and $E(n-1)$-acyclic. We prove by induction that $X$ is $Z(i)$-acyclic assuming it is $Z(i+1)$-acyclic. Notice that we have a fiber sequence 
		\[ \Sigma^{2(p^i-1)}Z(i)\xrightarrow{v_i} Z(i)\to Z(i+1) \]
		Thus by smashing $X$ to this fiber sequence, we have $v_i$ acts invertibly on $Z(i)\wedge X$. Thus $Z(i)\wedge X\simeq Z(i)[v_i^{-1}]\wedge X$. Thus it suffices to prove $X$ is $Z(i)[v_i^{-1}]$-acyclic. But then, since 
		\[ Z(i)=Z(0)\wedge_{MU_{(p)}}\bigwedge_{0\leq k<i}{_{MU_{(p)}}}M(p^k-1) \]
		it suffices to prove $X$ is $Z(0)[v_i^{-1}]$-acyclic, which is Bousfield equivalent to $E(i)$. But $i<n$, hence this follows from the assumption. 
	\end{proof}

	Then, we have the fact that $E(n-1)$ and $K(n)$ are somewhat disjoint. To prove this, we first need a highly non-trivial fact (which we will prove later), \cref{E_n_smashing}. 
	
	\begin{thm}\label{E_local_K_acyc}
		Let $X$ be any spectrum. Then $L_{E(n-1)}(X)$ is $K(n)$-acyclic. 
	\end{thm}
	\begin{proof}
	By \cref{E_n_smashing}, $L_{E(n-1)}$ is smashing, hence $L_{E(n-1)}(X)\wedge K(n)\simeq X\wedge L_{E(n-1)}(K(n))$. Thus it suffices to show $L_{E(n-1)}(K(n))\simeq*$. In other words, it suffices to prove $E(n-1)\wedge K(n)\simeq*$. This is obvious since $E(n-1)\wedge K(n)$ is complex orientable, and its associated formal group law must be both of height $\leq n-1$ and of height $n$. 
	\end{proof}

	We now present a characterization of the $E(n)$-localization with the help of the two preceding theorems. 
	
	\begin{thm}\label{Pullback_E_n_localization}
		For any spectrum $X$, the pullback diagram 
		% https://q.uiver.app/#q=WzAsNCxbMCwwLCJYJyJdLFswLDEsIkxfe0Uobi0xKX0oWCkiXSxbMSwxLCJMX3tFKG4tMSl9KExfe0sobil9KFgpKSJdLFsxLDAsIkxfe0sobil9KFgpIl0sWzAsM10sWzMsMl0sWzAsMV0sWzEsMl0sWzAsMiwiIiwxLHsic3R5bGUiOnsibmFtZSI6ImNvcm5lci1pbnZlcnNlIn19XV0=
		\[\begin{tikzcd}
			{X'} & {L_{K(n)}(X)} \\
			{L_{E(n-1)}(X)} & {L_{E(n-1)}(L_{K(n)}(X))}
			\arrow[from=1-1, to=1-2]
			\arrow[from=1-1, to=2-1]
			\arrow["\ulcorner"{anchor=center, pos=0.125}, draw=none, from=1-1, to=2-2]
			\arrow[from=1-2, to=2-2]
			\arrow[from=2-1, to=2-2]
		\end{tikzcd}\]
		equipped with the evident map $X\to X'$ exhibits $X'$ as the $E(n)$-localization of $X$. 
	\end{thm}
	\begin{proof}
		Since $E(n)$-acyclic implies $E(n-1)$-acyclic and $K(n)$-acyclic, $E(n-1)$-local and $K(n)$-local implies $E(n)$-local. Then, since $E(n)$-local spectra are closed under pullbacks, $X'$ is $E(n)$-local. It now suffices to prove the map $X\to X'$ is an $E(n)$-equivalence. By \cref{E_n_equiv_E_and_K}, it suffices to show that $X\to X'$ induces both a $K(n)$-equivalence and an $E(n-1)$-equivalence. In other words, we must show that the squares 
		% https://q.uiver.app/#q=WzAsOCxbMCwwLCJYXFx3ZWRnZSBLKG4pIl0sWzAsMSwiTF97RShuLTEpfShYKVxcd2VkZ2UgSyhuKSJdLFsxLDEsIkxfe0Uobi0xKX0oTF97SyhuKX0oWCkpXFx3ZWRnZSBLKG4pIl0sWzEsMCwiTF97SyhuKX0oWClcXHdlZGdlIEsobikiXSxbMCwyLCJYXFx3ZWRnZSBFKG4tMSkiXSxbMSwyLCJMX3tLKG4pfShYKVxcd2VkZ2UgRShuLTEpIl0sWzAsMywiTF97RShuLTEpfShYKVxcd2VkZ2UgRShuLTEpIl0sWzEsMywiTF97RShuLTEpfShMX3tLKG4pfShYKSlcXHdlZGdlIEUobi0xKSJdLFswLDNdLFszLDJdLFswLDFdLFsxLDJdLFs0LDVdLFs1LDddLFs0LDZdLFs2LDddXQ==
		\[\begin{tikzcd}
			{X\wedge K(n)} & {L_{K(n)}(X)\wedge K(n)} \\
			{L_{E(n-1)}(X)\wedge K(n)} & {L_{E(n-1)}(L_{K(n)}(X))\wedge K(n)} \\
			{X\wedge E(n-1)} & {L_{K(n)}(X)\wedge E(n-1)} \\
			{L_{E(n-1)}(X)\wedge E(n-1)} & {L_{E(n-1)}(L_{K(n)}(X))\wedge E(n-1)}
			\arrow[from=1-1, to=1-2]
			\arrow[from=1-1, to=2-1]
			\arrow[from=1-2, to=2-2]
			\arrow[from=2-1, to=2-2]
			\arrow[from=3-1, to=3-2]
			\arrow[from=3-1, to=4-1]
			\arrow[from=3-2, to=4-2]
			\arrow[from=4-1, to=4-2]
		\end{tikzcd}\]
		are pullback squares. Here, the second diagram obviously has two vertical morphisms by definition equivalences, hence is a pullback square. Now, the first square has the top horizontal morphism an equivalence, hence it suffices to prove the bottom morphism is an equivalence. But by \cref{E_local_K_acyc}, the bottom line is the zero morphism between $*$, hence an equivalence. 
	\end{proof}

	\begin{const}\label{Orthog_decomp_E_K}
		If $X$ is an $E(n)$-local spectrum, the above argument shows that $X\simeq L_{E(n-1)}(X)\vee_{L_{E(n-1)}(L_{K(n)})(X)}L_{K(n)}(X)$. Conversely, if we have a $K(n)$-local spectrum $Y$ and an $E(n-1)$-local spectrum $Z$, we can form a pullback diagram 
		% https://q.uiver.app/#q=WzAsNCxbMCwxLCJaIl0sWzEsMSwiTF97RShuLTEpfShZKSJdLFsxLDAsIlkiXSxbMCwwLCJYIl0sWzMsMl0sWzIsMV0sWzMsMF0sWzAsMV0sWzMsMSwiIiwxLHsic3R5bGUiOnsibmFtZSI6ImNvcm5lci1pbnZlcnNlIn19XV0=
		\[\begin{tikzcd}
			X & Y \\
			Z & {L_{E(n-1)}(Y)}
			\arrow[from=1-1, to=1-2]
			\arrow[from=1-1, to=2-1]
			\arrow["\ulcorner"{anchor=center, pos=0.125}, draw=none, from=1-1, to=2-2]
			\arrow[from=1-2, to=2-2]
			\arrow[from=2-1, to=2-2]
		\end{tikzcd}\]
		Then $X$ is $E(n)$-local. Moreover, $Y$ can be identified with $L_{K(n)}(X)$, and $Z$ can be identified with $L_{E(n-1)}(X)$. This implies that the category of $E(n)$-local spectra is equivalent to the category of triples $(Y,Z,\alpha:Z\to L_{E(n-1)}(Y))$. 
	\end{const}

\section{Fields and uniqueness of $K(n)$}

	Notice that in the construction of $M(k)$, the choice of the multiplication of $M(k)$ is not canonically determined. Moreover, since the choice of $t_k$ for $k\neq p^i-1$ is not canonical, the underlying spectrum of $M(k)$ is also not canonically determined. However, when they are put together, surprisingly we have the $K(n)$'s are independent of these choices, though the ring structures on them do depend on the ring structures of $M(k)$. 
	
	\begin{defn}
		Let $R$ be a commutative evenly graded ring. We say that $R$ is a graded field if every non-zero homogeneous element of $R$ is invertible. Equivalently, $R$ is a field iff $R$ is a field in the ordinary sense concentrated at degree $0$ or of the form $k[\beta^{\pm1}]$ for some $\beta$ of positive even degree. We say a ring spectrum $E$ is a field if $\pi_*(E)$ is a field in the above sense. 
	\end{defn}

	\begin{coro}\label{K_n_modules}
		For a field $R$, any graded module $M$ over $R$ admits a free basis of homogeneous elements. 
	\end{coro}
	\begin{proof}
		When $R$ is a field in the ordinary sense, this is obvious. When $R=k[\beta^{\pm1}]$ and $\beta$ has degree $d$, then a $k$-basis of $\bigoplus_{0\leq i<d}M_i$ is a $R$-basis of $M$. 
	\end{proof}

	\begin{coro}\label{Modules_split_over_fields}
		For a ring spectrum $E$ which is a field, every $E$-module $M$ is of the form $\bigvee_{\alpha}\Sigma^{k_\alpha}E$. 
	\end{coro}
	\begin{proof}
		Simply pick a homogeneous $\pi_*(E)$ basis of $\pi_*(M)$, and such a basis determines a map $\bigvee_{\alpha}\Sigma^{k_\alpha}E\to M$, which is obviously a homotopy equivalence. 
	\end{proof}

	\begin{coro}\label{Homology_fields_smash_tensor}
		For spectra $X$ and $Y$ and a field $E$, we have $E_*(X\wedge Y)\simeq E_*(X)\otimes_{\pi_*(E)}E_*(Y)$. 
	\end{coro}
	\begin{proof}
		Notice we have 
		\begin{align*}
			&E\wedge X\wedge Y \\
			=&(\bigvee_{\alpha}\Sigma^{k_\alpha}E)\wedge Y \\
			=&\bigvee_{\alpha}\Sigma^{k_\alpha}(E\wedge Y) \\
			=&\bigvee_{\alpha}\Sigma^{k_\alpha}(\bigvee_{\beta}\Sigma^{\ell_\beta}E) \\
			=&\bigvee_{\alpha,\beta}\Sigma^{k_\alpha+\ell_\beta}E
		\end{align*}
		Then taking homotopy groups gives the results. 
	\end{proof}
	
	\begin{coro}
		$K(n)$ is a field. In particular, $K(0)\simeq H\mathbb{Q}$ and $K(\infty)\simeq H\mathbb{F}_p$ are fields. 
	\end{coro}
	
	\begin{lemma}\label{K_n_map_modules}
		Let $f:X\to Y$ be a map of spectra. Then suppose that $X$ and $Y$ are $K(n)$-modules, we have $f_*:\pi_*(X)\to\pi_*(Y)$ is $\pi_*(K(n))=\mathbb{F}_p[v_n^{\pm}]$-linear. 
	\end{lemma}
	\begin{proof}
		We factor the map as a composition $X\to K(n)\wedge X\to K(n)\wedge Y\to Y$, where the middle spectra are regarded as a $K(n)$-module by the multiplication on $K(n)$. The second and third maps are obviously $K(n)$-module maps, hence it suffices to treat the case where $f$ is of the form $f:X\to K(n)\wedge X$. Since $K(n)$ is a field, $X$ is free, hence it suffices to prove the case for $X=K(n)$. Thus we must prove that the inclusion on the two factors $f,g:K(n)\to K(n)\wedge K(n)$ induces the same map on homotopy groups. In other words, we must show $f_*(v_n)=g_*(v_n)$. \\
		Now that $K(n)\wedge K(n)$ has two complex orientations obtained from each component, we have two formal group laws over $R=\pi_{even}(K(n)\wedge K(n))$. Obviously these formal group laws has $[p]$-series beginning with $f_*(v_n)t^{p^n}$ and $g_*(v_n)t^{p^n}$ by definition. Then, since the formal group laws differ by a coordinate change of the form $t\mapsto t+b_1t^2+\cdots$, we have $f_*(v_n)=g_*(v_n)$. 
	\end{proof}

	\begin{coro}
		For $X$ a spectrum admitting the structure of a $K(n)$-module. Then any retract of $X$, say $Y$, has the structrue of a $K(n)$-module. 
	\end{coro}
	\begin{proof}
		Consider the composition 
		\[ Y\xrightarrow{i}X\xrightarrow{r}Y\xrightarrow{i}X \]
		By definition, the composition of the first two maps is the identity $\id_Y$, and the composition of the last two maps is a map $f:X\to X$. By \cref{K_n_map_modules}, $f_*$ is a $\mathbb{F}_p[v_n^{\pm1}]$-module map from $\pi_*(X)$ to itself. Thus the image of $f_*$ is a graded submodule of $\pi_*(X)$, and has a basis of classes $\eta_\alpha\in\pi_{k_\alpha}X$. Then clearly we have $Y\simeq\bigvee\Sigma^{k_\alpha}K(n)$. 
	\end{proof}

	\begin{coro}
		Let $E$ be any field with $E\wedge K(n)\neq0$. Then $E$ admits the structure of a $K(n)$-modules. 
	\end{coro}
	\begin{proof}
		Notice the unit map of $E$-module structure $E\to E\wedge K(n)$ is non-zero, hence $E$ is a direct summand of $E\wedge K(n)$ since $E$ is a field. Thus by the preceding corollary, we have $E$ a retract of $E\wedge K(n)$, hence admits a $K(n)$-module structure. 
	\end{proof}

	\begin{coro}
		Let $E$ be a complex oriented cohomology theory having a formal group law of height exactly $n$. If $E\neq0$, then $E\wedge K(n)\neq0$. 
	\end{coro}
	\begin{proof}
		Notice $E\wedge E(n-1)$ is a complex oriented cohomology theory having formal group law both of height $\leq n-1$ and exactly $n$, hence is $0$. If $E\wedge K(n)=0$, then $E\wedge E(n)=0$ by \cref{Orthog_decomp_E_K}. Thus the quasi-coherent sheaf on $\mathcal{M}_{FG}$ associated to $E$, $\mathcal{F}_E$, vanishes when we restrict to the open substack $\mathcal{M}_{FG}^{\leq n}$. Since the formal group law on $E\wedge E(n)$ has height $\leq n$, we have $E\simeq 0$. 
	\end{proof}

	\begin{thm}\label{K_n_det_homotopy_grp}
		Let $E$ be any complex oriented ring spectrum whose formal group law is of height exactly $n$, and whose homotopy groups are given by $\pi_*(E)\simeq\mathbb{F}_p[v_n^{\pm1}]$. Then we have an equivalence of spectra $E\simeq K(n)$. 
	\end{thm}
	\begin{proof}
		It follows from the preceding two corollaries directly that $E$ is a free $K(n)$-module. By considering $\pi_*(E)$ we see this module is of rank $1$, hence $E\simeq K(n)$. 
	\end{proof}

	Moreover, we actually have any field $E$ is a $K(n)$-module for some $n$. To prove this, we need some preliminary results. 
	\begin{lemma}
		Let $\{E^\alpha\}$ denote a collection of ring spectra. The following conditions are equivalent. 
		\begin{enumerate}
			\item Let $R$ be a $p$-local ring spectrum. If $x\in\pi_m(R)$ is a homotopy class whose image in $E_0^\alpha(R)$ is zero for all $\alpha$, then $x$ is nilpotent in $\pi_*(R)$. 
			\item Let $R$ be a $p$-local ring spectrum. If $x\in\pi_0(R)$ is a homotopy class whose image in $E_0^\alpha(R)$ is zero for all $\alpha$, then $x$ is nilpotent in $\pi_0(R)$. 
			\item Let $X$ be any $p$-local spectrum. If $x\in\pi_0(X)$ has trivial image under the Hurewicz map $\pi_0(X)\to E_0^\alpha(X)$ for all $\alpha$, then the induced class $x^{\otimes n}$ in $\pi_0(X^{\wedge n})$ is zero for $n\gg 0$. 
			\item Let $X$ be any $p$-local spectrum, and let $F$ be a finite spectrum. If $f:F\to X$ is such that each composite map $F\to X\to X\wedge E^\alpha$ is null, then $f^{\otimes n}:F^{\wedge n}\to X^{\wedge n}$ is null for $n\gg 0$. 
		\end{enumerate}
	\end{lemma}
	\begin{proof}
		$(1)\implies(2)$ is obvious, and $(2)\implies(3)$ by taking $R$ to be $\bigvee_nX^{\wedge n}$. $(3)\implies(4)$ by considering the function spectrum $\Map(F,X)$. Finally, suppose $(4)$ holds. Then let $x\in\pi_m(R)$ be a class whose image vanishes in $E_m^\alpha(R)$ for all $\alpha$. We identify $x$ with a map $\Sigma^mS\to R$. Then $x^n$ can be identified with the composition $\Sigma^{mn}S\to R^{\wedge n}\to R$, where the second map is given by the multiplication. Since $x^{\otimes n}$ is null for $n\gg 0$, $x$ is nilpotent. 
	\end{proof}
	
	\begin{defn}
		We say a collection of ring spectra $\{E^\alpha\}$ detects nilpotence if the above equivalent conditions holds. 
	\end{defn}

	Recall the theorem \cref{Nil_pot_equiv_acyc_loc}. Then we have the following corollary. 
	\begin{coro}
		For an element $x\in\pi_0(X)$, we denote $x^{-1}X$ the colimit $\ilim(S\xrightarrow{x}X\xrightarrow{x}X^{\wedge 2}\to\cdots)$. Then let $E$ be any ring spectrum. Then the $E$-acyclicity of $T$ is equivalent to the condition that $x^{\otimes n}$ vanishes in $E_0(X^{\wedge n})$ for $n\gg0$. 
	\end{coro}
	\begin{proof}
		Consider the ring spectrum $\bigvee X^{\wedge n}$. 
	\end{proof}

	Recall that we have the following theorem. 
	\begin{thm}
		The spectrum $MU$ detects nilpotence. 
	\end{thm}

	This implies the following. 
	\begin{thm}\label{K_n_detect_nilpot}
		The collection $\{K(n)\}$ detects nilpotence. 
	\end{thm}
	\begin{proof}
		We verify the third equivalent condition. Now, we fix $x\in\pi_0(X)$ such that its image in $K(n)_0(X)$ is trivial for all $n$. We wish to prove that for some $n$ we have $x^{\otimes n}$ is trivial. By \cref{Nil_pot_thm}, it suffices to check that the the image of $x$ in $MU_0(X)$ is nilpotent. Thus, it suffices to prove that $MU_*(T)=0$. That is, we must show the quasi-coherent sheaf $\mathcal{F}_{\Sigma^kT}$ on $\mathcal{M}_{FG}$ vanishes for $k\in\mathbb{Z}$. \\
		We consider the fiber sequences $\Sigma^{2k}MU_{(p)}\xrightarrow{t_k}MU_{(p)}\to M(k)$. For $n\geq0$, we let $P(n)$ denote the smash product 
		\[ \bigwedge_{k\neq p^m-1(m\geq n)}{_{MU_{(p)}}}M(k) \]
		Thus we have $\pi_*(P(n))=\mathbb{Z}_{(p)}[v_1,\dots]/(v_0,v_1,\dots,v_{n-1})$. Therefore $P(0)=BP$ (recall the definition of $BP$ in \cref{BP_spectra_const}), which is Landweber-exact. Thus the map $\Spec\pi_*(P(0))\to\mathcal{M}_{FG}\times\Spec\mathbb{Z}_{(p)}$ is faithfully flat. Thus $P(0)_*(X)$ is the pullback of the quasi-coherent sheaf $\mathcal{F}_{\Sigma^{-*}X}$ on $\mathcal{M}_{FG}$. Thus it suffices to show $P(0)_*(T)=0$. Let $P(\infty)=\ilim P(n)$, hence $P(\infty)=H\mathbb{F}_p$. Thus the image of $x$ in $P(\infty)_0(X)$ vanishes, hence the image of $x$ in $P(n)_0(X)$ vanishes for some $n$. Thus, $P(n)_0(T)\simeq*$. \\
		We prove $P(m)_0(T)\simeq*$ for all $m$ using induction. Assume $P(m+1)_0(T)\simeq*$. Then we have the fiber sequence $\Sigma^{2(p^m-1)}P(m)\xrightarrow{v_m}P(m)\to P(m+1)$. Thus the multiplication of $v_m$ is invertible on $P(m)_*(T)$, thus $P(m)_*(T)\simeq P(m)[v_m^{-1}]_*(T)$. Thus it suffices to prove $T$ is $MU_{(p)}[v_m^{-1}]$-acyclic. Notice that this is a Landweber-exact theory whose associated formal group law has height $\leq m$, hence it suffices to prove $T$ is $E(m)$-acyclic. However, since $T$ is $K(k)$-acyclic, by assumption, this follows immediately. 
	\end{proof}

	\begin{coro}\label{p_loc_K_n_smash_nonzero}
		For a nonzero $p$-local ring spectrum $E$, $E\wedge K(n)$ is nonzero for some $n$. 
	\end{coro}
	\begin{proof}
		If not, then every element of $\pi_0(E)$ is nilpotent, in particular, the unit, which implies $E\simeq*$. 
	\end{proof}
	\begin{coro}
		Let $E$ be a ring spectrum such that $\pi_*(E)$ is a graded field. Then $E$ has the structure of a $K(n)$-module for some $p$ and $n$. 
	\end{coro}
	\begin{proof}
		Take $p$ to be the characteristic of $\pi_*(E)$. 
	\end{proof}

\section{The periodicity theorem}

	Instead of arbitrary spectra, we consider the category of $p$-local spectra in this section. 
	\begin{defn}\label{Thick_subcat}
		Let $\mathcal{T}$ be a full subcategory of some stable category $\mathcal{C}$. We say $\mathcal{T}$ is thick if it contains $0$, is closed under the formation of fibers and cofibers, and if every retract of a spectrum in $\mathcal{T}$ is in $\mathcal{T}$. 
	\end{defn}

	\begin{const}\label{Thick_subcat_gen}
		By definition the intersection of two thick subcategories is still thick. So for a collection $\mathcal{F}$ of objects in $\mathcal{C}$, it makes sense to talk of the minimal thick subcategory containing $\mathcal{F}$, denoted by $\mathcal{T}(\mathcal{F})$. This subcategory will be called the thick subcategory generated by $\mathcal{F}$, and is indeed generated by $\mathcal{F}$ in the following sense. 

		Let $\mathcal{F}'$ be the collection of objects in $\mathcal{C}$ which is a suspension of objects in $\mathcal{F}$. Then let $\mathcal{T}_0(\mathcal{F}')=0$, and $\mathcal{T}_1(\mathcal{F}')$ denote the full subcategory spanned by objects $\tilde{X}$ which are retracts of objects $X\in\mathcal{F}$. Then inductively, we define $\mathcal{T}_n(\mathcal{F}')$ to be the full subcategory spanned by objects $\tilde{X}_n$ which are retracts of objects $X_n$, and $X_n$ sits in a fiber sequence $Y\to X_n\to Z$, where $Y\in\mathcal{T}_1{\mathcal{F}'}$ and $Z\in\mathcal{T}_{n-1}(\mathcal{F}')$. Finally we let $\mathcal{T}(\mathcal{F})=\bigcup\mathcal{T}_n(\mathcal{F}')$. This is clearly the minimal thick subcategory containing $\mathcal{F}$. 

		Moreover, the above construction has the obvious property that for $Y\in\mathcal{T}_{n_1}(\mathcal{F})$ and $Z\in\mathcal{T}_{n_2}(\mathcal{F})$, if we have a fiber sequence $Y\to X\to Z$, then $X\in\mathcal{T}_{n_1+n_2}(\mathcal{F})$. 
	\end{const}

	\begin{lemma}\label{Thick_subcat_closed_smash}
		For any thick subcategory $\mathcal{T}$, and let $X\in\mathcal{T}$. Then for any $p$-local spectrum $Y$, we have $X\wedge_{S_{(p)}}Y\in\mathcal{T}$. 
	\end{lemma}
	\begin{proof}
		The subcategory consisting of $Y$ such that $X\wedge_{S_{(p)}}Y\in\mathcal{T}$ is also thick, and contains $S_{(p)}$ by definition. Thus this subcategory is the whole category of finite $p$-local spectra, since finite $p$-local spectra admit finite cell decompositions. 
	\end{proof}

	\begin{const}
		For a thick subcategory $\mathcal{T}$, let $\mathcal{C}$ denote the collection of $p$-local spectra generated by $\mathcal{T}$ under colimits. By the argument of presentable categories and compact objects, every object $X\in\mathcal{C}$ is a filtered colimit of objects in $\mathcal{T}$, and in particular, if $X$ is a finite $p$-local spectrum, then the identity map $X\to\ilim X_\alpha$ factors through some $X_\alpha$ by the compactness. Then $X$ is the retract of $X_\alpha$ hence $X\in\mathcal{T}$. Thus the construction $\mathcal{T}\mapsto\mathcal{C}$ determines a bijection between thick subcategories of finite $p$-local spectra and subcategories of $\mathcal{C}$ of the category of $p$-local spectra stable under desuspension and generated by finite $p$-local spectra under colimits. 
	\end{const}

	\begin{prop}\label{Thick_subcat_gen_smashing_loc}
		Let $\mathcal{C}$ denote a subcategory of $p$-local spectra stable under desuspension and generated by a thick subcategory $\mathcal{T}$ under colimits. Then the localization w.r.t $\mathcal{C}$ is smashing. 
	\end{prop}
	\begin{proof}
		It suffices to prove $\mathcal{C}^\perp$ is closed under colimits, by \cref{Cond_smashing_loc}. Let $X_\alpha\in\mathcal{C}^\perp$. Then for any $Y\in\mathcal{C}$, we prove $[Y,\ilim X_\alpha]_*$ is trivial. However, since $Y\in\mathcal{C}$, $Y=\ilim Y_\beta$ for $Y_\beta\in\mathcal{T}$. Then since $Y_\beta$ is finite, we have $[Y_\beta,\ilim X_\alpha]\simeq\ilim[Y_\beta,X_\alpha]$. Thus 
		\begin{align*}
			&[Y,\ilim X_\alpha] \\
			\simeq&[\ilim Y_\beta,\ilim X_\alpha] \\
			\simeq&\plim[Y_\beta,\ilim X_\alpha] \\
			\simeq&\plim\ilim[Y_\beta,X_\alpha] \\
			\simeq&0
		\end{align*}
	\end{proof}

	\begin{prop}\label{K_n_acyc_then_K_n-1}
		Let $X$ be a finite $p$-local spectrum. Suppose $K(n)_*(X)=0$ for some $n\geq0$. Then $K(n-1)_*(X)=0$. 
	\end{prop}
	\begin{proof}
		We let $R$ denote the ring spectrum 
		\[ \bigwedge_{k\neq p^n-1,p^{n-1}-1}{_{MU_{(p)}}}M(k)\wedge_{MU_{(p)}}MU_{(p)}[v_n^{-1}] \]
		We assume $n>1$ (when $n=1$, the proof is essentially the same, with some minor changes to the notations). Then, $\pi_*(R)\simeq\mathbb{F}_p[v_{n-1},v_n^{\pm1}]$. In particular, $\pi_0(R)$ is equivalent to the polynomial ring $\mathbb{F}_p[v_{n-1}^av_n^{-b}]$ where $(a,b)$ is the minimal solution to $a(p^{n-1}-1)=b(p^n-1)$. Then, for every integer $k$, $R_k(X)$ is a finitely generated module over $\pi_0(R)$. \\
		Consider the fiber sequence $\Sigma^{2(p^{n-1}-1)}R\xrightarrow{v_{n-1}}R\to K(n)$. Since $K(n)_*(X)=0$, the multiplication by $v_{n-1}$ is invertible on $R_*(X)$. Thus, the multiplication of $v_{n-1}^av_n^{-b}:=t$ is invertible on each $R_k(X)$. Notice that if $R_k(X)$ has a $\mathbb{F}_p[t]$-free part, then $t$ is obviously not invertible, hence $R_k(X)$ is a torsion $\mathbb{F}_p[t]$-module. In particular, there is a polynomial $f$ such that $f$ annihilates each $R_k(X)$ for $0\leq k<2(p^n-1)$, and since $\pi_*(R)$ is periodic of period $2(p^n-1)$, $f$ annihilates $\pi_*(R)$. Since $t$ has invertible action, we may assume $f$ is divisible by $t$. Then for $k\gg0$, the product $f(t)v_n^k$ can be written as a polynomial in $v_{n-1}$ and $v_n$, and hence comes from $\pi_*(MU_{(p)})$. Thus multiplication by $f(t)$ on $R$ is well-defined, and we may localize $R$ to $R[f(t)^{-1}]$ with $R[f(t)^{-1}]_*(X)\simeq R_*(X)[f(t)^{-1}]=0$. \\
		Then, $R[f(t)^{-1}]$ has a complex orientation, which induces a formal group law of height exactly $n-1$ (since $f(t)$ is divisible by $t$, $v_{n-1}$ is invertible in $\pi_*(R[f(t)^{-1}])$). Thus $R[f(t)^{-1}]\wedge K(m)$ vanishes for $m\neq n-1$, and hence by \cref{p_loc_K_n_smash_nonzero}, $R[f(t)^{-1}]\wedge K(n-1)\neq0$, and contains $K(n-1)$ as a retract. But $X\wedge R[f(t)^{-1}]\simeq0$, thus $X\wedge R[f(t)^{-1}]\wedge K(n-1)\simeq0$ and thus $X\wedge K(n-1)\simeq0$. 
	\end{proof}

	\begin{defn}
		We say a $p$-local spectrum $X$ has type $n$ if $K(n)_*(X)\neq0$ but $K(m)_*(X)=0$ for all $m<n$. By the preceding theorem, every spectrum is of type $n$ for some $n$. For exaxmple, $X$ has type 0 iff $H_*(X;\mathbb{Q})\neq0$, or equivalently iff $H_*(X;\mathbb{Z})$ is not torsion. From now on, we let $\mathcal{C}_{\geq n}$ denote the collection of finite $p$-local spectra of type $\geq n$. 
	\end{defn}

	\begin{lemma}
		$\mathcal{C}_{\geq n}$ are thick subcategories. 
	\end{lemma}
	\begin{proof}
		It is obviously closed under the formation of fibers and cofibers by the long exact sequence of $K(m)$-homology. Then, the retracts of spectra induce retracts of homology, thus this category is also closed under retracts. 
	\end{proof}

	\begin{thm}\label{Thick_subcat_thm}
		Every thick subcategory is of the form $\mathcal{C}_{\leq n}$ for some $n$. 
	\end{thm}
	\begin{proof}
		If $\mathcal{T}=0$, then we have $\mathcal{T}=\mathcal{C}_{\leq\infty}$. Otherwise, we claim that the thick subcategory $\mathcal{T}=\mathcal{C}_{\leq n}$, where $n$ is the minimal $n$ such that there exists a spectrum of type $n$ in $\mathcal{T}$. We prove that for $X\in\mathcal{T}$ of type $n$, and $Y$ in $\mathcal{T}$ of type $\geq n$, then $Y\in\mathcal{T}$. \\
		Let $DX$ denote the $p$-local Spanier-Whitehead dual of $X$. Then the identity map $X\to X$ is classified by a map $e:S_{(p)}\to X\wedge_{S_{(p)}}DX$. Since $X$ has type $n$, $e$ induces an injection for $m\geq n$ (by \cref{Homology_fields_smash_tensor})
		\[ K(m)_*(S_{(p)})\to K(m)_*(X\wedge_{S_{(p)}}DX)\simeq K(m)_*(X)\otimes_{\mathbb{F}_p[v_m^{\pm1}]}K(m)_*(X)^\wedge \]
		here $K(m)_*(X)^\wedge$ denote the algebraic dual of $K(m)_*(X)$. We consider the fiber $F=\fib(S_{(p)}\to X\wedge DX)$. Thus the map $f:F\to S_{(p)}$ induces a zero map $K(m)_*(F)\to K(m)_*(S_{(p)})$ for $m\geq n$, and we may consider the composite $g:F\to S_{(p)}\to Y\wedge DY$. Then $g$ induces the zero map $K(m)_*(F)\to K(m)_*(Y\wedge DY)$ for $m\geq n$ since $f$ is the zero map, and also when $m<n$ since $Y$ is of type $n$. Thus by \cref{K_n_detect_nilpot}, some smash product of $g^{\wedge k}:F^{\wedge k}\to(Y\wedge DY)^{\wedge k}$ is null. Then, since $Y\wedge DY$ is a ring spectra (with multiplication given by the counit map), we obtain a null-homotopic map $F^{\wedge k}\to Y\wedge DY$. By the definition of the counit maps, this map corresponds to the composition 
		\[ F^{\wedge k}\wedge Y\xrightarrow{f}F^{\wedge k-1}\wedge Y\to\cdots\to Y \]
		Thus $Y$ is a retract of the cofiber $Y/(F^{\wedge k}\wedge Y)$. Also, notice that $Y/(F^{\wedge k}\wedge Y)$ admit a finite filtration $(F^{\wedge a}\wedge Y)/(F^{\wedge a+1}\wedge Y)$. Thus it suffices to prove each $F^{\wedge a}\wedge Y/(F^{\wedge a+1}\wedge Y)$ is in $\mathcal{T}$, since it is thick. However, these spectra are of the form $F^{\wedge a}\wedge Y\wedge (S_{(p)}/F)\simeq F^{\wedge a}\wedge Y\wedge DX\wedge X$, and thus belongs to $\mathcal{T}$ since $X$ belongs to $\mathcal{T}$ and everything else are finite $p$-local, by \cref{Thick_subcat_closed_smash}. 
	\end{proof}

	We have seen that all thick subcategories are of the form $\mathcal{C}_{\geq n}$. However, we still need to prove that the $\mathcal{C}_{\geq n}$ are distinct. That is, we need to construct for each $n$ a finite $p$-local spectrum $X$ of type $n$. However, this, along with a stronger assumption, is a part of a famous construction. We will use this fact without proof, later. 

	\begin{defn}
		Let $X$ be a finite $p$-local spectrum, and let $n\geq1$. Then we define a $v_n$-self map is a map $f:\Sigma^kX\to X$ such that $f$ induces an isomorphism $K(n)_*(X)\to K(n)_*(X)$, and for $m\neq n$, the induced map $K(m)_*(X)\to K(m)_*(X)$ is nilpotent. \\
		Notice that $X\wedge DX$ carries a natural ring spectrum structure, and giving a $v_n$-self map of degree $k$ is simply giving an element in $\pi_k(X\wedge DX)$. Thus we define for a $p$-local ring spectrum, an element $x\in\pi_k(R)$ is an $v_n$-element iff the image of $x$ in $K(m)_*(R)$ is nilpotent for $m\neq n$, and invertible for $m=n$. \\
		Also, notice that a $v_n$-element induces a $v_n$-self map of $R$ given by multiplication. In particular, this implies that $R$ has type $>0$, hence has homotopy groups are torsion. 
	\end{defn}

	\begin{lemma}\label{v_n_implies_type_geq_n}
		If $X$ is a finite $p$-local spectrum that admits a $v_n$-self map $f$, $X$ has type $\geq n$. 
	\end{lemma}
	\begin{proof}
		If $X$ has type $m<n$, then since $f$ is not an isomorphism on $K(m)$-homology, we have $\cof(f)=X/f$ has non-vanishing $K(m)$-homology, but has vanishing $K(n)$-homology, which contradicts \cref{K_n_acyc_then_K_n-1}. 
	\end{proof}

	\begin{remark}
		The zero map $0:X\to X$ is a $v_n$-self map if $X$ is of type $>n$. 
	\end{remark}

	\begin{const}\label{type_n+1_from_n}
		Let $X$ be a spectrum of type $n$. Then, let $f:\Sigma^kX\to X$ be a $v_n$-self map. We consider $X/f:=\cof(f)$. By definition, $X/f$ is of type $\geq n$. However, since $f$ is a $v_n$-self map, it induces an isomorphism on $K(n)$-homology. Thus $X/f$ has trivial $K(n)$-homology, hence is of type $>n$. But $f$ does not induce an isomorphism on $K(n+1)$-homology by the definition of $v_n$-self maps, hence $X/f$ must be of type $n+1$. 
	\end{const}

	\begin{thm}\label{Periodicity_thm}
		Let $X$ be a finite $p$-local spectrum of type $\geq n$. Then $X$ admits a $v_n$-self map. 
	\end{thm}

	Again, we need a few lemmas before the proof. 

	\begin{lemma}\label{Periodicity_thm_lem_1}
		Let $R$ be a $\mathbb{Z}_{(p)}$-algebra, and let $x,y\in R$ be commuting elements such that $x-y$ is torsion and nilpotent. Then $x^{p^k}=y^{p^k}$ for $k\gg0$. 
	\end{lemma}
	\begin{proof}
		We have 
		\[ x^{p^k}=(y+(x-y))^{p^k}=y^{p^k}+\sum_{i=1}^{p^k}\binom{p^k}{i}y^{p^k-i}(x-y)^i \]
		Now, suppose $(x-y)$ is nilpotent of degree less than $p^a$ (that is, $(x-y)^{p^a}=0$), the right hand side could be rewritten as 
		\[ x^{p^k}=y^{p^k}+\sum_{i=1}^{p^a-1}\frac{p^k}{i}\binom{p^k-1}{i-1}y^{p^k-i}(x-y)^i \]
		Now, $p^k/i$ ($R$ is a $\mathbb{Z}_{(p)}$ algebra, so this makes sense) is divisible by $p^{k-a}$ for all $i$, hence for $k\gg0$, $p^k(x-y)/i=0$.
	\end{proof}

	\begin{lemma}\label{Periodicity_thm_lem_2}
		Let $R$ be a finite $p$-local spectrum, and let $x\in\pi_k(R)$ be a $v_n$-element. Then a power of $x$ satisfies that $x\mapsto v_n^a\in K(n)_*(R)$, and $x\mapsto 0\in K(m)_*(R)$ for $m\neq n$. This lemma strengthens the assumption of $v_n$-elements. 
	\end{lemma}
	\begin{proof}
		Recall that by \cref{K_n_homology} we have $K(m)_*(R)=H_*(R;\mathbb{F}_p)[v_m^{\pm1}]$ for $m\gg0$. Then, we have $x$ nilpotent in $H_*(R;\mathbb{F}_p)$. Hence for $m\gg0$, we have a power of $x$ mapping to $0\in K(m)_*(R)$. Then, there are finitely many remaining terms $k$ between $n$ and a sufficiently large number such that the image of $x$ in $K(k)_*(R)$ is nilpotent but not necessarily $0$. Since there are only finitely many terms, we may replace $x$ by the degree of nilpotency one by one. Thus we have some $b$ such that $x^b\mapsto 0$ for all $m\neq n$. From now on, we let $x$ be $x^b$. 
		Then, since $K(n)_*(R)$ is a finite module over $\pi_*(K(n))$, we have $K(n)_*(R)/(v_n-1)$ is finite. But note that the image of $x$ in $K(n)_*(R)/(v_n-1)$ is a unit, hence by finiteness a power of $x$ is $1$ in $K(n)_*(R)/(v_n-1)\simeq\bigoplus_{0\leq i<2(p^n-1)}K(n)_iR$. Then by lifting the image back, we have a power of $x$ is the power of $v_n$. 
	\end{proof}

	\begin{lemma}\label{Periodicity_thm_lem_3}
		Let $R$ be a finite $p$-local ring spectrum, and let $x\in\pi_k(R)$ be an $v_n$-element. After raising $x$ to a suitable power, we may assume $x$ is central in $\pi_*(R)$. 
	\end{lemma}
	\begin{proof}
		We first assume $x$ satisfies the additional properties of \cref{Periodicity_thm_lem_2}. Let $A=R\wedge DR$, and let $a,b\in A$ be determined by the self maps of $R$ given by left and right multiplication by $x$. Then obviously $a$ and $b$ commute. Since $A$ is of type $>0$, $a-b\in\pi_kA$ is torsion. Then, it suffices to prove that $a-b$ is nilpotent, thus by \cref{Periodicity_thm_lem_1}, we have $a^{p^k}=b^{p^k}$, and by taking $x$ to be $x^{p^k}$, we see that the left multiplication and right multiplication coincide. \\
		By \cref{K_n_detect_nilpot}, it suffices to prove the image of $a-b$ vanishes in all $K(m)_*(A)$. In other words, we must prove the compositions
		\[ R\to K(m)\wedge R\xrightarrow{x\times}K(m)\wedge R\hspace{5em}R\to K(m)\wedge R\xrightarrow{\times x}K(m)\wedge R \]
		agree. If $m\neq n$, this is obvious since $x$ has image $0$ in $K(m)_*(R)$. For $m=n$, it suffices to prove the left and right multiplication by $v_n^a$ induce the same self map of $K(m)\wedge R$. But this is clear, since $v_n$ is the in the image of the unit map $\pi_*(MU_{(p)})\to\pi_*(K(n))$, hence central. 
	\end{proof}

	\begin{lemma}\label{Periodicity_thm_lem_4}
		Let $R$ be a finite $p$-local ring spectrum and $x,y\in\pi_*(R)$ two $v_n$-elements. Then $x^a=y^b$ for some $a,b$. 
	\end{lemma}
	\begin{proof}
		By \cref{Periodicity_thm_lem_2} and \cref{Periodicity_thm_lem_3}, we may assume $x$ and $y$ satisfies $x,y\mapsto 0\in K(m)_*(R)$ for $m\neq n$ and $x,y\mapsto v_n^j\in K(n)_*(R)$. Moreover, $x$ and $y$ commutes. Since $R$ is of type $>0$, $x-y$ is again torsion. Finally, by the above construction, $x-y=0\in K(m)_*(R)$ for all $m$. Hence $x$ and $y$ satisfies the conditions of \cref{Periodicity_thm_lem_1}, hence $x^{p^j}=y^{p^j}$. Counting in the powers we raised in the above process, we obtain the results. 
	\end{proof}

	\begin{coro}\label{Periodicity_thm_lem_5}
		Let $X$ and $Y$ be spectra which admit $v_n$-self maps $f:\Sigma^aX\to X$ and $g:\Sigma^bY\to Y$. Let $h:X\to Y$ be any map. Then by replacing $f$ and $g$ by suitable powers, we may assume $a=b$, and the diagram 
		% https://q.uiver.app/#q=WzAsNCxbMCwwLCJcXFNpZ21hXmFYIl0sWzEsMCwiXFxTaWdtYV5hWSJdLFsxLDEsIlkiXSxbMCwxLCJYIl0sWzAsMSwiaCJdLFsxLDIsImciXSxbMCwzLCJmIiwyXSxbMywyLCJoIiwyXV0=
		\[\begin{tikzcd}
			{\Sigma^aX} & {\Sigma^aY} \\
			X & Y
			\arrow["h", from=1-1, to=1-2]
			\arrow["f"', from=1-1, to=2-1]
			\arrow["g", from=1-2, to=2-2]
			\arrow["h"', from=2-1, to=2-2]
		\end{tikzcd}\]
		commutes
	\end{coro}
	\begin{proof}
		We view $h$ as a map $e:S\to DX\wedge Y$. Then the commutativity of the diagram is equivalent to $(Df\wedge\id_Y)\circ e\simeq(\id_{DX}\wedge g)\circ e$. This corresponds to two $v_n$ elements in $DX\wedge Y$, hence follows from \cref{Periodicity_thm_lem_4}. 
	\end{proof}

	\begin{lemma}\label{Periodicity_thm_lem_6}
		Let $\mathcal{T}$ denote the collection of finite $p$-local spectra which admit a $v_n$-self map. Then $\mathcal{T}$ is closed. 
	\end{lemma}
	\begin{proof}
		Clearly $0\in\mathcal{T}$ and $\mathcal{T}$ is closed under suspension and desuspension. We first show that $\mathcal{T}$ is closed under the formation of cofibers. Let $h:X\to Y$ be a map of finite $p$-local spectra which admits $v_n$-self maps $f$ and $g$. Then by \cref{Periodicity_thm_lem_5}, we may assume that the diagram 
		% https://q.uiver.app/#q=WzAsNCxbMCwwLCJcXFNpZ21hXmFYIl0sWzEsMCwiXFxTaWdtYV5hWSJdLFsxLDEsIlkiXSxbMCwxLCJYIl0sWzAsMSwiaCJdLFsxLDIsImciXSxbMCwzLCJmIiwyXSxbMywyLCJoIiwyXV0=
		\[\begin{tikzcd}
			{\Sigma^kX} & {\Sigma^kY} \\
			X & Y
			\arrow["h", from=1-1, to=1-2]
			\arrow["f"', from=1-1, to=2-1]
			\arrow["g", from=1-2, to=2-2]
			\arrow["h"', from=2-1, to=2-2]
		\end{tikzcd}\]
		commutes up to homotopy. Further, by \cref{Periodicity_thm_lem_2}, we may assume $f$ and $g$ are zero maps on $K(m)$-homology for $m\neq n$. Then we consider the map $x:\Sigma^k(Y/X)\to Y/X$ given by the above diagram. We claim that $x$ is a $v_n$-self map. By the long exact sequence of $K(n)$-homology, $x$ induces an isomorphism on $K(n)$-homology. For $m\neq n$, we consider the diagram of exact sequences
		% https://q.uiver.app/#q=WzAsNixbMCwwLCJLKG0pX3sqLWt9KFkpIl0sWzEsMCwiSyhtKV97Ki1rfShZL1gpIl0sWzIsMCwiSyhtKV97Ki1rLTF9KFgpIl0sWzAsMSwiSyhtKV8qKFkpIl0sWzEsMSwiSyhtKV8qKFkvWCkiXSxbMiwxLCJLKG0pX3sqLTF9KFgpIl0sWzAsMV0sWzEsMl0sWzIsNSwiMCJdLFsxLDQsIngiLDJdLFswLDMsIjAiLDJdLFszLDRdLFs0LDVdXQ==
		\[\begin{tikzcd}
			{K(m)_{*-k}(Y)} & {K(m)_{*-k}(Y/X)} & {K(m)_{*-k-1}(X)} \\
			{K(m)_*(Y)} & {K(m)_*(Y/X)} & {K(m)_{*-1}(X)}
			\arrow[from=1-1, to=1-2]
			\arrow["0"', from=1-1, to=2-1]
			\arrow[from=1-2, to=1-3]
			\arrow["x"', from=1-2, to=2-2]
			\arrow["0", from=1-3, to=2-3]
			\arrow["{\varphi}", from=2-1, to=2-2]
			\arrow[from=2-2, to=2-3]
		\end{tikzcd}\]
		Then $x$ carries $K(m)_*(Y/X)$ to the image of $\varphi$, and annihilates the image of $\varphi$. Thus $x^2$ is trivial, hence $x$ is a $v_n$-self map. \\
		Finally we prove $\mathcal{T}$ is closed under retracts. For a retract diagram $X\to Y\to X$ and a $v_n$-self map $f:\Sigma^kY\to Y$, the map $\Sigma^kX\to\Sigma^kY\to Y\to X$ is a $v_n$-self map of $X$. 
	\end{proof}

	\begin{thm}\label{Existence_type_n_v_n}
		There is one finite $p$-local spectrum of type $n$ that admits a $v_n$-self map. 
	\end{thm}
	\begin{proof}
		Omitted. 
	\end{proof}
	\begin{proof}
		\cref{Periodicity_thm} is now obvious. Since $\mathcal{T}$ in \cref{Periodicity_thm_lem_6}, by \cref{Thick_subcat_thm}, coincides with some $\mathcal{C}_{\geq m}$, by \cref{Existence_type_n_v_n} and \cref{v_n_implies_type_geq_n}, implies that $\mathcal{T}=\mathcal{C}_{\geq n}$. 
	\end{proof}

\section{Localizations}

	In this section, all localization functors (smashing or not) is assumed to be in the category of $p$-local spectra for some $p$. 
	\begin{defn}\label{Telescope_const}
		Let $X$ be a finite $p$-local spectrum of type $\geq n$. By \cref{Periodicity_thm}, we see that $X$ admits a $v_n$-self map, say $f:\Sigma^kX\to X$. Moreover, any two $v_n$-self maps of $X$ coincide after raising to a suitable power. Thus it makes sense to define $X[f^{-1}]$ to be the colimit 
		\[ \ilim(X\xrightarrow{f}\Sigma^{-k}X\xrightarrow{f}\Sigma^{-2k}\to\cdots) \]
		This is independent of the choice of $f$. 
	\end{defn}

	\begin{prop}\label{Tel_type_loc_coincide}
		By \cref{Adj_func_thm}, the thick subcategory $\mathcal{C}_{\geq n+1}$ (the subcategory of finite $p$-local spectrums of type $>n$) generates a subcategory of $p$-local spectra $\mathcal{D}_{\geq n+1}$ under colimits, hence determines a localization functor from the category of $p$-local spectra $\mathcal{D}$ to $\mathcal{D}_{\geq n+1}$. Then, for a finite $p$-local spectrum of type $\geq n$, we have the canonical map $u:X\to X[f^{-1}]$ exhibiting $X[f^{-1}]$ a $\mathcal{D}_{\geq n+1}$-localization of $X$. We denote this localization $L_n^f$. In short, $X[f^{-1}]\simeq L_n^f(X)$. 
	\end{prop}
	\begin{proof}
		By \cref{Bousfield_localization}, we only need to check two things, namely, the fiber of $u:X\to X[f^{-1}]$ is in $\mathcal{D}_{\geq n+1}$, and $X[f^{-1}]$ is $\mathcal{D}_{\geq n+1}$ local. Notice that $\mathcal{C}_{\geq n+1}$ is itself closed under finite colimits, hence it equivalently generates $\mathcal{D}_{\geq n+1}$ under filtered colimits. Thus $\fib(u)\in\mathcal{D}_{\geq n+1}$ iff $\fib(u)$ is a filtered colimit of objects in $\mathcal{C}_{\geq n+1}$. Also, since $\mathcal{C}_{\geq n+1}$ generates $\mathcal{D}_{\geq n+1}$ under filtered colimits, for $X$ in $\mathcal{D}_{\geq n+1}$ we write $X=\ilim X_\alpha$, then $Y$ is $\mathcal{D}_{\geq n+1}$-local iff $[X,Y]_*=0$ for all $X$, iff $\ilim[X_\alpha,Y]_*=0$ for all $X_\alpha$, and iff $Y$ is $\mathcal{C}_{\geq n+1}$-local. \\
		We first check $\fib(u)$ is a filtered colimit of objects in $\mathcal{C}_{\geq n+1}$. This is since the forming of cofibers commute with filtered colimits, hence we have 
		\[ \cof(u)\simeq\ilim(0\to(\Sigma^{-k}X)/X\to(\Sigma^{-2k}X)/X\to\cdots) \]
		Noticce that $(\Sigma^{-nk}X)/X$ is the cofiber of a power of the $v_n$-self map, hence has type $>n$ by \cref{type_n+1_from_n}. Thus the fiber, as a shift of the cofiber, is the filtered colimit of objects in $\mathcal{C}_{\geq n+1}$. \\
		We then check $X[f^{-1}]$ is $\mathcal{C}_{\geq n+1}$-local. Let $Y$ be a spectrum of type $>n$, then to show every map $Y\to X[f^{-1}]$ is null, it suffices to prove $DY\wedge X[f^{-1}]$ is null-homotopic. By \cref{Periodicity_thm_lem_2}, we may assume $f$ induces the zero map on $K(m)_*(X[f^{-1}])$ for all $m\neq n$. Thus by \cref{K_n_homology}, we have $K(m)_*(DY\wedge X[f^{-1}])=0$ for all $m\neq n$. Thus it now suffices to prove $K(n)_*(DY\wedge X[f^{-1}])=0$. But since $Y$ has type $>n$, this is also zero. 
	\end{proof}

	\begin{coro}
		The localization $X[f^{-1}]$ or $L_n^f(X)$ is smashing. 
	\end{coro}
	\begin{proof}
		By \cref{Thick_subcat_gen_smashing_loc}. 
	\end{proof}

	\begin{prop}
		For a finite $p$-local spectrum $X$, we have a fiber sequence 
		\[ \ilim_{k_0,\dots,k_n}\Sigma^{-n-1}X/(v_0^{k_0},\dots,v_n^{k_n})\to X\to L_n^f(X) \]
		where $v_n$ denote any $v_n$-self map $\Sigma^{\alpha_n}X\to X$. 
	\end{prop}
	\begin{proof}
		We prove by induction. First, for the case $n=0$, we have $L_0^t(X)\simeq X[p^{-1}]$. Therefore we may form the cofiber of diagrams 
		% https://q.uiver.app/#q=WzAsMTIsWzAsMCwiWCJdLFsxLDAsIlgiXSxbMiwwLCJYIl0sWzMsMCwiXFxjZG90cyJdLFswLDEsIlgiXSxbMSwxLCJYIl0sWzIsMSwiWCJdLFszLDEsIlxcY2RvdHMiXSxbMCwyLCIqIl0sWzEsMiwiWC9wIl0sWzIsMiwiWC9wXjIiXSxbMywyLCJcXGNkb3RzIl0sWzAsMSwiXFxpZCJdLFsxLDIsIlxcaWQiXSxbMiwzXSxbNCw1LCJwIl0sWzUsNiwicCJdLFs2LDddLFs4LDldLFs5LDEwXSxbMTAsMTFdLFswLDQsIlxcaWQiLDJdLFsxLDVdLFsyLDZdLFs0LDhdLFs1LDldLFs2LDEwXV0=
		\[\begin{tikzcd}
			X & X & X & \cdots \\
			X & X & X & \cdots \\
			{*} & {X/p} & {X/p^2} & \cdots
			\arrow["\id", from=1-1, to=1-2]
			\arrow["\id"', from=1-1, to=2-1]
			\arrow["\id", from=1-2, to=1-3]
			\arrow[from=1-2, to=2-2]
			\arrow[from=1-3, to=1-4]
			\arrow[from=1-3, to=2-3]
			\arrow["p", from=2-1, to=2-2]
			\arrow[from=2-1, to=3-1]
			\arrow["p", from=2-2, to=2-3]
			\arrow[from=2-2, to=3-2]
			\arrow[from=2-3, to=2-4]
			\arrow[from=2-3, to=3-3]
			\arrow[from=3-1, to=3-2]
			\arrow[from=3-2, to=3-3]
			\arrow[from=3-3, to=3-4]
		\end{tikzcd}\]
		Thus by taking colimits and shifting the resulting fiber sequence of colimits, we have 
		\[ \ilim_k\Sigma^{-1}X/p^k\to X\to L_0^t(X) \]
		a fiber sequence. \\
		Then, suppose we have proven the case $n$. For $n+1$, we form the following diagram 
		% https://q.uiver.app/#q=WzAsNixbMCwwLCJcXGlsaW1cXFNpZ21hXnstbi0xfVgvKHZfMF57a18wfSxcXGRvdHMsdl9uXntrX259KSJdLFsxLDAsIlgiXSxbMiwwLCJMX25edChYKSJdLFswLDEsIkxfe24rMX1edFxcaWxpbVxcU2lnbWFeey1uLTF9WC8odl8wXntrXzB9LFxcZG90cyx2X25ee2tfbn0pIl0sWzEsMSwiTF97bisxfV50KFgpIl0sWzIsMSwiTF97bisxfV50TF9uXnQoWCkiXSxbMCwxXSxbMSwyXSxbMCwzXSxbMSw0XSxbMiw1XSxbMyw0XSxbNCw1XV0=
		\[\begin{tikzcd}
			{\ilim_{k_0,\dots,k_n}\Sigma^{-n-1}X/(v_0^{k_0},\dots,v_n^{k_n})} & X & {L_n^f(X)} \\
			{L_{n+1}^t\ilim_{k_0,\dots,k_n}\Sigma^{-n-1}X/(v_0^{k_0},\dots,v_n^{k_n})} & {L_{n+1}^t(X)} & {L_{n+1}^tL_n^f(X)}
			\arrow[from=1-1, to=1-2]
			\arrow[from=1-1, to=2-1]
			\arrow[from=1-2, to=1-3]
			\arrow[from=1-2, to=2-2]
			\arrow[from=1-3, to=2-3]
			\arrow[from=2-1, to=2-2]
			\arrow[from=2-2, to=2-3]
		\end{tikzcd}\]
		Since $\mathcal{C}_{\geq n+1}\supset\mathcal{C}_{\geq n+2}$, the right map is an equivalence. Then, since $L_n^f$ is smashing, it commutes with colimits. Thus the fiber of the left map and the middle map coincide. Thus 
		\[ \fib(X\to L_{n+1}^t(X))\simeq\ilim_{k_0,\dots,k_n}\fib(\Sigma^{-n-1}X/(v_0^{k_0},\dots,v_n^{k_n})\to L_{n+1}^t\Sigma^{-n-1}X/(v_0^{k_0},\dots,v_n^{k_n})) \]
		Then, by definition, $X/(v_0^{k_0},\dots,v_n^{k_n})$ is of type $>n$, hence $L_{n+1}^t\Sigma^{-n-1}X/(v_0^{k_0},\dots,v_n^{k_n})$ is given by $\Sigma^{-n-1}X/(v_0^{k_0},\dots,v_n^{k_n})[v_{n+1}^{-1}]$. Therefore the fiber can be identified with 
		\[ \ilim_{k_0,\dots,k_{n+1}}\Sigma^{-n-2}X/(v_0^{k_0},\dots,v_{n+1}^{k_{n+1}}) \]
	\end{proof}

	\begin{const}
		Notice $\mathcal{C}_{\geq n+1}^\perp$ is closed under colimits (as in the proof of \cref{Tel_type_loc_coincide}). Then we consider the Bousfield localization functor induced by $\mathcal{C}_{\geq n+1}^\perp$. We define the colocalization functor to be $D(X)$ and the localization functor $R(X)$. Then the fiber sequence $D(X)\to X\to R(X)$ coincides with the 
		\todo{Finish the localization w.r.t $\mathcal{C}_{\geq n+1}^\perp$.}
	\end{const}

	Then we prove the following theorem characterizing smashing localizations. 
	\begin{thm}\label{Smash_loc_ker_type}
		Let $L$ be a smashing localization functor. Then $L$ satisfies one of the following conditions. 
		\begin{enumerate}
			\item $L$ is the zero functor. 
			\item $L$ is the identity functor. 
			\item $\ker L_n^f\subseteq\ker L\subseteq\ker L_{E(n)}$. 
		\end{enumerate}
	\end{thm}

	\begin{lemma}\label{Smash_loc_ker_type_lem_1}
		Let $L$ be a localization functor. For $0\leq n\leq\infty$, we either have $LK(n)\simeq*$ or $LK(n)\simeq K(n)$. 
	\end{lemma}
	\begin{proof}
		By the definition of localization, we have a functor $K(n)\to LK(n)$, which gives $LK(n)$ a structure of a $K(n)$-module. Moreover, by \cref{K_n_modules}, if $LK(n)\neq*$ we have $K(n)$ a retract of $LK(n)$. Thus since the retract of a local object is still local, we have $K(n)$ is local, hence $LK(n)\simeq K(n)$. 
	\end{proof}

	\begin{lemma}\label{Smash_loc_ker_type_lem_2}
		Let $L$ be a smashing localization functor and let $E$ be a nonzero complex-oriented cohomology theory whose formal group has height exactly $n$, then $LE\simeq*$ iff $LK(n)\simeq*$. 
	\end{lemma}
	\begin{proof}
		If $LE\simeq*$, then $*\simeq K(n)\wedge LE\simeq LK(n)\wedge E$. Since $K(n)\wedge E\neq0$, $LK(n)$ must be trivial by \cref{Smash_loc_ker_type_lem_1}. 
	\end{proof}

	\begin{lemma}\label{Smash_loc_ker_type_lem_3}
		Let $L$ be a smashing localization functor. If $LK(m)\simeq*$, then $LK(n)\simeq*$ for $n>m$. 
	\end{lemma}
	\begin{proof}
		For $k\geq 0$, recall $M(k)$ is the cofiber of the map $t_k:\Sigma^{2k}MU_{(p)}\to MU_{(p)}$, and let $R$ be the following ring spectrum 
		\[ R:=\bigwedge_{k\neq p^m-1,p^n-1}{_{MU_{(p)}}}M(k)\wedge_{MU_{(p)}}MU_{(p)}[v_n^{-1}] \]
		We assume $0<m<n<\infty$, so that we can write $\pi_*(R)\simeq\mathbb{F}_p[v_m,v_n^{\pm1}]$. The other cases are essentially the same. Notice $R[v_m^{-1}]$ is a ring spectrum whose formal group law has height exactly $m$. Since $L$ is smashing, it commutes with colimits, hence $LR[v_m^{-1}]$ is the following colimit. 
		\[ \ilim(LR\xrightarrow{v_m}\Sigma^{-2(p^m-1)}LR\to\cdots) \]
		But by \cref{Smash_loc_ker_type_lem_2}, $LR[v_m^{-1}]\simeq*$. Thus $1\in\pi_0LR$ vanishes in $\pi_{2k(p^m-1)}(LR)$ for $k\gg0$. Or equivalently, $v_m$ is nilpotent in $\pi_*(LR)$, say $v_m^k=0$. Let $R'$ denote the cofiber of $v_m^{k+1}:\Sigma^{2(k+1)(p^m-1)}R\to R$, hence $v_m^k$ vanishes in $\pi_*(LR')$ (by the long exact sequence, and notice that the homotopy groups of $R$ is concentrated on even degrees). Thus the map on homotopy groups $\pi_*(R')\to\pi_*(LR')$ is not surjective, hence $R'$ is not local. But $R'$ can be obtained as a successive extension of $k+1$ copies of $R/v_m\simeq K(n)$, thus $K(n)$ is not local. By \cref{Smash_loc_ker_type_lem_1}. 
	\end{proof}

	\begin{lemma}\label{Smash_loc_ker_type_lem_4}
		Let $L$ be a smashing localization, and $LK(n)\simeq K(n)$. Then every spectrum belonging to $\ker L$ is $E(n)$-acyclic. 
	\end{lemma}
	\begin{proof}
		Recall $E(n)$ is Bousfield equivalent to $K(0)\vee\cdots\vee K(n)$. And by \cref{Smash_loc_ker_type_lem_3}, we have $LK(m)\simeq K(m)$ for all $m\leq n$. Finally, we have $K(m)\wedge X\simeq LK(m)\wedge X\simeq K(m)\wedge LX\simeq 0$ for $X\in\ker L$. 
	\end{proof}

	Unfortunately, we will need a theorem that we will prove later (in fact, we still need to prove that the localization w.r.t $E(n)$ is smashing) to analyze the structure of the smashing localization. Namely, we will need \cref{Chromatic_conv_thm}. 
	\begin{lemma}\label{Smash_loc_ker_type_lem_5}
		Let $L$ be a smashing localization such that $LK(n)\simeq K(n)$ for $0\leq n<\infty$. Then $L$ is equivalent to the identity functor. 
	\end{lemma}
	\begin{proof}
		By \cref{Chromatic_conv_thm}, we have 
		\[ \plim(\cdots\to L_{E(1)}S_{(p)}\to L_{E(0)}S_{(p)})\simeq S_{(p)} \]
		Then, by \cref{Smash_loc_ker_type_lem_4}, we have $L_{E(n)}S_{(p)}$ is $E(n)$-local, hence is local w.r.t $L$. Thus $S_{(p)}$, a limit of $L$-local spectra, is $L$-local. Thus $LX\simeq X\wedge LS_{(p)}\simeq X\wedge S_{(p)}\simeq X$. 
	\end{proof}

	\begin{thm}\label{Smash_loc_ker_type_lem_6}
		Let $L$ be a smashing localization functor, and let $n\geq 0$ be an integer. The following conditions are equivalent. 
		\begin{enumerate}
			\item $LK(n)\simeq*$. 
			\item $LK(m)\simeq*$ for all $n\leq m\leq\infty$. 
			\item Every finite $p$-local spectrum $X$ of type $\geq n$ belongs to $\ker L$. 
			\item There is one finite $p$-local spectrum $X$ of type $\geq n$ belonging to $\ker L$. 
		\end{enumerate}
	\end{thm}
	\begin{proof}
		$(1)\implies(2)$ and $(3)\implies(4)$: Obvious. \\
		$(4)\implies(1)$: Notice $LX\simeq*$ implies $LX\wedge K(n)\simeq X\wedge LK(n)\simeq*$, but since $X$ is of type $n$, $X\wedge K(n)\neq*$, thus $LK(n)\simeq*$, by \cref{Smash_loc_ker_type_lem_1}. \\
		$(2)\implies(3)$: Let $X$ be a finite $p$-local spectrum of type $\geq n$. We prove $LX\simeq*$. Then let $R=X\wedge DX$. Since $LX$ is a $LR$-module, it suffices to prove $LR\simeq*$. Since $LR$ is a ring spectrum, by \cref{K_n_detect_nilpot}, it suffices to prove $LR\wedge K(m)\simeq*$ for all $m$. If $m<n$, then obviously $LR\wedge K(m)\simeq L(R\wedge K(m))=0$, and $R$ is of type $n>m$. For $m\geq n$, this is obvious by the assumption. 
	\end{proof}

	\begin{proof}
		We now prove \cref{Smash_loc_ker_type}. \\
		By \cref{Smash_loc_ker_type_lem_6}, we see that every smashing localization functor is of one of the three cases. 
		\begin{enumerate}
			\item $LK(m)\simeq*$ for all $0\leq m<\infty$. 
			\item $LK(m)\simeq K(m)$ for all $0\leq m<\infty$. 
			\item There is some $n\geq0$ such that $LK(n)\simeq K(n)$ but $LK(n+1)=0$. 
		\end{enumerate}
		For the first case, $L$ is the zero functor by \cref{Smash_loc_ker_type_lem_2}. For the second case, this follows from \cref{Smash_loc_ker_type_lem_5}. For the third case, \cref{Smash_loc_ker_type_lem_6} guarantees that $\ker L$ contains every finite spectrum of type $>n$. Conversely, if $X$ is a finite $p$-local spectrum such that $LX\simeq0$, we have 
		\[ 0\simeq K(n)\wedge LX\simeq LK(n)\wedge X\simeq K(n)\wedge X \]
		hence $X$ has type $>n$. Therefore, we have finite $p$-local spectra belong to $\ker L$ iff it is of type $>n$, that is, iff it is $E(n)$-acyclic. Then by \cref{Smash_loc_ker_type_lem_4}, we have $\ker L_n^f\subseteq\ker L\subseteq\ker L_{E(n)}$. 
	\end{proof}

	\begin{remark}
		The last case may seem useless, so Ravenel came up with the telescope conjecture, stating that $L_n^f\simeq L_{E(n)}$. However, this conjecture was disproved in October 2023, by methods of K-theory. 
	\end{remark}

\chapter{The Chromatic Spectral Sequence}

\section{Pro-objects and spectral sequences}

	\begin{const}
		Recall for a spectrum $X$ and a ring spectrum $E$, we have an associated cosimplicial ring spectrum $[n]\mapsto X\wedge E^{\wedge n+1}$, and we have the Adams tower of $X$ converging to $\Tot X^\bullet$. That is, we have
		\[ \Tot X^\bullet\simeq\plim(\cdots\to\Tot_2X^\bullet\to\Tot_1X^\bullet\to\Tot_0X^\bullet) \]
		We have an obvious map $X\to\Tot X^\bullet$. We have seen in the case where $E\simeq H\mathbb{F}_p$ that this map gives the information of the $p$-primary part of the homotopy groups of $X$. In general, we consider how this map is close to a homotopy equivalence. 
	\end{const}

	\begin{prop}\label{Cond_Adams_E_loc}
		For $E$ a ring spectrum and $X$ be any spectrum, then the canonical map $X\to\Tot X^\bullet$ exhibits $\Tot X^\bullet$ as an $E$-localization of $X$ iff $E\wedge\Tot X^\bullet\simeq\Tot(E\wedge X^\bullet)$. Or equivalently, since smash products commutes with finite limits, iff 
		\[ E\wedge\plim\Tot_nX^\bullet\simeq\plim E\wedge\Tot_nX^\bullet \]
	\end{prop}
	\begin{proof}
		Firstly, notice that for $X\to Y$ an $E$-equivalence, we have a natural homotopy equivalence $E\wedge X\to E\wedge Y$, hence $X^\bullet$ depends only on $L_E(X)$. Then, since $\Tot X^\bullet$ is a inverse limit of $E$-local spectra ($E$-modules, in particular), $\Tot X^\bullet$ is $E$-local. Thus $\Tot X^\bullet$ is an $E$-localization of $X$ iff $X\to\Tot X^\bullet$ induces an isomorphism on $E$-homology. Equivalently, iff $E\wedge X\to E\wedge\Tot X^\bullet$ is an homotopy equivalence. \\
		Obviously, there are natural maps $\Tot X^\bullet\to\Tot_i X^\bullet$, hence a natural map $E\wedge\Tot X^\bullet\to\Tot(E\wedge X^\bullet)$. However, notice that $E\wedge X^\bullet$ is a split augmented cosimplicial object with the $-1$-th space $E\wedge X$. Thus the composition $E\wedge X\to E\wedge\Tot X^\bullet\to\Tot(E\wedge X^\bullet)$ is a homotopy equivalence. 
	\end{proof}

	\begin{defn}
		Recall that the pro-category is formed by formally adjoining to a category all filtered limits. In particular, we define $\Pro(\Sp)$ to be the category with objects $\plim X_\alpha$ and morphisms $\Hom(\plim X_\alpha,\plim Y_\beta)\simeq\plim_\beta\ilim_\alpha(X_\alpha,Y_\beta)$ (this notion makes sense since the essential image of the Yoneda embedding is cocompact). We have a natural forgetful functor $U:\Pro(\Sp)\to\Sp$ given by $\plim X_\alpha\mapsto\plim X_\alpha$, where the second limit is taken in $\Sp$. \\
		We say a pro-spectrum $\plim X_\alpha$ is constant if it is equivalent to the constant tower $X$. Or, formally, $\plim X_\alpha\simeq X$. \\
		The monoidal structure on $\Sp$ somewhat extends to $\Pro(\Sp)$ by defining $E\wedge\plim X_\alpha=\plim(E\wedge X_\alpha)$. Since the smash product does not necessarily commutes with limits, the map $E\wedge U(\plim X_\alpha)\to U(E\wedge\plim X_\alpha)$ is not always an equivalence, but in particular it is an equivalence when $X_\alpha$ is constant. 
	\end{defn}

	\begin{coro}\label{Cond_Adams_E_loc_2}
		The equivalent conditions of \cref{Cond_Adams_E_loc} is satisfied when the tower 
		\[ \cdots\to\Tot_2X^\bullet\to\Tot_1X^\bullet\to\Tot_0X^\bullet \]
		is constant as a pro-spectrum. 
	\end{coro}

	So, to be able to apply the corollary, we need a condition to determine if a tower is constant as a pro-spectrum. This is the following criterion given by Bousfield. 
	\begin{thm}\label{Cond_tower_const_pro}
		Suppose we have a tower of spectra 
		\[ \cdots\to Y(2)\to Y(1)\to Y(0) \]
		Then, suppose there is some $s\geq 1$ such that for every finite spectrum $F$, let $\{E_r^{p,q},d_r\}$ be the spectral sequence (for a full account see \href{https://ncatlab.org/nlab/show/spectral+sequence+of+a+filtered+complex}{nlab} or of the dual sequence see Higher Algebra 1.2.2, but I will give a brief introduction in the proof) associated to the tower 
		\[ \cdots\to F\wedge Y(2)\to F\wedge Y(1)\to F\wedge Y(0) \]
		then the groups $E_s^{p,q}$ vanish for $p\geq s$, then the tower $Y$ is constant as a pro-spectrum. 
	\end{thm}

	Before the proof of the theorem, we need a technical lemma. 
	\begin{lemma}\label{Cond_tower_const_pro_lem_1}
		The composition of two phantom maps is null. 
	\end{lemma}
	\begin{proof}
		Fix a spectrum $X$, and let $u:\bigvee F_\alpha\to X$ be the direct sum of all maps from finite spectra to $X$. Then $X':=\fib(u)$ is the retract of a sum of finite spectra. Then, suppose $f:X\to Y$ and $g:Y\to Z$ be phantom. Since $f$ is phantom, $f\circ u=0$, and hence $f$ is equivalent to the composition $X\to\Sigma X'\to Y$. Then, $g\circ f$ factors as $X\to\Sigma X'\to Y\to Z$. Since $g$ is phantom and $\Sigma X'$ is a retract of a sum of finite spectra, this map is null. 
	\end{proof}

	\begin{proof}
		We first construct the spectral sequence. For $m\leq n$, we let $F(n,m)$ denote the homotopy fiber of the map $Y(n)\to Y(m)$, and let $E_r^{p,q}$ be defined as the image of the map $\pi_q(F(p+r-1,p-1))\to\pi_q(F(p,p-r))$, and $\der_r$ carries $E_r^{p,q}$ to $E_r^{p+r,q-1}\subseteq\pi_{q-1}(F(p+r,p))$ induced by the fiber sequence $F(p+r,p)\to F(p+r,p-r)\to F(p,p-r)$. Notice that if $p<0$, then $F(p,p-r)$ is contractible and thus $E_r^{p,q}$ vanishes. If $p\geq s$, then $E_r^{p,q}$ vanishes by assumption. Thus, if $r\geq s$, then at least one of the groups $E_r^{p,q}$ and $E_r^{p+r,q-1}$ vanishes, hence $d_r$ is always zero. Thus we have the spectral sequence collapses at the $s$-th space. \\
		Now, we suppose $r>p$. Since $F(p,p-r)\simeq Y(p)$, we have $\pi_q(F(p,p-r))\simeq\pi_q(Y(p))$. Also, note that $E_r^{p,q}$ is the image of the composite map 
		\[ \pi_q(F(p+r-1,p-1))\to\pi_q(Y(p+r-1))\to\pi_q(Y(p)) \]
		By definition, the image of the first map is the kernel of $\pi_q(Y(p+r-1))\to\pi_q(Y(p-1))$. And notice the map $\pi_q(Y(p+r-1))\to\pi_q(Y(p-1))$ factors through the map $\pi_q(Y(p))\to\pi_q(Y(p-1))$. It follows that for $r>p$, the group $E_r^{p,q}$ is the intersection $\im(\pi_q(Y(p+r-1))\to\pi_q(Y(p)))\cap\ker(\pi_q(Y(p))\to\pi_q(Y(p-1)))$. Thus for $r>p,s$, the intersecion is independent for $r$. \\
		Then, we prove by induction that the intersection $\im(\pi_q(Y(p+r))\to\pi_q(Y(p)))\cap\ker(\pi_q(Y(p))\to\pi_q(Y(p-k)))$ is independent of $r\gg0$ for all $k$. The case $k=0$ is trivial, and we assume $k>0$. Suppose $r>p,s$, and $x\in\pi_q(Y(p+r))$ has trivial image in $\pi_q(Y(p-k))$. Let $y$ be the image of $x$ in $\pi_q(Y(p))$. We wish to show that $y$ lifts to $\pi_q(Y(p+r+1))$. Let $y'$ denote the image of $y$ in $\pi_q(Y(p-1))$. Then $y'$ belongs to the kernel of the map $\pi_q(Y(p-1))\to\pi_q(Y(p-k))$. Since $y'$ lifts to $\pi_q(Y(p+r))$, the induction hypothesis shows that $y'$ can be lifted to $x'\in\pi_q(Y(p+r+1))$. Subtracting the image of $x'$ from $x$ we may assume $y'=0$, and thus $y\in\ker(\pi_q(Y(p))\to\pi_q(Y(p-1)))$, and thus we have proven the result. \\
		Then we take $k=p-1$, we see that the image of the map $\pi_*(Y(p+r))\to\pi_*(Y(p))$ is independent of $r$ as long as $r\geq p,s$. We denote this image by $A(p)_*$. By construction, we have a sequence of surjections 
		\[ \cdots\to A(2)_*\to A(1)_*\to A(0)_* \]
		Thus by the construction of the spectral sequence, they fit in exact sequences 
		\[ 0\lto E_\infty^{p,*}\lto A(p)_*\lto A(p-1)_*\lto 0 \]
		By assumption, $E_\infty^{p,*}$ vanishes for $p\geq s$. Thus $A(p)_*\to A(p')_*$ are isomorphisms for $p\geq p'\geq s$. This encourages us to consider the tower 
		\[ \cdots\to\pi_*(Y(4s))\xrightarrow{\theta_2}\pi_*(Y(2s))\xrightarrow{\theta_1}\pi_*(Y(s)) \]
		For $m\geq0$, let $Z(m)_*\subseteq\pi_*(Y(2^ms))$ be the kernel of $\theta_m$. Thus $Z(m)_*\cap A(2^ms)_*=0$, since $\theta_m$ restricts to an isomorphism $A(2^ms)_*\to A(2^{m-1}s)_*$. Then by definition, $\theta_m$ maps onto $A(2^{m-1}s)_*$. Thus we have a direct sum decomposition $\pi_*(Y(2^ms))\simeq A(2^ms)_*\oplus Z(m)_*$. Thus $\{\pi_*(Y(2^ms))\}$, as a pro-object in graded Abelian groups, is equivalent to the constant group $A(s)_*$. \\
		Let $Y=\plim Y(p)\simeq\plim_mY(2^ms)$. The Milnor exact sequence 
		\[ 0\lto\plim^1\pi_{*+1}(Y(p))\to\pi_*(Y)\to\plim\pi_*Y(p)\to0 \]
		gives $\pi_*(Y)=A(s)_*$. For each integer $p\geq0$, let $Y(p)/Y$ denote the cofiber of the canonical map $Y\to Y(p)$. Thus by considering the direct sum decomposition, we have the tower
		\[ \cdots\to Y(4s)/Y\to Y(2s)/Y\to Y(s)/Y \]
		has each map in the tower trivial on all homotopy groups. By applying the above argument to $(Y(2^ms)/Y)_*F$ for all finite $F$, we see that each map is a phantom. Now, by \cref{Cond_tower_const_pro_lem_1}, we have the tower 
		\[ \cdots\to Y(16s)/Y\to Y(4s)/Y\to Y(s)/Y \]
		consisting of null maps, hence we have this tower equivalent to $*$. Thus $\{Y(p)\}$ is equivalent to $Y$. 
	\end{proof}

	\begin{coro}\label{Cond_Adams_E_loc_3}
		Let $E$ be a ring spectrum and $X$ be any spectrum. Suppose that there is an integer $s\geq 1$ such that for every finite spectrum $F$, the $E$-based Adams spectral sequence $\{E^{p,q}_r,\der_r\}$ for $X\wedge F$ has $E_r^{p,q}=0$ for $p\geq s$. Then the Adams tower $\{\Tot_nX^\bullet\}$ is equivalent as a pro-object to a constant tower. Thus, the Adams tower exhibits $\Tot(X^\bullet)$ as $L_E(X)$. Also, the canonical map $L_E(X)\wedge Y\to L_E(X\wedge Y)$ is an equivalence for any other $Y$. 
	\end{coro}
	\begin{proof}
		Apply \cref{Cond_tower_const_pro} and \cref{Cond_Adams_E_loc_2} to the Adams tower. 
	\end{proof}

	\begin{lemma}\label{Det_smash_loc_via_Adams}
		Let $E$ be a $p$-local ring spectrum, and suppose that there is one finite $p$-local spectrum $X$ of type $0$ such that the condition in \cref{Cond_Adams_E_loc_3} holds. Then $L_E$ is a smashing localization. 
	\end{lemma}
	\begin{proof}
		Let $\mathcal{T}$ be the collection of finite $p$-local spectra such that for every $p$-local spectrum $Y$ the map $L_E(X)\wedge Y\to L_E(X\wedge Y)$ is an equivalence. Clearly $\mathcal{T}$ is thick. Then, since $\mathcal{T}$ contains $X$, which is of type $0$, by \cref{Thick_subcat_thm}, we have $\mathcal{T}$ contains every finite $p$-local spectrum. Thus $S_{(p)}\in\mathcal{T}$, and $L_E(Y)\simeq L_E(S_{(p)})\wedge Y$ for all $Y$, hence $L_E$ is smashing. 
	\end{proof}

	We now prove \cref{E_n_smashing}. 
	\begin{thm}\label{E_n_smashing}
		The localization $L_{E(n)}$ is smashing. 
	\end{thm}

	\begin{lemma}\label{E_n_smashing_lem_1}
		
	\end{lemma}







\section{The chromatic tower}

	\begin{thm}\label{Chromatic_conv_thm}
		
	\end{thm}

\section{The J-homomorphism}















	
	


	









	



	

	






\chapter{Appendix: Formal Schemes}

\section{Basic definitions}

	\begin{const}
		Recall that we have the category of affine schemes $\Aff$ dual to the category $\Ring$. Therefore, by the Yoneda lemma, we may regard the category $\Aff$ as the full subcategory of $\Fun(\Ring,\Set)$ consisting of representable objects. Obviously $\Aff\simeq\Ring^{\op}$. The following chapter will mainly be focused on affine schemes in this viewpoint. 
	\end{const}
	
	We begin with some basic affine schemes. 
	\begin{exam}
		For a ring $A$, we define $\Spec A$ to be the functor $B\mapsto\Hom(A,B)$. Clearly, all affine schemes are of this form. \\
		We define the scheme $\mathbb{A}^1$ to be the forgetful functor $\Ring\to\Set$. This is representable by the ring $\mathbb{Z}[x]$. \\
		We define the scheme $\mathbb{G}_m$ to be the functor $R\mapsto R^*$, the units in $R$. This is representable by the ring $\mathbb{Z}[x,x^{-1}]$. \\
		We define the scheme $\FGL$ to be the functor $R\mapsto\FGL(R)$. Recall from \cref{Lazard_ring} that we have $\FGL$ is indeed a scheme. 
	\end{exam}
	
	Then we imitate the definition of the structure sheaf on a scheme $X$. 
	\begin{defn}
		For any affine scheme $X$, we define $\mathcal{O}_X$ to be the set $\Hom_{\Aff}(X,\mathbb{A}^1)$. Notice $\mathbb{A}^1$ is a ring object: the scheme $\mathbb{A}^1\times\mathbb{A}^1=\mathbb{A}^2$ sends a ring $R$ to two elements in $R$, and the structure map is given by the addition and multiplication of the two elements. Then $\Hom_{\Aff}(X,\mathbb{A}^1)$ inherits a natural ring structure. 
	\end{defn}
	\begin{remark}
		By the identification $\Aff\simeq\Ring^{\op}$, for $X=\Spec A$, we clearly have $\Hom_{\Aff}(X,\mathbb{A}^1)\simeq\Hom_{\Ring}(\mathbb{Z}[x],A)\simeq A$. 
	\end{remark}

	\begin{lemma}
		$X(-)\simeq\Hom_{\Ring}(\mathcal{O}_X,-)$. 
	\end{lemma}
	\begin{proof}
		Suppose $X$ is represented by a ring $R$, then $X(-)\simeq\Hom_{\Ring}(R,-)$. But then, $\Hom_{\Aff}(X,\mathbb{A}^1)\simeq\Hom_{\Ring}(\mathbb{Z}[x],R)\simeq R$. 
	\end{proof}

	\begin{defn}
		For an affine scheme $X$, we define a point of $X$ to be a ring $R$ and an element $x\in X(R)$. For any scheme, we may define its category of points $\Pt(X)$ with object pairs of the form $(R,x)$, and we have a morphism $f:(R,x)\to (S,y)$ if $f:R\to S$ is a morphism of rings and $X(f)(x)=y$. For a point $(R,x)$, we let $\mathcal{O}_x$ denote its underlying ring. 
	\end{defn}
	\begin{defn}
		Notice that for $(R,x)$ a point of $X$, and $f\in\mathcal{O}_X$, by definition $f$ is a morphism from $X$ to $\mathbb{A}^1$. Thus $(R,x)$ push forwards to a point of $\mathbb{A}^1$, namely, an element of $R$. We will denote this element by $f(x)$. 
	\end{defn}
	\begin{defn}
		For a scheme $X$ and an ideal $I\leq\mathcal{O}_X$, we define a scheme $V(I)$ by 
		\[ V(I)(R)=\{x\in X(R)|f(x)=0,\forall f\in I\} \]
		We call schemes of this form closed subschemes of $X$. \\
		Similarly, for a scheme $X$ and an element $f\in\mathcal{O}_X$, we define a scheme $D(f)$ by 
		\[ D(f)(R)=\{x\in X(R)|f(x)\in R^*\} \]
		We call these schemes the basic open subschemes of $X$. \\
		A locally closed subscheme is a basic open subscheme of a closed subscheme, which is of the form $D(f)\cap V(I)=D(f)\times_XV(I)$. \\
		Finally, we have the intersection of subschemes $D(f)\cap D(g)=D(fg)$, $V(I)\cap V(J)=V(I+J)$. We define the union of closed subschemes to be $V(I)\cup V(J)=V(I\cap J)$, and the schematic union by $V(I)+V(J)=V(IJ)$. 
	\end{defn}
	\begin{remark}
		The union and the schematic union does not coincide in most cases. The schematic union counts multiplicity by infinitesimal thickening. For example, $V(I)+V(I)=V(I^2)$. 
	\end{remark}
	
	We finally recall some basic definitions in algebraic geometry. 
	\begin{defn}
		We define a scheme $X$ to be reduced if $\mathcal{O}_X$ has no non-zero nilpotents, and define $X_{red}=\Spec(\mathcal{O}_X/\sqrt{0})$. We have a natural inclusion $X_{red}\inj X$ induced by the quotient map $\mathcal{O}_X\to\mathcal{O}_X/\sqrt{0}$. \\
		We define $X$ to be connected if it cannot be split into $Y\amalg Z$, equivalently if there is no non-trivial idempotents in $\mathcal{O}_X$. \\
		We define $X$ to be integral if $\mathcal{O}_X$ is an integral domain, and irreducible if $X_{red}$ is integral. \\
		We define $X$ to be Noetherian if $\mathcal{O}_X$ is Noetherian. Then, $X_{red}=\bigcup Y_i$, where $Y_i$ is integral and closed. This decomposition is unique, with $Y_i$ corresponding to the schemes $V(\mathfrak{p}_i)$, where $\mathfrak{p}_i$ are minimal prime ideals of $\mathcal{O}_X$. 
	\end{defn}

\section{Sheaves}

	\begin{defn}
		For a functor $X\in\Fun(\Ring,\Set)$, we define a sheaf $M$ over $X$ with the following data. 
		\begin{enumerate}
			\item For each $(R,x)\in\Pt(X)$, an $R$-module $M_x$. 
			\item For each map $f:(R,x)\to(S,y)\in\Pt(X)$, an isomorphism $\theta(f):S\otimes_RM_x\to M_y$, satisfying the natural functorality conditions. 
		\end{enumerate}
		For sheaves $M$ and $N$ over $X$, we define a map $f:M\to N$ to be maps $f_x:M_x\to N_x$ for all points of $x$ and commuting with all the maps $\theta(f)$. With this definition, we denote $\Sh(X)$ for the category of all sheaves over $X$, 
	\end{defn}

	\begin{lemma}
		The category $\Sh(X)$ has direct sums and tensor products. Moreover, the unit of the direct sum is the zero sheaf, with $M_x=0$ for all $x$. The unit of the tensor product is the sheaf $\mathcal{O}$ defined on each $(R,x)$ to be $R$. 
	\end{lemma}
	\begin{proof}
		Direct by using the pointwise definitions. 
	\end{proof}

	Similar to the classical methods in algebraic geometry, we can form a sheaf on a scheme using a module. 
	\begin{defn}
		Let $X$ be a scheme and $N$ be an $\mathcal{O}_X$-module. Then we may define a sheaf $\tilde{N}$ over $X$ by letting $N_x=R\otimes_{\mathcal{O}_x}N$, where $R$ is given a $\mathcal{O}_X$-algebra structure by $f(r)=f(x)r$. 
	\end{defn}
	Also, we may form the vector bundle scheme using a sheaf on $X$. Through this definition, we may define the sections of a sheaf. 
	\begin{defn}
		For a scheme $X$ and a sheaf $M$ over $X$, we define $\mathbb{A}(M):\Ring\to\Set$ (not necessarily an affine scheme) as $\mathbb{A}(M)(R)=\bigoplus_{x\in X(R)}M_x$. Notice that we have a natural transformation $\mathbb{A}(M)\to X$ by sending $M_x$ to $x$. Therefore we have $\mathbb{A}(M)\in\Fun(\Ring,\Set)/X$. \\
		For a scheme $Y$ over $X$ and a sheaf $M$, we define $\Gamma(Y,M)=\Hom_{\Fun(\Ring,\Set)/X}(Y,\mathbb{A}(M))$. 
	\end{defn}
	\begin{prop}
		When $X$ is an affine scheme, we have $\Gamma(X,-):\Sh(X)\to\Mod(\mathcal{O}_X)$ an equivalence of categories. 
	\end{prop}
	\begin{proof}
		We claim $\Gamma(X,M)\simeq M_{x_0}$, where $x_0$ is the canonical element in $X(\mathcal{O}_X)$. Firstly, the isomorphism is given by $u\mapsto $
	\end{proof}


\chapter{Appendix: Algebraic Stacks}

\section{Basic Definitions}
	As is well-known, stacks are generalizations to sheaves, similar to groupoids generalizing sets. Namely, a stack is a sheaf taking values in groupoids. We give the formal definition as follows (we will focus only on the theory of stacks here, but the arguments apply easily to $\infty$-stacks). 
	\begin{defn}\label{Defn_sieves}
		Let $\mathcal{C}$ be a category, we define a sieve on $\mathcal{C}$ to be a full subcategory $\mathcal{C}'$ such that if $f:c\to d$ is a morphism in $\mathcal{C}$, and $d\in\mathcal{C}'$, then we have $c\in\mathcal{C}'$. For $c\in\mathcal{C}$ an object, we define a sieve on $c$ to be a sieve on the over category $\mathcal{C}_{/c}$. \\
		Notice that if we have a functor between categories $F:\mathcal{C}\to\mathcal{D}$, then for a sieve $\mathcal{D}'$ on $\mathcal{D}$, the category $\mathcal{D}'\times_{\mathcal{D}}\mathcal{C}$ is a sieve on $\mathcal{C}$. We denote this sieve by $F^*\mathcal{D}'$. This construction applies to a morphism $f:c\to d$ and a sieve $S$ on $d$, since $f$ induces a functor $\mathcal{C}_{/c}\to\mathcal{C}_{/d}$. 
	\end{defn}
	\begin{defn}\label{Defn_sites}
		A category $\mathcal{C}$ equipped with sieves on objects (called covering sieves) is called a site if the three following properties are satisfied. 
		\begin{enumerate}
			\item For $c\in\mathcal{C}$, $\mathcal{C}_{/c}$ is a covering sieve. 
			\item For $c\in\mathcal{C}$ and $S$ a covering sieve on $c$, then for any $f:d\to c$, $f^*S$ is a covering sieve on $d$. 
			\item For $c\in\mathcal{C}$, $S$ a covering sieve on $c$, and $T$ an arbitrary sieve on $c$, if for all $f:d\to c$ in $S$ we have $f^*T$ a covering sieve on $d$, then $T$ is a covering sieve on $c$. 
		\end{enumerate}
	\end{defn}
	\begin{const}
		For a sieve $S$ on $c$, we define a subobject ($(-1)$-truncated morphism) of $j(c)$ (here $j$ is the Yoneda embedding) $j(S)$ as follows. We first recall that for a object $c\in\mathcal{C}$, the functor $F:\mathcal{P}(\mathcal{C}_{/c})\to\mathcal{P}(\mathcal{C})_{/j(c)}$ extending $\mathcal{C}_{/c}\to\mathcal{P}(\mathcal{C})_{/j(c)}$ is an equivalence of categories. Then, a sieve $S$ on $c$ induces a functor $\mathcal{C}_{/c}\to\Delta^1$, where the elements in $S$ are sent to $0$ and others to $1$. We consider the inclusion $\Delta^1\inj\iGrpd^{\op}$ given by $0\mapsto*$ and $1\mapsto\phi$. Then we obtain a presheaf $\mathcal{P}(\mathcal{C}_{/c})$. By applying $F$ to this element, we have an element in $\mathcal{P}(\mathcal{C})_{/j(c)}$. Also, notice that $\Delta^1$ is the $(-1)$-truncated object classifyer, the element we obtained is a $(-1)$-truncated morphism $j(S)\inj j(c)$. 
	\end{const}
	\begin{defn}
		We call a functor (presheaf) $\mathcal{F}:\mathcal{C}^{\op}\to\iGrpd$ a stack if we have for all $c$ and all sieves $S$ on $c$, we have $\Hom_{\mathcal{P}(\mathcal{C})}(j(c),\mathcal{F})\to\Hom_{\mathcal{P}(\mathcal{C})}(j(S),\mathcal{F})$ is an isomorphism. 
	\end{defn}
	However, by the straightening and unstraightening theorem, we have another perspective on stacks. 
	\begin{defn}
		We define a functor $F:\mathcal{F}\to\mathcal{C}$ to be a left fibration if it has the right lifting properties with respect to inclusions $\Lambda^n_i\inj\Delta^n$, where $0\leq i<n$. This implies the fiber of $c\in\mathcal{C}$ is a groupoid, which we will denote by $\mathcal{F}(c)$. 
	\end{defn}
	\begin{const}
		Let $\mathcal{F}\to\mathcal{C}$ be a left fibration. Then let $S$ be a covering sieve on $c$ (or equivalently, let $c_i\to c$ be a family of morphisms generating the covering sieve), then we define the category of descent data $\Desc(S,\mathcal{F})$. This category has objects as systems $(X_i,\theta_{i,j})$, where $X_i$ are objects of $\mathcal{F}$ over $c_i$, and $\theta_{i,j}:X_i|_{c_{i,j}}\to X_j|_{c_{i,j}}$ are isomorphisms. The maps $\theta_{i,j}$ are required to satisfy the cocycle conditions. The morphisms from $(X,\theta)$ to $(Y,\rho)$ are families of arrows $f_i:X_i\to Y_i$ over $c_i$ such that $\rho_{i,j}f_j=f_i\theta_{i,j}$. \\
		Then, notice we have an obvious functor $\mathcal{F}(c)\to\Desc(S,\mathcal{F})$ given by restriction. 
	\end{const}
	\begin{defn}
		We define a stack to be a left fibration $F:\mathcal{F}\to\mathcal{C}$ such that the functor $\mathcal{F}(c)\to\Desc(S,\mathcal{F})$ is an equivalence for all $c$ and all covering sieves on $c$. 
	\end{defn}

	\begin{thm}
		The two definitions of stacks are equivalent. 
	\end{thm}
	\begin{proof}
		The equivalence is given by the straightening and unstraightening construction, and we see that the condition $\mathcal{F}(c)\simeq\Desc(S,\mathcal{F})$ is equivalent to the locality (on sieves) condition. 
	\end{proof}

	Although the writer of this note prefers the first definition, we will adapt the second due to historical reasons. 
	We begin with a theorem determining whether a left fibration is a stack. 
	\begin{prop}\label{Cond_stack}
		The above definition is equivalent to the two following conditions. 
		\begin{enumerate}
			\item For all $X$, $Y$ over $c$, the functor $\mathrm{Iso}_c(X,Y):\mathcal{C}_{/c}\to\Set$ sending $f:c'\to c$ to the set of isomorphisms between $f^*X$ and $f^*Y$ is a sheaf. 
			\item For a covering $c_i\to c$, and a collection $X_i$ over $c_i$ with isomorphisms $\theta_{i,j}:X_i|_{c_{i,j}}\to X_j|_{c_{i,j}}$, there is a $X$ over $c$ with $X|_{c_i}\simeq X_i$ inducing the above isomorphisms. 
		\end{enumerate}
	\end{prop}
	\begin{proof}
		One condition corresponds to the fully faithfulness of the functor, and the other essentially surjectiveness. 
	\end{proof}

	\begin{defn}
		For stacks $\mathcal{F}$, $\mathcal{G}$, and $\mathcal{H}$ over a category $\mathcal{C}$ with structure morphism $p_{\mathcal{F}}$, $p_{\mathcal{G}}$, and $p_{\mathcal{H}}$, we define a morphism between $\mathcal{F}$ and $\mathcal{G}$ to be a functor $f:\mathcal{F}\to\mathcal{G}$ such that $p_{\mathcal{G}}\circ f=p_{\mathcal{F}}$. We define $f\circ g=h:\mathcal{F}\to\mathcal{G}\to\mathcal{H}$ if there is a natural isomorphism between the two functors. \\
		We define the category $\Hom(\mathcal{F},\mathcal{G})$ to have objects morphisms, and morphisms natural transformations. 
	\end{defn}

	From now on, we restrict our discussion to algebraic geometry, in which $\mathcal{C}$ is always the category of schemes or the category of schemes over $S$, equipped with some topology (for example Zariski, etale, fppf). 

	\begin{defn}
		For a scheme $X\in\Sch_{/S}$, we define the stack associated to $X$ to be the left fibration $p_X:\Sch_{/X}\to\Sch_{/S}$, where $p_X$ is the natural projection. Similarly we may define a stack associated to an algebraic space (this is a special case of a stack where the preimage of $A\in\Sch_{/S}$ is discrete). In this case, we say the stack is representable by a scheme or algebraic space. 
	\end{defn}

	\begin{const}
		The fiber product of stacks are given by homotopy pullbacks of categories (or, as is more well-known, the comma category). 
	\end{const}
	\begin{defn}
		We define a morphism of stacks $f:\mathcal{F}\to\mathcal{G}$ is representable if for all $U\in\Sch$ and morphisms $U\to\mathcal{G}$ (note here $U$ denote the stack associated to $U$, and we will use this notation repeatedly), the fiber product $U\times_{\mathcal{G}}\mathcal{F}$ is representable by an algebraic space. \\
		Moreover, for a property $\mathfrak{P}$ of morphisms of schemes, if $\mathfrak{P}$ is local w.r.t our chosen topology on $\Sch$ and is stable under basechange (for example separated, quasi-compact), then we say a morphism of stacks $f:\mathcal{F}\to\mathcal{G}$ is $\mathfrak{P}$ iff for all $U$ and morphisms $U\to\mathcal{G}$, the morphism $U\times_{\mathcal{G}}\mathcal{F}$ is $\mathfrak{P}$. 
	\end{defn}

	\begin{defn}
		We define a Deligne-Mumford (Artin) stack to be a stack $\mathcal{F}$ on $\Sch/S$ with the etale (fppf) topology, satisfying the two following properties. 
		\begin{enumerate}
			\item The diagonal $\mathcal{F}\to\mathcal{F}\times_S\mathcal{F}$ is representable, quasi-compact and separated. 
			\item There is a scheme $U$ (called an atlas, and has the intuition of an atlas) and a etale surjective morphism $U\to\mathcal{F}$. 
		\end{enumerate}
	\end{defn}

	We finish this section with an important example (especially important for this notes). 
	\begin{exam}
		For a scheme $X$ over $S$ and an affine flat group $G$ over $S$, we define the stack $[X/G]$ as follows. This stack has the objects over $B\in\Sch/S$ defined as principal $G$-bundles $\pi:E\to B$ equipped with a $G$-equivariant map $f:E\to X$. Morphisms over $g:B\to B'$ are commutative diagrams 
		% https://q.uiver.app/#q=WzAsNSxbMCwwLCJFIl0sWzAsMSwiQiJdLFsxLDAsIkUnIl0sWzEsMSwiQiciXSxbMiwwLCJYIl0sWzEsM10sWzAsMV0sWzIsM10sWzAsMl0sWzIsNF0sWzAsNCwiIiwyLHsiY3VydmUiOi0zfV1d
		\[\begin{tikzcd}
			E & {E'} & X \\
			B & {B'}
			\arrow[from=1-1, to=1-2]
			\arrow[curve={height=-18pt}, from=1-1, to=1-3]
			\arrow[from=1-1, to=2-1]
			\arrow[from=1-2, to=1-3]
			\arrow[from=1-2, to=2-2]
			\arrow[from=2-1, to=2-2]
		\end{tikzcd}\]
		Firstly, it is a stack by the criterion \cref{Cond_stack}: $Iso_B(e,e')$ is non-empty iff $e\simeq e'$ as principal bundles, and when it is non-empty, it is classified by the stablizer of the identity map $X\to X$, hence is a sheaf represented by a scheme. The second condition follows by the fppf descent of principal bundles. Then, we claim it is an Artin stack. Firstly, for any map $U\to[X/G]\times_S[X/G]$ (equivalent to the data of two objects $Y_1$, $Y_2$ over $U$), notice that the pullback $U\times_{[X/G]\times_S[X/G]}[X/G]\to U$ is equivalent to the morphism $\mathrm{Iso}_U(Y_1,Y_2)\to U$. However recall by the above construction that $\mathrm{Iso}$
		\todo{finish the proof!}
	\end{exam}

\section{Constructions}

	\begin{defn}
		For a stack $\mathcal{F}$, we define a stack $\mathcal{G}$ to be a substack of $\mathcal{F}$ if it is a full subcategory of $\mathcal{F}$, and 
		\begin{enumerate}
			\item If $X\in\mathcal{F}$ is in $\mathcal{G}$, then isomorphic objects are also in $\mathcal{G}$. 
			\item For all morphism of schemes $f:U\to V$, if $X$ over $V$ is in $\mathcal{G}$, then so is $f^*X$. 
			\item Let $U_i\to U$ be a cover, then $X$ over $U$ is in $\mathcal{G}$ iff $X|_{U_i}$ is in $\mathcal{G}$ for all $i$. 
		\end{enumerate}
	\end{defn}

	\begin{defn}
		We define $\mathcal{G}$ to be an open (closed) substack of $\mathcal{F}$ if it is a substack and the inclusion functor $\mathcal{G}\to\mathcal{F}$ is representable and an open (closed) immersion. 
	\end{defn}

	\begin{defn}
		For an algebraic stack $\mathcal{F}$ over $S$, the set of points of $\mathcal{F}$ is defined to be the set of equivalence classes of pairs $(k,x)$, where $k$ is a field over $S$, and $x:\Spec k\to\mathcal{F}$ a morphism of stacks over $S$. Two pairs $(k,x)$ and $(k',x')$ are equivalent if there is a common field extension $k''$ and a map $x'':\Spec k''\to\mathcal{F}$ extending $x$ and $x'$. Clearly every point of an algebraic stack is an image of a point of its atlas. 
	\end{defn}

	\begin{defn}
		For an algebraic stack $\mathcal{F}$ with atlas $u:X\to\mathcal{F}$, and $\xi$ a point of $\mathcal{F}$ coming from $x\in X$, we then define the dimension of $\mathcal{F}$ at $\xi$ to be 
		\[ \dim_\xi\mathcal{F}=\dim_xX-\dim_xf \]
		where $\dim_xf$ is defined to be the relative dimension of the pullback $f^{-1}(\xi)\to\xi$. 
		We define the dimension of the algebraic stack to be the maximum dimension of its points. 
	\end{defn}

	\begin{remark}
		The dimension may be negative. 
	\end{remark}

	\begin{defn}\label{QCoh_sheaf_stack}
		We define a quasi-coherent sheaf $\mathcal{S}$ on an algebraic stack $\mathcal{F}$ to consist of the following data. For each morphism from a scheme $X\to\mathcal{F}$, we have a quasi-coherent sheaf $\mathcal{S}_X$ on $X$, and for each $f:X\to Y$ commuting with the morphisms to $\mathcal{F}$, we have an isomorphism $\varphi_f:\mathcal{S}_X\to f^*\mathcal{S}_Y$. Moreover, for morphisms $g\circ f:X\to Y\to Z$, we must have $(f^*\varphi_g)\circ\varphi_f=\varphi_{g\circ f}$. \\
		We say $\mathcal{S}$ is coherent (flat, faithfully flat, finite type, finite presentation, locally free) if $\mathcal{S}_X$ satisfy the corresponding property for all $X$. \\
		The morphisms are defined obviously by $h_X:\mathcal{S}_X\to\mathcal{S}'_X$ commuting with the isomophisms $\varphi$. 
	\end{defn}
	\begin{exam}
		We define the structure sheaf $\mathcal{O}_{\mathcal{F}}$ of a algebraic stack $\mathcal{F}$ to be $(\mathcal{O}_{\mathcal{F}})_X=\mathcal{O}_X$. 
	\end{exam}
	\begin{thm}
		The category of quasi-coherent sheaves over an algebraic stack $\mathcal{F}$, denoted by $\QCoh(\mathcal{F})$, is an Abelian category. 
	\end{thm}
	\begin{proof}
		The proof is similar to the fact that the quasi-coherent sheaves over a scheme is an Abelian category. 
	\end{proof}


\chapter{Appendix: Deformation Theory}
